%% Ngân hàng đề Toán 12 (KNTT) -- Chương I, Bài 3, Dạng 03
%% Dạng: Tiệm cận xiên của đồ thị hàm số
%% Phương án ĐẦU TIÊN trong \cauTN luôn là đáp án đúng (gói trondethi tự đảo).
%% Phần II: \y{nội dung}{d|s}. Phần III: đáp án đặt ở đối số cuối.

\cauTN[id=C1-B3-D03-NB-001,muc=NB,loigiai={Xét hàm số $y=-2x+3+\dfrac{3}{x-1}$. Chia đa thức: $y=-2x+3+\dfrac{3}{x - 1}$; $a=\lim\limits_{x\to +\infty}\dfrac{y}{x}=-2$, $b=\lim\limits_{x\to +\infty}(y+2x)=3$ nên $y=-2x+3$ là tiệm cận xiên (khi $x\to +\infty$); $a=\lim\limits_{x\to -\infty}\dfrac{y}{x}=-2$, $b=\lim\limits_{x\to -\infty}(y+2x)=3$ nên $y=-2x+3$ là tiệm cận xiên (khi $x\to -\infty$); $\lim\limits_{x\to {1}^{+}}y=+\infty$ nên $x=1$ là tiệm cận đứng.}]{Tiệm cận xiên của đồ thị hàm số $y=-2x+3+\dfrac{3}{x-1}$ là đường thẳng}
  {$y=-2x+3$.}
  {$y=-2x-3$.}{$y=2x+3$.}{$y=-2x+4$.}

\cauTN[id=C1-B3-D03-NB-002,muc=NB,loigiai={Xét hàm số $y=\dfrac{1}{2}x+4+\dfrac{1}{x+2}$. Chia đa thức: $y=\dfrac{1}{2}x+4+\dfrac{2}{2 x + 4}$; $a=\lim\limits_{x\to +\infty}\dfrac{y}{x}=\dfrac{1}{2}$, $b=\lim\limits_{x\to +\infty}(y-\dfrac{1}{2}x)=4$ nên $y=\dfrac{1}{2}x+4$ là tiệm cận xiên (khi $x\to +\infty$); $a=\lim\limits_{x\to -\infty}\dfrac{y}{x}=\dfrac{1}{2}$, $b=\lim\limits_{x\to -\infty}(y-\dfrac{1}{2}x)=4$ nên $y=\dfrac{1}{2}x+4$ là tiệm cận xiên (khi $x\to -\infty$); $\lim\limits_{x\to {-2}^{+}}y=+\infty$ nên $x=-2$ là tiệm cận đứng.}]{Tiệm cận xiên của đồ thị hàm số $y=\dfrac{1}{2}x+4+\dfrac{1}{x+2}$ đi qua điểm nào dưới đây?}
  {$M(-1;\dfrac{7}{2})$.}
  {$M(-1;\dfrac{9}{2})$.}{$M(-1;\dfrac{5}{2})$.}{$M(-1;\dfrac{11}{2})$.}

\cauTN[id=C1-B3-D03-NB-003,muc=NB,loigiai={Xét hàm số $y=\dfrac{3x^{2}+4x-2}{-x}$. Chia đa thức: $y=-3x-4+\dfrac{2}{x}$; $a=\lim\limits_{x\to +\infty}\dfrac{y}{x}=-3$, $b=\lim\limits_{x\to +\infty}(y+3x)=-4$ nên $y=-3x-4$ là tiệm cận xiên (khi $x\to +\infty$); $a=\lim\limits_{x\to -\infty}\dfrac{y}{x}=-3$, $b=\lim\limits_{x\to -\infty}(y+3x)=-4$ nên $y=-3x-4$ là tiệm cận xiên (khi $x\to -\infty$); $\lim\limits_{x\to {0}^{+}}y=+\infty$ nên $x=0$ là tiệm cận đứng.}]{Tiệm cận xiên của đồ thị hàm số $y=\dfrac{3x^{2}+4x-2}{-x}$ đi qua điểm nào dưới đây?}
  {$M(2;-10)$.}
  {$M(2;-9)$.}{$M(2;-11)$.}{$M(2;-8)$.}

\cauTN[id=C1-B3-D03-NB-004,muc=NB,loigiai={Xét hàm số $y=\dfrac{3x^{2}+3x}{x-2}$. Chia đa thức: $y=3x+9+\dfrac{18}{x - 2}$; $a=\lim\limits_{x\to +\infty}\dfrac{y}{x}=3$, $b=\lim\limits_{x\to +\infty}(y-3x)=9$ nên $y=3x+9$ là tiệm cận xiên (khi $x\to +\infty$); $a=\lim\limits_{x\to -\infty}\dfrac{y}{x}=3$, $b=\lim\limits_{x\to -\infty}(y-3x)=9$ nên $y=3x+9$ là tiệm cận xiên (khi $x\to -\infty$); $\lim\limits_{x\to {2}^{+}}y=+\infty$ nên $x=2$ là tiệm cận đứng.}]{Tiệm cận xiên của đồ thị hàm số $y=\dfrac{3x^{2}+3x}{x-2}$ là đường thẳng}
  {$y=3x+9$.}
  {$y=3x-9$.}{$y=-3x+9$.}{$y=3x+10$.}

\cauTN[id=C1-B3-D03-NB-005,muc=NB,loigiai={Xét hàm số $y=\dfrac{1}{2}x+3-\dfrac{2}{x+3}$. Chia đa thức: $y=\dfrac{1}{2}x+3+\dfrac{-4}{2 x + 6}$; $a=\lim\limits_{x\to +\infty}\dfrac{y}{x}=\dfrac{1}{2}$, $b=\lim\limits_{x\to +\infty}(y-\dfrac{1}{2}x)=3$ nên $y=\dfrac{1}{2}x+3$ là tiệm cận xiên (khi $x\to +\infty$); $a=\lim\limits_{x\to -\infty}\dfrac{y}{x}=\dfrac{1}{2}$, $b=\lim\limits_{x\to -\infty}(y-\dfrac{1}{2}x)=3$ nên $y=\dfrac{1}{2}x+3$ là tiệm cận xiên (khi $x\to -\infty$); $\lim\limits_{x\to {-3}^{+}}y=-\infty$ nên $x=-3$ là tiệm cận đứng.}]{Tiệm cận xiên của đồ thị hàm số $y=\dfrac{1}{2}x+3-\dfrac{2}{x+3}$ là đường thẳng}
  {$y=\dfrac{1}{2}x+3$.}
  {$y=\dfrac{1}{2}x-3$.}{$y=-\dfrac{1}{2}x+3$.}{$y=\dfrac{1}{2}x+4$.}

\cauTN[id=C1-B3-D03-NB-006,muc=NB,loigiai={Xét hàm số $y=x+3+\dfrac{4}{x+3}$. Chia đa thức: $y=x+3+\dfrac{4}{x + 3}$; $a=\lim\limits_{x\to +\infty}\dfrac{y}{x}=1$, $b=\lim\limits_{x\to +\infty}(y-x)=3$ nên $y=x+3$ là tiệm cận xiên (khi $x\to +\infty$); $a=\lim\limits_{x\to -\infty}\dfrac{y}{x}=1$, $b=\lim\limits_{x\to -\infty}(y-x)=3$ nên $y=x+3$ là tiệm cận xiên (khi $x\to -\infty$); $\lim\limits_{x\to {-3}^{+}}y=+\infty$ nên $x=-3$ là tiệm cận đứng.}]{Tiệm cận xiên của đồ thị hàm số $y=x+3+\dfrac{4}{x+3}$ là đường thẳng}
  {$y=x+3$.}
  {$y=x-3$.}{$y=-x+3$.}{$y=x+4$.}

\cauTN[id=C1-B3-D03-NB-007,muc=NB,loigiai={Xét hàm số $y=x-1+\dfrac{1}{x+2}$. Chia đa thức: $y=x-1+\dfrac{1}{x + 2}$; $a=\lim\limits_{x\to +\infty}\dfrac{y}{x}=1$, $b=\lim\limits_{x\to +\infty}(y-x)=-1$ nên $y=x-1$ là tiệm cận xiên (khi $x\to +\infty$); $a=\lim\limits_{x\to -\infty}\dfrac{y}{x}=1$, $b=\lim\limits_{x\to -\infty}(y-x)=-1$ nên $y=x-1$ là tiệm cận xiên (khi $x\to -\infty$); $\lim\limits_{x\to {-2}^{+}}y=+\infty$ nên $x=-2$ là tiệm cận đứng.}]{Tiệm cận xiên của đồ thị hàm số $y=x-1+\dfrac{1}{x+2}$ đi qua điểm nào dưới đây?}
  {$M(3;2)$.}
  {$M(3;3)$.}{$M(3;1)$.}{$M(3;4)$.}

\cauTN[id=C1-B3-D03-NB-008,muc=NB,loigiai={Xét hàm số $y=\dfrac{3x^{2}-2x-4}{-x-1}$. Chia đa thức: $y=-3x+5+\dfrac{1}{- x - 1}$; $a=\lim\limits_{x\to +\infty}\dfrac{y}{x}=-3$, $b=\lim\limits_{x\to +\infty}(y+3x)=5$ nên $y=-3x+5$ là tiệm cận xiên (khi $x\to +\infty$); $a=\lim\limits_{x\to -\infty}\dfrac{y}{x}=-3$, $b=\lim\limits_{x\to -\infty}(y+3x)=5$ nên $y=-3x+5$ là tiệm cận xiên (khi $x\to -\infty$); $\lim\limits_{x\to {-1}^{+}}y=-\infty$ nên $x=-1$ là tiệm cận đứng.}]{Tiệm cận xiên của đồ thị hàm số $y=\dfrac{3x^{2}-2x-4}{-x-1}$ là đường thẳng}
  {$y=-3x+5$.}
  {$y=-3x-5$.}{$y=3x+5$.}{$y=-3x+6$.}

\cauTN[id=C1-B3-D03-NB-009,muc=NB,loigiai={Xét hàm số $y=\dfrac{-2x^{2}-2x+5}{x-2}$. Chia đa thức: $y=-2x-6+\dfrac{-7}{x - 2}$; $a=\lim\limits_{x\to +\infty}\dfrac{y}{x}=-2$, $b=\lim\limits_{x\to +\infty}(y+2x)=-6$ nên $y=-2x-6$ là tiệm cận xiên (khi $x\to +\infty$); $a=\lim\limits_{x\to -\infty}\dfrac{y}{x}=-2$, $b=\lim\limits_{x\to -\infty}(y+2x)=-6$ nên $y=-2x-6$ là tiệm cận xiên (khi $x\to -\infty$); $\lim\limits_{x\to {2}^{+}}y=-\infty$ nên $x=2$ là tiệm cận đứng.}]{Tiệm cận xiên của đồ thị hàm số $y=\dfrac{-2x^{2}-2x+5}{x-2}$ là đường thẳng}
  {$y=-2x-6$.}
  {$y=-2x+6$.}{$y=2x-6$.}{$y=-2x-5$.}

\cauTN[id=C1-B3-D03-NB-010,muc=NB,loigiai={Xét hàm số $y=-2x+4+\dfrac{3}{x+3}$. Chia đa thức: $y=-2x+4+\dfrac{3}{x + 3}$; $a=\lim\limits_{x\to +\infty}\dfrac{y}{x}=-2$, $b=\lim\limits_{x\to +\infty}(y+2x)=4$ nên $y=-2x+4$ là tiệm cận xiên (khi $x\to +\infty$); $a=\lim\limits_{x\to -\infty}\dfrac{y}{x}=-2$, $b=\lim\limits_{x\to -\infty}(y+2x)=4$ nên $y=-2x+4$ là tiệm cận xiên (khi $x\to -\infty$); $\lim\limits_{x\to {-3}^{+}}y=+\infty$ nên $x=-3$ là tiệm cận đứng.}]{Tiệm cận xiên của đồ thị hàm số $y=-2x+4+\dfrac{3}{x+3}$ là đường thẳng}
  {$y=-2x+4$.}
  {$y=-2x-4$.}{$y=2x+4$.}{$y=-2x+5$.}

\cauTN[id=C1-B3-D03-NB-011,muc=NB,loigiai={Xét hàm số $y=3x+4-\dfrac{1}{x-2}$. Chia đa thức: $y=3x+4+\dfrac{-1}{x - 2}$; $a=\lim\limits_{x\to +\infty}\dfrac{y}{x}=3$, $b=\lim\limits_{x\to +\infty}(y-3x)=4$ nên $y=3x+4$ là tiệm cận xiên (khi $x\to +\infty$); $a=\lim\limits_{x\to -\infty}\dfrac{y}{x}=3$, $b=\lim\limits_{x\to -\infty}(y-3x)=4$ nên $y=3x+4$ là tiệm cận xiên (khi $x\to -\infty$); $\lim\limits_{x\to {2}^{+}}y=-\infty$ nên $x=2$ là tiệm cận đứng.}]{Tiệm cận xiên của đồ thị hàm số $y=3x+4-\dfrac{1}{x-2}$ là đường thẳng}
  {$y=3x+4$.}
  {$y=3x-4$.}{$y=-3x+4$.}{$y=3x+5$.}

\cauTN[id=C1-B3-D03-NB-012,muc=NB,loigiai={Xét hàm số $y=\dfrac{3x^{2}-4x+2}{x-4}$. Chia đa thức: $y=3x+8+\dfrac{34}{x - 4}$; $a=\lim\limits_{x\to +\infty}\dfrac{y}{x}=3$, $b=\lim\limits_{x\to +\infty}(y-3x)=8$ nên $y=3x+8$ là tiệm cận xiên (khi $x\to +\infty$); $a=\lim\limits_{x\to -\infty}\dfrac{y}{x}=3$, $b=\lim\limits_{x\to -\infty}(y-3x)=8$ nên $y=3x+8$ là tiệm cận xiên (khi $x\to -\infty$); $\lim\limits_{x\to {4}^{+}}y=+\infty$ nên $x=4$ là tiệm cận đứng.}]{Tiệm cận xiên của đồ thị hàm số $y=\dfrac{3x^{2}-4x+2}{x-4}$ là đường thẳng}
  {$y=3x+8$.}
  {$y=3x-8$.}{$y=-3x+8$.}{$y=3x+9$.}

\cauTN[id=C1-B3-D03-NB-013,muc=NB,loigiai={Xét hàm số $y=\dfrac{x^{2}+4x}{x-2}$. Chia đa thức: $y=x+6+\dfrac{12}{x - 2}$; $a=\lim\limits_{x\to +\infty}\dfrac{y}{x}=1$, $b=\lim\limits_{x\to +\infty}(y-x)=6$ nên $y=x+6$ là tiệm cận xiên (khi $x\to +\infty$); $a=\lim\limits_{x\to -\infty}\dfrac{y}{x}=1$, $b=\lim\limits_{x\to -\infty}(y-x)=6$ nên $y=x+6$ là tiệm cận xiên (khi $x\to -\infty$); $\lim\limits_{x\to {2}^{+}}y=+\infty$ nên $x=2$ là tiệm cận đứng.}]{Tiệm cận xiên của đồ thị hàm số $y=\dfrac{x^{2}+4x}{x-2}$ là đường thẳng}
  {$y=x+6$.}
  {$y=x-6$.}{$y=-x+6$.}{$y=x+7$.}

\cauTN[id=C1-B3-D03-NB-014,muc=NB,loigiai={Xét hàm số $y=\dfrac{2x^{2}-2x+5}{2x-1}$. Chia đa thức: $y=x-\dfrac{1}{2}+\dfrac{\dfrac{9}{2}}{2 x - 1}$; $a=\lim\limits_{x\to +\infty}\dfrac{y}{x}=1$, $b=\lim\limits_{x\to +\infty}(y-x)=-\dfrac{1}{2}$ nên $y=x-\dfrac{1}{2}$ là tiệm cận xiên (khi $x\to +\infty$); $a=\lim\limits_{x\to -\infty}\dfrac{y}{x}=1$, $b=\lim\limits_{x\to -\infty}(y-x)=-\dfrac{1}{2}$ nên $y=x-\dfrac{1}{2}$ là tiệm cận xiên (khi $x\to -\infty$); $\lim\limits_{x\to {\dfrac{1}{2}}^{+}}y=+\infty$ nên $x=\dfrac{1}{2}$ là tiệm cận đứng.}]{Tiệm cận xiên của đồ thị hàm số $y=\dfrac{2x^{2}-2x+5}{2x-1}$ là đường thẳng}
  {$y=x-\dfrac{1}{2}$.}
  {$y=x+\dfrac{1}{2}$.}{$y=-x-\dfrac{1}{2}$.}{$y=x-\dfrac{3}{2}$.}

\cauTN[id=C1-B3-D03-NB-015,muc=NB,loigiai={Xét hàm số $y=\dfrac{-2x^{2}+6}{x-2}$. Chia đa thức: $y=-2x-4+\dfrac{-2}{x - 2}$; $a=\lim\limits_{x\to +\infty}\dfrac{y}{x}=-2$, $b=\lim\limits_{x\to +\infty}(y+2x)=-4$ nên $y=-2x-4$ là tiệm cận xiên (khi $x\to +\infty$); $a=\lim\limits_{x\to -\infty}\dfrac{y}{x}=-2$, $b=\lim\limits_{x\to -\infty}(y+2x)=-4$ nên $y=-2x-4$ là tiệm cận xiên (khi $x\to -\infty$); $\lim\limits_{x\to {2}^{+}}y=-\infty$ nên $x=2$ là tiệm cận đứng.}]{Tiệm cận xiên của đồ thị hàm số $y=\dfrac{-2x^{2}+6}{x-2}$ đi qua điểm nào dưới đây?}
  {$M(3;-10)$.}
  {$M(3;-9)$.}{$M(3;-11)$.}{$M(3;-8)$.}

\cauTN[id=C1-B3-D03-NB-016,muc=NB,loigiai={Xét hàm số $y=-2x-3-\dfrac{2}{x}$. Chia đa thức: $y=-2x-3+\dfrac{-2}{x}$; $a=\lim\limits_{x\to +\infty}\dfrac{y}{x}=-2$, $b=\lim\limits_{x\to +\infty}(y+2x)=-3$ nên $y=-2x-3$ là tiệm cận xiên (khi $x\to +\infty$); $a=\lim\limits_{x\to -\infty}\dfrac{y}{x}=-2$, $b=\lim\limits_{x\to -\infty}(y+2x)=-3$ nên $y=-2x-3$ là tiệm cận xiên (khi $x\to -\infty$); $\lim\limits_{x\to {0}^{+}}y=-\infty$ nên $x=0$ là tiệm cận đứng.}]{Tiệm cận xiên của đồ thị hàm số $y=-2x-3-\dfrac{2}{x}$ là đường thẳng}
  {$y=-2x-3$.}
  {$y=-2x+3$.}{$y=2x-3$.}{$y=-2x-2$.}

\cauTN[id=C1-B3-D03-NB-017,muc=NB,loigiai={Xét hàm số $y=-x+4+\dfrac{2}{x-1}$. Chia đa thức: $y=-x+4+\dfrac{2}{x - 1}$; $a=\lim\limits_{x\to +\infty}\dfrac{y}{x}=-1$, $b=\lim\limits_{x\to +\infty}(y+x)=4$ nên $y=-x+4$ là tiệm cận xiên (khi $x\to +\infty$); $a=\lim\limits_{x\to -\infty}\dfrac{y}{x}=-1$, $b=\lim\limits_{x\to -\infty}(y+x)=4$ nên $y=-x+4$ là tiệm cận xiên (khi $x\to -\infty$); $\lim\limits_{x\to {1}^{+}}y=+\infty$ nên $x=1$ là tiệm cận đứng.}]{Tiệm cận xiên của đồ thị hàm số $y=-x+4+\dfrac{2}{x-1}$ đi qua điểm nào dưới đây?}
  {$M(-2;6)$.}
  {$M(-2;7)$.}{$M(-2;5)$.}{$M(-2;8)$.}

\cauTN[id=C1-B3-D03-NB-018,muc=NB,loigiai={Xét hàm số $y=3x-4+\dfrac{1}{x}$. Chia đa thức: $y=3x-4+\dfrac{1}{x}$; $a=\lim\limits_{x\to +\infty}\dfrac{y}{x}=3$, $b=\lim\limits_{x\to +\infty}(y-3x)=-4$ nên $y=3x-4$ là tiệm cận xiên (khi $x\to +\infty$); $a=\lim\limits_{x\to -\infty}\dfrac{y}{x}=3$, $b=\lim\limits_{x\to -\infty}(y-3x)=-4$ nên $y=3x-4$ là tiệm cận xiên (khi $x\to -\infty$); $\lim\limits_{x\to {0}^{+}}y=+\infty$ nên $x=0$ là tiệm cận đứng.}]{Tiệm cận xiên của đồ thị hàm số $y=3x-4+\dfrac{1}{x}$ là đường thẳng}
  {$y=3x-4$.}
  {$y=3x+4$.}{$y=-3x-4$.}{$y=3x-3$.}

\cauTN[id=C1-B3-D03-NB-019,muc=NB,loigiai={Xét hàm số $y=-2x+3+\dfrac{2}{x-2}$. Chia đa thức: $y=-2x+3+\dfrac{2}{x - 2}$; $a=\lim\limits_{x\to +\infty}\dfrac{y}{x}=-2$, $b=\lim\limits_{x\to +\infty}(y+2x)=3$ nên $y=-2x+3$ là tiệm cận xiên (khi $x\to +\infty$); $a=\lim\limits_{x\to -\infty}\dfrac{y}{x}=-2$, $b=\lim\limits_{x\to -\infty}(y+2x)=3$ nên $y=-2x+3$ là tiệm cận xiên (khi $x\to -\infty$); $\lim\limits_{x\to {2}^{+}}y=+\infty$ nên $x=2$ là tiệm cận đứng.}]{Tiệm cận xiên của đồ thị hàm số $y=-2x+3+\dfrac{2}{x-2}$ là đường thẳng}
  {$y=-2x+3$.}
  {$y=-2x-3$.}{$y=2x+3$.}{$y=-2x+4$.}

\cauTN[id=C1-B3-D03-NB-020,muc=NB,loigiai={Xét hàm số $y=x-4-\dfrac{1}{x+3}$. Chia đa thức: $y=x-4+\dfrac{-1}{x + 3}$; $a=\lim\limits_{x\to +\infty}\dfrac{y}{x}=1$, $b=\lim\limits_{x\to +\infty}(y-x)=-4$ nên $y=x-4$ là tiệm cận xiên (khi $x\to +\infty$); $a=\lim\limits_{x\to -\infty}\dfrac{y}{x}=1$, $b=\lim\limits_{x\to -\infty}(y-x)=-4$ nên $y=x-4$ là tiệm cận xiên (khi $x\to -\infty$); $\lim\limits_{x\to {-3}^{+}}y=-\infty$ nên $x=-3$ là tiệm cận đứng.}]{Tiệm cận xiên của đồ thị hàm số $y=x-4-\dfrac{1}{x+3}$ là đường thẳng}
  {$y=x-4$.}
  {$y=x+4$.}{$y=-x-4$.}{$y=x-3$.}

\cauTN[id=C1-B3-D03-NB-021,muc=NB,loigiai={Xét hàm số $y=\dfrac{-2x^{2}-2x-6}{-x+2}$. Chia đa thức: $y=2x+6+\dfrac{-18}{2 - x}$; $a=\lim\limits_{x\to +\infty}\dfrac{y}{x}=2$, $b=\lim\limits_{x\to +\infty}(y-2x)=6$ nên $y=2x+6$ là tiệm cận xiên (khi $x\to +\infty$); $a=\lim\limits_{x\to -\infty}\dfrac{y}{x}=2$, $b=\lim\limits_{x\to -\infty}(y-2x)=6$ nên $y=2x+6$ là tiệm cận xiên (khi $x\to -\infty$); $\lim\limits_{x\to {2}^{+}}y=+\infty$ nên $x=2$ là tiệm cận đứng.}]{Tiệm cận xiên của đồ thị hàm số $y=\dfrac{-2x^{2}-2x-6}{-x+2}$ là đường thẳng}
  {$y=2x+6$.}
  {$y=2x-6$.}{$y=-2x+6$.}{$y=2x+7$.}

\cauTN[id=C1-B3-D03-NB-022,muc=NB,loigiai={Xét hàm số $y=\dfrac{x^{2}+5x-1}{-x+1}$. Chia đa thức: $y=-x-6+\dfrac{5}{1 - x}$; $a=\lim\limits_{x\to +\infty}\dfrac{y}{x}=-1$, $b=\lim\limits_{x\to +\infty}(y+x)=-6$ nên $y=-x-6$ là tiệm cận xiên (khi $x\to +\infty$); $a=\lim\limits_{x\to -\infty}\dfrac{y}{x}=-1$, $b=\lim\limits_{x\to -\infty}(y+x)=-6$ nên $y=-x-6$ là tiệm cận xiên (khi $x\to -\infty$); $\lim\limits_{x\to {1}^{+}}y=-\infty$ nên $x=1$ là tiệm cận đứng.}]{Tiệm cận xiên của đồ thị hàm số $y=\dfrac{x^{2}+5x-1}{-x+1}$ là đường thẳng}
  {$y=-x-6$.}
  {$y=-x+6$.}{$y=x-6$.}{$y=-x-5$.}

\cauTN[id=C1-B3-D03-NB-023,muc=NB,loigiai={Xét hàm số $y=x+3+\dfrac{3}{x+3}$. Chia đa thức: $y=x+3+\dfrac{3}{x + 3}$; $a=\lim\limits_{x\to +\infty}\dfrac{y}{x}=1$, $b=\lim\limits_{x\to +\infty}(y-x)=3$ nên $y=x+3$ là tiệm cận xiên (khi $x\to +\infty$); $a=\lim\limits_{x\to -\infty}\dfrac{y}{x}=1$, $b=\lim\limits_{x\to -\infty}(y-x)=3$ nên $y=x+3$ là tiệm cận xiên (khi $x\to -\infty$); $\lim\limits_{x\to {-3}^{+}}y=+\infty$ nên $x=-3$ là tiệm cận đứng.}]{Tiệm cận xiên của đồ thị hàm số $y=x+3+\dfrac{3}{x+3}$ là đường thẳng}
  {$y=x+3$.}
  {$y=x-3$.}{$y=-x+3$.}{$y=x+4$.}

\cauTN[id=C1-B3-D03-NB-024,muc=NB,loigiai={Xét hàm số $y=\dfrac{1}{2}x+3+\dfrac{4}{x+1}$. Chia đa thức: $y=\dfrac{1}{2}x+3+\dfrac{8}{2 x + 2}$; $a=\lim\limits_{x\to +\infty}\dfrac{y}{x}=\dfrac{1}{2}$, $b=\lim\limits_{x\to +\infty}(y-\dfrac{1}{2}x)=3$ nên $y=\dfrac{1}{2}x+3$ là tiệm cận xiên (khi $x\to +\infty$); $a=\lim\limits_{x\to -\infty}\dfrac{y}{x}=\dfrac{1}{2}$, $b=\lim\limits_{x\to -\infty}(y-\dfrac{1}{2}x)=3$ nên $y=\dfrac{1}{2}x+3$ là tiệm cận xiên (khi $x\to -\infty$); $\lim\limits_{x\to {-1}^{+}}y=+\infty$ nên $x=-1$ là tiệm cận đứng.}]{Tiệm cận xiên của đồ thị hàm số $y=\dfrac{1}{2}x+3+\dfrac{4}{x+1}$ là đường thẳng}
  {$y=\dfrac{1}{2}x+3$.}
  {$y=\dfrac{1}{2}x-3$.}{$y=-\dfrac{1}{2}x+3$.}{$y=\dfrac{1}{2}x+4$.}

\cauTN[id=C1-B3-D03-NB-025,muc=NB,loigiai={Xét hàm số $y=3x-2-\dfrac{2}{x+3}$. Chia đa thức: $y=3x-2+\dfrac{-2}{x + 3}$; $a=\lim\limits_{x\to +\infty}\dfrac{y}{x}=3$, $b=\lim\limits_{x\to +\infty}(y-3x)=-2$ nên $y=3x-2$ là tiệm cận xiên (khi $x\to +\infty$); $a=\lim\limits_{x\to -\infty}\dfrac{y}{x}=3$, $b=\lim\limits_{x\to -\infty}(y-3x)=-2$ nên $y=3x-2$ là tiệm cận xiên (khi $x\to -\infty$); $\lim\limits_{x\to {-3}^{+}}y=-\infty$ nên $x=-3$ là tiệm cận đứng.}]{Tiệm cận xiên của đồ thị hàm số $y=3x-2-\dfrac{2}{x+3}$ là đường thẳng}
  {$y=3x-2$.}
  {$y=3x+2$.}{$y=-3x-2$.}{$y=3x-1$.}

\cauTN[id=C1-B3-D03-NB-026,muc=NB,loigiai={Xét hàm số $y=\dfrac{x^{2}+x}{x-3}$. Chia đa thức: $y=x+4+\dfrac{12}{x - 3}$; $a=\lim\limits_{x\to +\infty}\dfrac{y}{x}=1$, $b=\lim\limits_{x\to +\infty}(y-x)=4$ nên $y=x+4$ là tiệm cận xiên (khi $x\to +\infty$); $a=\lim\limits_{x\to -\infty}\dfrac{y}{x}=1$, $b=\lim\limits_{x\to -\infty}(y-x)=4$ nên $y=x+4$ là tiệm cận xiên (khi $x\to -\infty$); $\lim\limits_{x\to {3}^{+}}y=+\infty$ nên $x=3$ là tiệm cận đứng.}]{Tiệm cận xiên của đồ thị hàm số $y=\dfrac{x^{2}+x}{x-3}$ là đường thẳng}
  {$y=x+4$.}
  {$y=x-4$.}{$y=-x+4$.}{$y=x+5$.}

\cauTN[id=C1-B3-D03-NB-027,muc=NB,loigiai={Xét hàm số $y=3x+\dfrac{1}{x+3}$. Chia đa thức: $y=3x+\dfrac{1}{x + 3}$; $a=\lim\limits_{x\to +\infty}\dfrac{y}{x}=3$, $b=\lim\limits_{x\to +\infty}(y-3x)=0$ nên $y=3x$ là tiệm cận xiên (khi $x\to +\infty$); $a=\lim\limits_{x\to -\infty}\dfrac{y}{x}=3$, $b=\lim\limits_{x\to -\infty}(y-3x)=0$ nên $y=3x$ là tiệm cận xiên (khi $x\to -\infty$); $\lim\limits_{x\to {-3}^{+}}y=+\infty$ nên $x=-3$ là tiệm cận đứng.}]{Tiệm cận xiên của đồ thị hàm số $y=3x+\dfrac{1}{x+3}$ là đường thẳng}
  {$y=3x$.}
  {$y=-3x$.}{$y=3x+1$.}{$y=3x-1$.}

\cauTN[id=C1-B3-D03-NB-028,muc=NB,loigiai={Xét hàm số $y=2x-\dfrac{1}{x+1}$. Chia đa thức: $y=2x+\dfrac{-1}{x + 1}$; $a=\lim\limits_{x\to +\infty}\dfrac{y}{x}=2$, $b=\lim\limits_{x\to +\infty}(y-2x)=0$ nên $y=2x$ là tiệm cận xiên (khi $x\to +\infty$); $a=\lim\limits_{x\to -\infty}\dfrac{y}{x}=2$, $b=\lim\limits_{x\to -\infty}(y-2x)=0$ nên $y=2x$ là tiệm cận xiên (khi $x\to -\infty$); $\lim\limits_{x\to {-1}^{+}}y=-\infty$ nên $x=-1$ là tiệm cận đứng.}]{Tiệm cận xiên của đồ thị hàm số $y=2x-\dfrac{1}{x+1}$ đi qua điểm nào dưới đây?}
  {$M(-2;-4)$.}
  {$M(-2;-3)$.}{$M(-2;-5)$.}{$M(-2;-2)$.}

\cauTN[id=C1-B3-D03-NB-029,muc=NB,loigiai={Xét hàm số $y=3x-3+\dfrac{3}{x-3}$. Chia đa thức: $y=3x-3+\dfrac{3}{x - 3}$; $a=\lim\limits_{x\to +\infty}\dfrac{y}{x}=3$, $b=\lim\limits_{x\to +\infty}(y-3x)=-3$ nên $y=3x-3$ là tiệm cận xiên (khi $x\to +\infty$); $a=\lim\limits_{x\to -\infty}\dfrac{y}{x}=3$, $b=\lim\limits_{x\to -\infty}(y-3x)=-3$ nên $y=3x-3$ là tiệm cận xiên (khi $x\to -\infty$); $\lim\limits_{x\to {3}^{+}}y=+\infty$ nên $x=3$ là tiệm cận đứng.}]{Tiệm cận xiên của đồ thị hàm số $y=3x-3+\dfrac{3}{x-3}$ là đường thẳng}
  {$y=3x-3$.}
  {$y=3x+3$.}{$y=-3x-3$.}{$y=3x-2$.}

\cauTN[id=C1-B3-D03-NB-030,muc=NB,loigiai={Xét hàm số $y=\dfrac{2x^{2}-4x-1}{x-2}$. Chia đa thức: $y=2x+\dfrac{-1}{x - 2}$; $a=\lim\limits_{x\to +\infty}\dfrac{y}{x}=2$, $b=\lim\limits_{x\to +\infty}(y-2x)=0$ nên $y=2x$ là tiệm cận xiên (khi $x\to +\infty$); $a=\lim\limits_{x\to -\infty}\dfrac{y}{x}=2$, $b=\lim\limits_{x\to -\infty}(y-2x)=0$ nên $y=2x$ là tiệm cận xiên (khi $x\to -\infty$); $\lim\limits_{x\to {2}^{+}}y=-\infty$ nên $x=2$ là tiệm cận đứng.}]{Tiệm cận xiên của đồ thị hàm số $y=\dfrac{2x^{2}-4x-1}{x-2}$ đi qua điểm nào dưới đây?}
  {$M(1;2)$.}
  {$M(1;3)$.}{$M(1;1)$.}{$M(1;4)$.}

\cauTN[id=C1-B3-D03-NB-031,muc=NB,loigiai={Xét hàm số $y=3x+2-\dfrac{1}{x-3}$. Chia đa thức: $y=3x+2+\dfrac{-1}{x - 3}$; $a=\lim\limits_{x\to +\infty}\dfrac{y}{x}=3$, $b=\lim\limits_{x\to +\infty}(y-3x)=2$ nên $y=3x+2$ là tiệm cận xiên (khi $x\to +\infty$); $a=\lim\limits_{x\to -\infty}\dfrac{y}{x}=3$, $b=\lim\limits_{x\to -\infty}(y-3x)=2$ nên $y=3x+2$ là tiệm cận xiên (khi $x\to -\infty$); $\lim\limits_{x\to {3}^{+}}y=-\infty$ nên $x=3$ là tiệm cận đứng.}]{Tiệm cận xiên của đồ thị hàm số $y=3x+2-\dfrac{1}{x-3}$ là đường thẳng}
  {$y=3x+2$.}
  {$y=3x-2$.}{$y=-3x+2$.}{$y=3x+3$.}

\cauTN[id=C1-B3-D03-NB-032,muc=NB,loigiai={Xét hàm số $y=3x-1+\dfrac{3}{x-3}$. Chia đa thức: $y=3x-1+\dfrac{3}{x - 3}$; $a=\lim\limits_{x\to +\infty}\dfrac{y}{x}=3$, $b=\lim\limits_{x\to +\infty}(y-3x)=-1$ nên $y=3x-1$ là tiệm cận xiên (khi $x\to +\infty$); $a=\lim\limits_{x\to -\infty}\dfrac{y}{x}=3$, $b=\lim\limits_{x\to -\infty}(y-3x)=-1$ nên $y=3x-1$ là tiệm cận xiên (khi $x\to -\infty$); $\lim\limits_{x\to {3}^{+}}y=+\infty$ nên $x=3$ là tiệm cận đứng.}]{Tiệm cận xiên của đồ thị hàm số $y=3x-1+\dfrac{3}{x-3}$ đi qua điểm nào dưới đây?}
  {$M(-2;-7)$.}
  {$M(-2;-6)$.}{$M(-2;-8)$.}{$M(-2;-5)$.}

\cauTN[id=C1-B3-D03-NB-033,muc=NB,loigiai={Xét hàm số $y=\dfrac{2x^{2}-x+4}{x+1}$. Chia đa thức: $y=2x-3+\dfrac{7}{x + 1}$; $a=\lim\limits_{x\to +\infty}\dfrac{y}{x}=2$, $b=\lim\limits_{x\to +\infty}(y-2x)=-3$ nên $y=2x-3$ là tiệm cận xiên (khi $x\to +\infty$); $a=\lim\limits_{x\to -\infty}\dfrac{y}{x}=2$, $b=\lim\limits_{x\to -\infty}(y-2x)=-3$ nên $y=2x-3$ là tiệm cận xiên (khi $x\to -\infty$); $\lim\limits_{x\to {-1}^{+}}y=+\infty$ nên $x=-1$ là tiệm cận đứng.}]{Tiệm cận xiên của đồ thị hàm số $y=\dfrac{2x^{2}-x+4}{x+1}$ là đường thẳng}
  {$y=2x-3$.}
  {$y=2x+3$.}{$y=-2x-3$.}{$y=2x-2$.}

\cauTN[id=C1-B3-D03-NB-034,muc=NB,loigiai={Xét hàm số $y=\dfrac{1}{2}x+1+\dfrac{1}{x-1}$. Chia đa thức: $y=\dfrac{1}{2}x+1+\dfrac{2}{2 x - 2}$; $a=\lim\limits_{x\to +\infty}\dfrac{y}{x}=\dfrac{1}{2}$, $b=\lim\limits_{x\to +\infty}(y-\dfrac{1}{2}x)=1$ nên $y=\dfrac{1}{2}x+1$ là tiệm cận xiên (khi $x\to +\infty$); $a=\lim\limits_{x\to -\infty}\dfrac{y}{x}=\dfrac{1}{2}$, $b=\lim\limits_{x\to -\infty}(y-\dfrac{1}{2}x)=1$ nên $y=\dfrac{1}{2}x+1$ là tiệm cận xiên (khi $x\to -\infty$); $\lim\limits_{x\to {1}^{+}}y=+\infty$ nên $x=1$ là tiệm cận đứng.}]{Tiệm cận xiên của đồ thị hàm số $y=\dfrac{1}{2}x+1+\dfrac{1}{x-1}$ là đường thẳng}
  {$y=\dfrac{1}{2}x+1$.}
  {$y=\dfrac{1}{2}x-1$.}{$y=-\dfrac{1}{2}x+1$.}{$y=\dfrac{1}{2}x+2$.}

\cauTN[id=C1-B3-D03-NB-035,muc=NB,loigiai={Xét hàm số $y=3x+2+\dfrac{4}{x+3}$. Chia đa thức: $y=3x+2+\dfrac{4}{x + 3}$; $a=\lim\limits_{x\to +\infty}\dfrac{y}{x}=3$, $b=\lim\limits_{x\to +\infty}(y-3x)=2$ nên $y=3x+2$ là tiệm cận xiên (khi $x\to +\infty$); $a=\lim\limits_{x\to -\infty}\dfrac{y}{x}=3$, $b=\lim\limits_{x\to -\infty}(y-3x)=2$ nên $y=3x+2$ là tiệm cận xiên (khi $x\to -\infty$); $\lim\limits_{x\to {-3}^{+}}y=+\infty$ nên $x=-3$ là tiệm cận đứng.}]{Tiệm cận xiên của đồ thị hàm số $y=3x+2+\dfrac{4}{x+3}$ là đường thẳng}
  {$y=3x+2$.}
  {$y=3x-2$.}{$y=-3x+2$.}{$y=3x+3$.}

\cauTN[id=C1-B3-D03-NB-036,muc=NB,loigiai={Xét hàm số $y=\dfrac{2x^{2}-4x-1}{-x-1}$. Chia đa thức: $y=-2x+6+\dfrac{5}{- x - 1}$; $a=\lim\limits_{x\to +\infty}\dfrac{y}{x}=-2$, $b=\lim\limits_{x\to +\infty}(y+2x)=6$ nên $y=-2x+6$ là tiệm cận xiên (khi $x\to +\infty$); $a=\lim\limits_{x\to -\infty}\dfrac{y}{x}=-2$, $b=\lim\limits_{x\to -\infty}(y+2x)=6$ nên $y=-2x+6$ là tiệm cận xiên (khi $x\to -\infty$); $\lim\limits_{x\to {-1}^{+}}y=-\infty$ nên $x=-1$ là tiệm cận đứng.}]{Tiệm cận xiên của đồ thị hàm số $y=\dfrac{2x^{2}-4x-1}{-x-1}$ là đường thẳng}
  {$y=-2x+6$.}
  {$y=-2x-6$.}{$y=2x+6$.}{$y=-2x+7$.}

\cauDS[id=C1-B3-D03-NB-037,muc=NB,loigiai={Xét hàm số $y=\dfrac{3x^{2}-5x-6}{-x+4}$. Chia đa thức: $y=-3x-7+\dfrac{22}{4 - x}$; $a=\lim\limits_{x\to +\infty}\dfrac{y}{x}=-3$, $b=\lim\limits_{x\to +\infty}(y+3x)=-7$ nên $y=-3x-7$ là tiệm cận xiên (khi $x\to +\infty$); $a=\lim\limits_{x\to -\infty}\dfrac{y}{x}=-3$, $b=\lim\limits_{x\to -\infty}(y+3x)=-7$ nên $y=-3x-7$ là tiệm cận xiên (khi $x\to -\infty$); $\lim\limits_{x\to {4}^{+}}y=-\infty$ nên $x=4$ là tiệm cận đứng. a) Sai: tiệm cận xiên là $y=-3x-7$. b) Đúng. c) Sai: có $2$ đường tiệm cận. d) Sai: giới hạn của $y$ tại vô cực không hữu hạn.}]{Cho hàm số $y=\dfrac{3x^{2}-5x-6}{-x+4}$. Xét tính đúng sai của các mệnh đề sau:}{%
  \y{Đường thẳng $y=-3x$ là tiệm cận xiên của đồ thị hàm số.}{s}
  \y{Đồ thị hàm số có tiệm cận đứng $x=4$.}{d}
  \y{Đồ thị hàm số có tất cả $3$ đường tiệm cận.}{s}
  \y{Đồ thị hàm số có tiệm cận ngang.}{s}}

\cauDS[id=C1-B3-D03-NB-038,muc=NB,loigiai={Xét hàm số $y=x-4-\dfrac{1}{x+2}$. Chia đa thức: $y=x-4+\dfrac{-1}{x + 2}$; $a=\lim\limits_{x\to +\infty}\dfrac{y}{x}=1$, $b=\lim\limits_{x\to +\infty}(y-x)=-4$ nên $y=x-4$ là tiệm cận xiên (khi $x\to +\infty$); $a=\lim\limits_{x\to -\infty}\dfrac{y}{x}=1$, $b=\lim\limits_{x\to -\infty}(y-x)=-4$ nên $y=x-4$ là tiệm cận xiên (khi $x\to -\infty$); $\lim\limits_{x\to {-2}^{+}}y=-\infty$ nên $x=-2$ là tiệm cận đứng. a) Đúng. b) Sai. c) Đúng. d) Đúng.}]{Cho hàm số $y=x-4-\dfrac{1}{x+2}$. Xét tính đúng sai của các mệnh đề sau:}{%
  \y{Đồ thị hàm số có tiệm cận đứng $x=-2$.}{d}
  \y{Đường thẳng $y=x-2$ là tiệm cận xiên của đồ thị hàm số.}{s}
  \y{Đường thẳng $y=x-4$ là tiệm cận xiên của đồ thị hàm số.}{d}
  \y{Đồ thị hàm số có tất cả $2$ đường tiệm cận.}{d}}

\cauDS[id=C1-B3-D03-NB-039,muc=NB,loigiai={Xét hàm số $y=-2x+2-\dfrac{1}{x-1}$. Chia đa thức: $y=-2x+2+\dfrac{-1}{x - 1}$; $a=\lim\limits_{x\to +\infty}\dfrac{y}{x}=-2$, $b=\lim\limits_{x\to +\infty}(y+2x)=2$ nên $y=-2x+2$ là tiệm cận xiên (khi $x\to +\infty$); $a=\lim\limits_{x\to -\infty}\dfrac{y}{x}=-2$, $b=\lim\limits_{x\to -\infty}(y+2x)=2$ nên $y=-2x+2$ là tiệm cận xiên (khi $x\to -\infty$); $\lim\limits_{x\to {1}^{+}}y=-\infty$ nên $x=1$ là tiệm cận đứng. a) Sai: giới hạn của $y$ tại vô cực không hữu hạn. b) Đúng. c) Đúng. d) Sai: có $2$ đường tiệm cận.}]{Cho hàm số $y=-2x+2-\dfrac{1}{x-1}$. Xét tính đúng sai của các mệnh đề sau:}{%
  \y{Đồ thị hàm số có tiệm cận ngang.}{s}
  \y{Đồ thị hàm số có tất cả $2$ đường tiệm cận.}{d}
  \y{Giao điểm của tiệm cận đứng và tiệm cận xiên là $I(1;0)$.}{d}
  \y{Đồ thị hàm số có tất cả $3$ đường tiệm cận.}{s}}

\cauDS[id=C1-B3-D03-NB-040,muc=NB,loigiai={Xét hàm số $y=2x-2+\dfrac{3}{x+2}$. Chia đa thức: $y=2x-2+\dfrac{3}{x + 2}$; $a=\lim\limits_{x\to +\infty}\dfrac{y}{x}=2$, $b=\lim\limits_{x\to +\infty}(y-2x)=-2$ nên $y=2x-2$ là tiệm cận xiên (khi $x\to +\infty$); $a=\lim\limits_{x\to -\infty}\dfrac{y}{x}=2$, $b=\lim\limits_{x\to -\infty}(y-2x)=-2$ nên $y=2x-2$ là tiệm cận xiên (khi $x\to -\infty$); $\lim\limits_{x\to {-2}^{+}}y=+\infty$ nên $x=-2$ là tiệm cận đứng. a) Sai: có $2$ đường tiệm cận. b) Đúng. c) Sai: giới hạn của $y$ tại vô cực không hữu hạn. d) Sai.}]{Cho hàm số $y=2x-2+\dfrac{3}{x+2}$. Xét tính đúng sai của các mệnh đề sau:}{%
  \y{Đồ thị hàm số có tất cả $3$ đường tiệm cận.}{s}
  \y{Đồ thị hàm số có tiệm cận đứng $x=-2$.}{d}
  \y{Đồ thị hàm số có tiệm cận ngang.}{s}
  \y{Đường thẳng $y=2x-3$ là tiệm cận xiên của đồ thị hàm số.}{s}}

\cauDS[id=C1-B3-D03-NB-041,muc=NB,loigiai={Xét hàm số $y=3x+1+\dfrac{4}{x+3}$. Chia đa thức: $y=3x+1+\dfrac{4}{x + 3}$; $a=\lim\limits_{x\to +\infty}\dfrac{y}{x}=3$, $b=\lim\limits_{x\to +\infty}(y-3x)=1$ nên $y=3x+1$ là tiệm cận xiên (khi $x\to +\infty$); $a=\lim\limits_{x\to -\infty}\dfrac{y}{x}=3$, $b=\lim\limits_{x\to -\infty}(y-3x)=1$ nên $y=3x+1$ là tiệm cận xiên (khi $x\to -\infty$); $\lim\limits_{x\to {-3}^{+}}y=+\infty$ nên $x=-3$ là tiệm cận đứng. a) Đúng. b) Sai. c) Sai: có $2$ đường tiệm cận. d) Sai: giới hạn của $y$ tại vô cực không hữu hạn.}]{Cho hàm số $y=3x+1+\dfrac{4}{x+3}$. Xét tính đúng sai của các mệnh đề sau:}{%
  \y{Đồ thị hàm số có tiệm cận đứng $x=-3$.}{d}
  \y{Đường thẳng $y=3x+3$ là tiệm cận xiên của đồ thị hàm số.}{s}
  \y{Đồ thị hàm số có tất cả $3$ đường tiệm cận.}{s}
  \y{Đồ thị hàm số có tiệm cận ngang.}{s}}

\cauDS[id=C1-B3-D03-NB-042,muc=NB,loigiai={Xét hàm số $y=3x+2+\dfrac{3}{x-1}$. Chia đa thức: $y=3x+2+\dfrac{3}{x - 1}$; $a=\lim\limits_{x\to +\infty}\dfrac{y}{x}=3$, $b=\lim\limits_{x\to +\infty}(y-3x)=2$ nên $y=3x+2$ là tiệm cận xiên (khi $x\to +\infty$); $a=\lim\limits_{x\to -\infty}\dfrac{y}{x}=3$, $b=\lim\limits_{x\to -\infty}(y-3x)=2$ nên $y=3x+2$ là tiệm cận xiên (khi $x\to -\infty$); $\lim\limits_{x\to {1}^{+}}y=+\infty$ nên $x=1$ là tiệm cận đứng. a) Đúng. b) Đúng. c) Đúng. d) Sai: tiệm cận xiên là $y=3x+2$.}]{Cho hàm số $y=3x+2+\dfrac{3}{x-1}$. Xét tính đúng sai của các mệnh đề sau:}{%
  \y{Đồ thị hàm số có tất cả $2$ đường tiệm cận.}{d}
  \y{Đồ thị hàm số có tiệm cận đứng $x=1$.}{d}
  \y{Giao điểm của tiệm cận đứng và tiệm cận xiên là $I(1;5)$.}{d}
  \y{Đường thẳng $y=3x$ là tiệm cận xiên của đồ thị hàm số.}{s}}

\cauDS[id=C1-B3-D03-NB-043,muc=NB,loigiai={Xét hàm số $y=x-4-\dfrac{1}{x}$. Chia đa thức: $y=x-4+\dfrac{-1}{x}$; $a=\lim\limits_{x\to +\infty}\dfrac{y}{x}=1$, $b=\lim\limits_{x\to +\infty}(y-x)=-4$ nên $y=x-4$ là tiệm cận xiên (khi $x\to +\infty$); $a=\lim\limits_{x\to -\infty}\dfrac{y}{x}=1$, $b=\lim\limits_{x\to -\infty}(y-x)=-4$ nên $y=x-4$ là tiệm cận xiên (khi $x\to -\infty$); $\lim\limits_{x\to {0}^{+}}y=-\infty$ nên $x=0$ là tiệm cận đứng. a) Đúng. b) Đúng. c) Sai. d) Sai: tiệm cận xiên là $y=x-4$.}]{Cho hàm số $y=x-4-\dfrac{1}{x}$. Xét tính đúng sai của các mệnh đề sau:}{%
  \y{Giao điểm của tiệm cận đứng và tiệm cận xiên là $I(0;-4)$.}{d}
  \y{Đồ thị hàm số có tất cả $2$ đường tiệm cận.}{d}
  \y{Đường thẳng $y=x-2$ là tiệm cận xiên của đồ thị hàm số.}{s}
  \y{Đường thẳng $y=x$ là tiệm cận xiên của đồ thị hàm số.}{s}}

\cauDS[id=C1-B3-D03-NB-044,muc=NB,loigiai={Xét hàm số $y=-2x+1-\dfrac{1}{x-3}$. Chia đa thức: $y=-2x+1+\dfrac{-1}{x - 3}$; $a=\lim\limits_{x\to +\infty}\dfrac{y}{x}=-2$, $b=\lim\limits_{x\to +\infty}(y+2x)=1$ nên $y=-2x+1$ là tiệm cận xiên (khi $x\to +\infty$); $a=\lim\limits_{x\to -\infty}\dfrac{y}{x}=-2$, $b=\lim\limits_{x\to -\infty}(y+2x)=1$ nên $y=-2x+1$ là tiệm cận xiên (khi $x\to -\infty$); $\lim\limits_{x\to {3}^{+}}y=-\infty$ nên $x=3$ là tiệm cận đứng. a) Đúng. b) Đúng. c) Sai: tiệm cận xiên là $y=-2x+1$. d) Sai: có $2$ đường tiệm cận.}]{Cho hàm số $y=-2x+1-\dfrac{1}{x-3}$. Xét tính đúng sai của các mệnh đề sau:}{%
  \y{Đồ thị hàm số có tiệm cận đứng $x=3$.}{d}
  \y{Giao điểm của tiệm cận đứng và tiệm cận xiên là $I(3;-5)$.}{d}
  \y{Đường thẳng $y=-2x$ là tiệm cận xiên của đồ thị hàm số.}{s}
  \y{Đồ thị hàm số có tất cả $3$ đường tiệm cận.}{s}}

\cauDS[id=C1-B3-D03-NB-045,muc=NB,loigiai={Xét hàm số $y=x-3-\dfrac{1}{x}$. Chia đa thức: $y=x-3+\dfrac{-1}{x}$; $a=\lim\limits_{x\to +\infty}\dfrac{y}{x}=1$, $b=\lim\limits_{x\to +\infty}(y-x)=-3$ nên $y=x-3$ là tiệm cận xiên (khi $x\to +\infty$); $a=\lim\limits_{x\to -\infty}\dfrac{y}{x}=1$, $b=\lim\limits_{x\to -\infty}(y-x)=-3$ nên $y=x-3$ là tiệm cận xiên (khi $x\to -\infty$); $\lim\limits_{x\to {0}^{+}}y=-\infty$ nên $x=0$ là tiệm cận đứng. a) Đúng. b) Sai: tiệm cận xiên là $y=x-3$. c) Sai. d) Đúng.}]{Cho hàm số $y=x-3-\dfrac{1}{x}$. Xét tính đúng sai của các mệnh đề sau:}{%
  \y{Đồ thị hàm số có tất cả $2$ đường tiệm cận.}{d}
  \y{Đường thẳng $y=x$ là tiệm cận xiên của đồ thị hàm số.}{s}
  \y{Đường thẳng $y=x-2$ là tiệm cận xiên của đồ thị hàm số.}{s}
  \y{Đồ thị hàm số có tiệm cận đứng $x=0$.}{d}}

\cauDS[id=C1-B3-D03-NB-046,muc=NB,loigiai={Xét hàm số $y=2x+1+\dfrac{3}{x-3}$. Chia đa thức: $y=2x+1+\dfrac{3}{x - 3}$; $a=\lim\limits_{x\to +\infty}\dfrac{y}{x}=2$, $b=\lim\limits_{x\to +\infty}(y-2x)=1$ nên $y=2x+1$ là tiệm cận xiên (khi $x\to +\infty$); $a=\lim\limits_{x\to -\infty}\dfrac{y}{x}=2$, $b=\lim\limits_{x\to -\infty}(y-2x)=1$ nên $y=2x+1$ là tiệm cận xiên (khi $x\to -\infty$); $\lim\limits_{x\to {3}^{+}}y=+\infty$ nên $x=3$ là tiệm cận đứng. a) Sai. b) Đúng. c) Đúng. d) Đúng.}]{Cho hàm số $y=2x+1+\dfrac{3}{x-3}$. Xét tính đúng sai của các mệnh đề sau:}{%
  \y{Đường thẳng $y=2x$ là tiệm cận xiên của đồ thị hàm số.}{s}
  \y{Đường thẳng $y=2x+1$ là tiệm cận xiên của đồ thị hàm số.}{d}
  \y{Giao điểm của tiệm cận đứng và tiệm cận xiên là $I(3;7)$.}{d}
  \y{Đồ thị hàm số có tiệm cận đứng $x=3$.}{d}}

\cauDS[id=C1-B3-D03-NB-047,muc=NB,loigiai={Xét hàm số $y=-x-4+\dfrac{2}{x-1}$. Chia đa thức: $y=-x-4+\dfrac{2}{x - 1}$; $a=\lim\limits_{x\to +\infty}\dfrac{y}{x}=-1$, $b=\lim\limits_{x\to +\infty}(y+x)=-4$ nên $y=-x-4$ là tiệm cận xiên (khi $x\to +\infty$); $a=\lim\limits_{x\to -\infty}\dfrac{y}{x}=-1$, $b=\lim\limits_{x\to -\infty}(y+x)=-4$ nên $y=-x-4$ là tiệm cận xiên (khi $x\to -\infty$); $\lim\limits_{x\to {1}^{+}}y=+\infty$ nên $x=1$ là tiệm cận đứng. a) Sai: tiệm cận xiên là $y=-x-4$. b) Đúng. c) Sai: giới hạn của $y$ tại vô cực không hữu hạn. d) Sai.}]{Cho hàm số $y=-x-4+\dfrac{2}{x-1}$. Xét tính đúng sai của các mệnh đề sau:}{%
  \y{Đường thẳng $y=-x$ là tiệm cận xiên của đồ thị hàm số.}{s}
  \y{Đồ thị hàm số có tất cả $2$ đường tiệm cận.}{d}
  \y{Đồ thị hàm số có tiệm cận ngang.}{s}
  \y{Đường thẳng $y=-x-3$ là tiệm cận xiên của đồ thị hàm số.}{s}}

\cauDS[id=C1-B3-D03-NB-048,muc=NB,loigiai={Xét hàm số $y=2x+\dfrac{1}{x+2}$. Chia đa thức: $y=2x+\dfrac{1}{x + 2}$; $a=\lim\limits_{x\to +\infty}\dfrac{y}{x}=2$, $b=\lim\limits_{x\to +\infty}(y-2x)=0$ nên $y=2x$ là tiệm cận xiên (khi $x\to +\infty$); $a=\lim\limits_{x\to -\infty}\dfrac{y}{x}=2$, $b=\lim\limits_{x\to -\infty}(y-2x)=0$ nên $y=2x$ là tiệm cận xiên (khi $x\to -\infty$); $\lim\limits_{x\to {-2}^{+}}y=+\infty$ nên $x=-2$ là tiệm cận đứng. a) Đúng. b) Sai: giới hạn của $y$ tại vô cực không hữu hạn. c) Đúng. d) Đúng.}]{Cho hàm số $y=2x+\dfrac{1}{x+2}$. Xét tính đúng sai của các mệnh đề sau:}{%
  \y{Đồ thị hàm số có tiệm cận đứng $x=-2$.}{d}
  \y{Đồ thị hàm số có tiệm cận ngang.}{s}
  \y{Giao điểm của tiệm cận đứng và tiệm cận xiên là $I(-2;-4)$.}{d}
  \y{Đường thẳng $y=2x$ là tiệm cận xiên của đồ thị hàm số.}{d}}

\cauTLN[id=C1-B3-D03-NB-049,muc=NB,loigiai={Xét hàm số $y=\dfrac{-2x^{2}-3x+5}{x+3}$. Chia đa thức: $y=-2x+3+\dfrac{-4}{x + 3}$; $a=\lim\limits_{x\to +\infty}\dfrac{y}{x}=-2$, $b=\lim\limits_{x\to +\infty}(y+2x)=3$ nên $y=-2x+3$ là tiệm cận xiên (khi $x\to +\infty$); $a=\lim\limits_{x\to -\infty}\dfrac{y}{x}=-2$, $b=\lim\limits_{x\to -\infty}(y+2x)=3$ nên $y=-2x+3$ là tiệm cận xiên (khi $x\to -\infty$); $\lim\limits_{x\to {-3}^{+}}y=-\infty$ nên $x=-3$ là tiệm cận đứng. Kết quả là $-7$.}]{Tiệm cận xiên của đồ thị hàm số $y=\dfrac{-2x^{2}-3x+5}{x+3}$ có phương trình $y=ax+b$. Tính $2a-b$.}{-7}

\cauTLN[id=C1-B3-D03-NB-050,muc=NB,loigiai={Xét hàm số $y=\dfrac{3x^{2}+5x+2}{-x-4}$. Chia đa thức: $y=-3x+7+\dfrac{30}{- x - 4}$; $a=\lim\limits_{x\to +\infty}\dfrac{y}{x}=-3$, $b=\lim\limits_{x\to +\infty}(y+3x)=7$ nên $y=-3x+7$ là tiệm cận xiên (khi $x\to +\infty$); $a=\lim\limits_{x\to -\infty}\dfrac{y}{x}=-3$, $b=\lim\limits_{x\to -\infty}(y+3x)=7$ nên $y=-3x+7$ là tiệm cận xiên (khi $x\to -\infty$); $\lim\limits_{x\to {-4}^{+}}y=-\infty$ nên $x=-4$ là tiệm cận đứng. Kết quả là $-13$.}]{Tiệm cận xiên của đồ thị hàm số $y=\dfrac{3x^{2}+5x+2}{-x-4}$ có phương trình $y=ax+b$. Tính $2a-b$.}{-13}

\cauTLN[id=C1-B3-D03-NB-051,muc=NB,loigiai={Xét hàm số $y=\dfrac{-x^{2}-4x-1}{-x+1}$. Chia đa thức: $y=x+5+\dfrac{-6}{1 - x}$; $a=\lim\limits_{x\to +\infty}\dfrac{y}{x}=1$, $b=\lim\limits_{x\to +\infty}(y-x)=5$ nên $y=x+5$ là tiệm cận xiên (khi $x\to +\infty$); $a=\lim\limits_{x\to -\infty}\dfrac{y}{x}=1$, $b=\lim\limits_{x\to -\infty}(y-x)=5$ nên $y=x+5$ là tiệm cận xiên (khi $x\to -\infty$); $\lim\limits_{x\to {1}^{+}}y=+\infty$ nên $x=1$ là tiệm cận đứng. Kết quả là $6$.}]{Tiệm cận xiên của đồ thị hàm số $y=\dfrac{-x^{2}-4x-1}{-x+1}$ có phương trình $y=ax+b$. Tính $a+b$.}{6}

\cauTLN[id=C1-B3-D03-NB-052,muc=NB,loigiai={Xét hàm số $y=\dfrac{-2x^{2}+2x-5}{x-1}$. Chia đa thức: $y=-2x+\dfrac{-5}{x - 1}$; $a=\lim\limits_{x\to +\infty}\dfrac{y}{x}=-2$, $b=\lim\limits_{x\to +\infty}(y+2x)=0$ nên $y=-2x$ là tiệm cận xiên (khi $x\to +\infty$); $a=\lim\limits_{x\to -\infty}\dfrac{y}{x}=-2$, $b=\lim\limits_{x\to -\infty}(y+2x)=0$ nên $y=-2x$ là tiệm cận xiên (khi $x\to -\infty$); $\lim\limits_{x\to {1}^{+}}y=-\infty$ nên $x=1$ là tiệm cận đứng. Kết quả là $-4$.}]{Tiệm cận xiên của đồ thị hàm số $y=\dfrac{-2x^{2}+2x-5}{x-1}$ có phương trình $y=ax+b$. Tính $2a-b$.}{-4}

\cauTLN[id=C1-B3-D03-NB-053,muc=NB,loigiai={Xét hàm số $y=3x-1+\dfrac{2}{x+3}$. Chia đa thức: $y=3x-1+\dfrac{2}{x + 3}$; $a=\lim\limits_{x\to +\infty}\dfrac{y}{x}=3$, $b=\lim\limits_{x\to +\infty}(y-3x)=-1$ nên $y=3x-1$ là tiệm cận xiên (khi $x\to +\infty$); $a=\lim\limits_{x\to -\infty}\dfrac{y}{x}=3$, $b=\lim\limits_{x\to -\infty}(y-3x)=-1$ nên $y=3x-1$ là tiệm cận xiên (khi $x\to -\infty$); $\lim\limits_{x\to {-3}^{+}}y=+\infty$ nên $x=-3$ là tiệm cận đứng. Kết quả là $-1$.}]{Tiệm cận xiên của đồ thị hàm số $y=3x-1+\dfrac{2}{x+3}$ cắt trục tung tại điểm có tung độ bằng bao nhiêu?}{-1}

\cauTLN[id=C1-B3-D03-NB-054,muc=NB,loigiai={Xét hàm số $y=-2x+4-\dfrac{2}{x-3}$. Chia đa thức: $y=-2x+4+\dfrac{-2}{x - 3}$; $a=\lim\limits_{x\to +\infty}\dfrac{y}{x}=-2$, $b=\lim\limits_{x\to +\infty}(y+2x)=4$ nên $y=-2x+4$ là tiệm cận xiên (khi $x\to +\infty$); $a=\lim\limits_{x\to -\infty}\dfrac{y}{x}=-2$, $b=\lim\limits_{x\to -\infty}(y+2x)=4$ nên $y=-2x+4$ là tiệm cận xiên (khi $x\to -\infty$); $\lim\limits_{x\to {3}^{+}}y=-\infty$ nên $x=3$ là tiệm cận đứng. Kết quả là $2$.}]{Tiệm cận xiên của đồ thị hàm số $y=-2x+4-\dfrac{2}{x-3}$ có phương trình $y=ax+b$. Tính $a+b$.}{2}

\cauTLN[id=C1-B3-D03-NB-055,muc=NB,loigiai={Xét hàm số $y=\dfrac{-2x^{2}-4x+3}{x-4}$. Chia đa thức: $y=-2x-12+\dfrac{-45}{x - 4}$; $a=\lim\limits_{x\to +\infty}\dfrac{y}{x}=-2$, $b=\lim\limits_{x\to +\infty}(y+2x)=-12$ nên $y=-2x-12$ là tiệm cận xiên (khi $x\to +\infty$); $a=\lim\limits_{x\to -\infty}\dfrac{y}{x}=-2$, $b=\lim\limits_{x\to -\infty}(y+2x)=-12$ nên $y=-2x-12$ là tiệm cận xiên (khi $x\to -\infty$); $\lim\limits_{x\to {4}^{+}}y=-\infty$ nên $x=4$ là tiệm cận đứng. Kết quả là $-14$.}]{Tiệm cận xiên của đồ thị hàm số $y=\dfrac{-2x^{2}-4x+3}{x-4}$ có phương trình $y=ax+b$. Tính $a+b$.}{-14}

\cauTLN[id=C1-B3-D03-NB-056,muc=NB,loigiai={Xét hàm số $y=\dfrac{2x^{2}+2}{2x-2}$. Chia đa thức: $y=x+1+\dfrac{2}{x - 1}$; $a=\lim\limits_{x\to +\infty}\dfrac{y}{x}=1$, $b=\lim\limits_{x\to +\infty}(y-x)=1$ nên $y=x+1$ là tiệm cận xiên (khi $x\to +\infty$); $a=\lim\limits_{x\to -\infty}\dfrac{y}{x}=1$, $b=\lim\limits_{x\to -\infty}(y-x)=1$ nên $y=x+1$ là tiệm cận xiên (khi $x\to -\infty$); $\lim\limits_{x\to {1}^{+}}y=+\infty$ nên $x=1$ là tiệm cận đứng. Kết quả là $1$.}]{Tiệm cận xiên của đồ thị hàm số $y=\dfrac{2x^{2}+2}{2x-2}$ cắt trục tung tại điểm có tung độ bằng bao nhiêu?}{1}

\cauTLN[id=C1-B3-D03-NB-057,muc=NB,loigiai={Xét hàm số $y=\dfrac{2x^{2}+4x-4}{x+3}$. Chia đa thức: $y=2x-2+\dfrac{2}{x + 3}$; $a=\lim\limits_{x\to +\infty}\dfrac{y}{x}=2$, $b=\lim\limits_{x\to +\infty}(y-2x)=-2$ nên $y=2x-2$ là tiệm cận xiên (khi $x\to +\infty$); $a=\lim\limits_{x\to -\infty}\dfrac{y}{x}=2$, $b=\lim\limits_{x\to -\infty}(y-2x)=-2$ nên $y=2x-2$ là tiệm cận xiên (khi $x\to -\infty$); $\lim\limits_{x\to {-3}^{+}}y=+\infty$ nên $x=-3$ là tiệm cận đứng. Kết quả là $6$.}]{Tiệm cận xiên của đồ thị hàm số $y=\dfrac{2x^{2}+4x-4}{x+3}$ có phương trình $y=ax+b$. Tính $2a-b$.}{6}

\cauTLN[id=C1-B3-D03-NB-058,muc=NB,loigiai={Xét hàm số $y=-x-3-\dfrac{1}{x}$. Chia đa thức: $y=-x-3+\dfrac{-1}{x}$; $a=\lim\limits_{x\to +\infty}\dfrac{y}{x}=-1$, $b=\lim\limits_{x\to +\infty}(y+x)=-3$ nên $y=-x-3$ là tiệm cận xiên (khi $x\to +\infty$); $a=\lim\limits_{x\to -\infty}\dfrac{y}{x}=-1$, $b=\lim\limits_{x\to -\infty}(y+x)=-3$ nên $y=-x-3$ là tiệm cận xiên (khi $x\to -\infty$); $\lim\limits_{x\to {0}^{+}}y=-\infty$ nên $x=0$ là tiệm cận đứng. Kết quả là $1$.}]{Tiệm cận xiên của đồ thị hàm số $y=-x-3-\dfrac{1}{x}$ có phương trình $y=ax+b$. Tính $2a-b$.}{1}

\cauTLN[id=C1-B3-D03-NB-059,muc=NB,loigiai={Xét hàm số $y=\dfrac{2x^{2}+2x+3}{x}$. Chia đa thức: $y=2x+2+\dfrac{3}{x}$; $a=\lim\limits_{x\to +\infty}\dfrac{y}{x}=2$, $b=\lim\limits_{x\to +\infty}(y-2x)=2$ nên $y=2x+2$ là tiệm cận xiên (khi $x\to +\infty$); $a=\lim\limits_{x\to -\infty}\dfrac{y}{x}=2$, $b=\lim\limits_{x\to -\infty}(y-2x)=2$ nên $y=2x+2$ là tiệm cận xiên (khi $x\to -\infty$); $\lim\limits_{x\to {0}^{+}}y=+\infty$ nên $x=0$ là tiệm cận đứng. Kết quả là $2$.}]{Tiệm cận xiên của đồ thị hàm số $y=\dfrac{2x^{2}+2x+3}{x}$ có phương trình $y=ax+b$. Tính $2a-b$.}{2}

\cauTLN[id=C1-B3-D03-NB-060,muc=NB,loigiai={Xét hàm số $y=x+\dfrac{3}{x-2}$. Chia đa thức: $y=x+\dfrac{3}{x - 2}$; $a=\lim\limits_{x\to +\infty}\dfrac{y}{x}=1$, $b=\lim\limits_{x\to +\infty}(y-x)=0$ nên $y=x$ là tiệm cận xiên (khi $x\to +\infty$); $a=\lim\limits_{x\to -\infty}\dfrac{y}{x}=1$, $b=\lim\limits_{x\to -\infty}(y-x)=0$ nên $y=x$ là tiệm cận xiên (khi $x\to -\infty$); $\lim\limits_{x\to {2}^{+}}y=+\infty$ nên $x=2$ là tiệm cận đứng. Kết quả là $0$.}]{Tiệm cận xiên của đồ thị hàm số $y=x+\dfrac{3}{x-2}$ cắt trục tung tại điểm có tung độ bằng bao nhiêu?}{0}

\cauTN[id=C1-B3-D03-TH-001,muc=TH,loigiai={Xét hàm số $y=\dfrac{-x^{3}-3x^{2}+3}{x^{2}+1}$. Chia đa thức: $y=-x-3+\dfrac{x + 6}{x^{2} + 1}$; $a=\lim\limits_{x\to +\infty}\dfrac{y}{x}=-1$, $b=\lim\limits_{x\to +\infty}(y+x)=-3$ nên $y=-x-3$ là tiệm cận xiên (khi $x\to +\infty$); $a=\lim\limits_{x\to -\infty}\dfrac{y}{x}=-1$, $b=\lim\limits_{x\to -\infty}(y+x)=-3$ nên $y=-x-3$ là tiệm cận xiên (khi $x\to -\infty$).}]{Tiệm cận xiên của đồ thị hàm số $y=\dfrac{-x^{3}-3x^{2}+3}{x^{2}+1}$ đi qua điểm nào dưới đây?}
  {$M(-2;-1)$.}
  {$M(-2;0)$.}{$M(-2;-2)$.}{$M(-2;1)$.}

\cauTN[id=C1-B3-D03-TH-002,muc=TH,loigiai={Xét hàm số $y=\dfrac{-2x^{2}-5x-1}{2x+2}$. Chia đa thức: $y=-x-\dfrac{3}{2}+\dfrac{2}{2 x + 2}$; $a=\lim\limits_{x\to +\infty}\dfrac{y}{x}=-1$, $b=\lim\limits_{x\to +\infty}(y+x)=-\dfrac{3}{2}$ nên $y=-x-\dfrac{3}{2}$ là tiệm cận xiên (khi $x\to +\infty$); $a=\lim\limits_{x\to -\infty}\dfrac{y}{x}=-1$, $b=\lim\limits_{x\to -\infty}(y+x)=-\dfrac{3}{2}$ nên $y=-x-\dfrac{3}{2}$ là tiệm cận xiên (khi $x\to -\infty$); $\lim\limits_{x\to {-1}^{+}}y=+\infty$ nên $x=-1$ là tiệm cận đứng.}]{Tổng số đường tiệm cận (đứng, ngang, xiên) của đồ thị hàm số $y=\dfrac{-2x^{2}-5x-1}{2x+2}$ là}
  {$2$.}
  {$3$.}{$1$.}{$4$.}

\cauTN[id=C1-B3-D03-TH-003,muc=TH,loigiai={Xét hàm số $y=\dfrac{x^{2}-5x-1}{2x-3}$. Chia đa thức: $y=\dfrac{1}{2}x-\dfrac{7}{4}+\dfrac{- \dfrac{25}{4}}{2 x - 3}$; $a=\lim\limits_{x\to +\infty}\dfrac{y}{x}=\dfrac{1}{2}$, $b=\lim\limits_{x\to +\infty}(y-\dfrac{1}{2}x)=-\dfrac{7}{4}$ nên $y=\dfrac{1}{2}x-\dfrac{7}{4}$ là tiệm cận xiên (khi $x\to +\infty$); $a=\lim\limits_{x\to -\infty}\dfrac{y}{x}=\dfrac{1}{2}$, $b=\lim\limits_{x\to -\infty}(y-\dfrac{1}{2}x)=-\dfrac{7}{4}$ nên $y=\dfrac{1}{2}x-\dfrac{7}{4}$ là tiệm cận xiên (khi $x\to -\infty$); $\lim\limits_{x\to {\dfrac{3}{2}}^{+}}y=-\infty$ nên $x=\dfrac{3}{2}$ là tiệm cận đứng.}]{Tiệm cận xiên của đồ thị hàm số $y=\dfrac{x^{2}-5x-1}{2x-3}$ đi qua điểm nào dưới đây?}
  {$M(-2;-\dfrac{11}{4})$.}
  {$M(-2;-\dfrac{7}{4})$.}{$M(-2;-\dfrac{15}{4})$.}{$M(-2;-\dfrac{3}{4})$.}

\cauTN[id=C1-B3-D03-TH-004,muc=TH,loigiai={Xét hàm số $y=\dfrac{1}{2}x-3+\dfrac{4}{x}$. Chia đa thức: $y=\dfrac{1}{2}x-3+\dfrac{8}{2 x}$; $a=\lim\limits_{x\to +\infty}\dfrac{y}{x}=\dfrac{1}{2}$, $b=\lim\limits_{x\to +\infty}(y-\dfrac{1}{2}x)=-3$ nên $y=\dfrac{1}{2}x-3$ là tiệm cận xiên (khi $x\to +\infty$); $a=\lim\limits_{x\to -\infty}\dfrac{y}{x}=\dfrac{1}{2}$, $b=\lim\limits_{x\to -\infty}(y-\dfrac{1}{2}x)=-3$ nên $y=\dfrac{1}{2}x-3$ là tiệm cận xiên (khi $x\to -\infty$); $\lim\limits_{x\to {0}^{+}}y=+\infty$ nên $x=0$ là tiệm cận đứng.}]{Tổng số đường tiệm cận (đứng, ngang, xiên) của đồ thị hàm số $y=\dfrac{1}{2}x-3+\dfrac{4}{x}$ là}
  {$2$.}
  {$3$.}{$1$.}{$4$.}

\cauTN[id=C1-B3-D03-TH-005,muc=TH,loigiai={Xét hàm số $y=-2x-1+\dfrac{3}{x-2}$. Chia đa thức: $y=-2x-1+\dfrac{3}{x - 2}$; $a=\lim\limits_{x\to +\infty}\dfrac{y}{x}=-2$, $b=\lim\limits_{x\to +\infty}(y+2x)=-1$ nên $y=-2x-1$ là tiệm cận xiên (khi $x\to +\infty$); $a=\lim\limits_{x\to -\infty}\dfrac{y}{x}=-2$, $b=\lim\limits_{x\to -\infty}(y+2x)=-1$ nên $y=-2x-1$ là tiệm cận xiên (khi $x\to -\infty$); $\lim\limits_{x\to {2}^{+}}y=+\infty$ nên $x=2$ là tiệm cận đứng.}]{Tổng số đường tiệm cận (đứng, ngang, xiên) của đồ thị hàm số $y=-2x-1+\dfrac{3}{x-2}$ là}
  {$2$.}
  {$3$.}{$1$.}{$4$.}

\cauTN[id=C1-B3-D03-TH-006,muc=TH,loigiai={Xét hàm số $y=\dfrac{-x^{2}+2x+6}{-x-4}$. Chia đa thức: $y=x-6+\dfrac{-18}{- x - 4}$; $a=\lim\limits_{x\to +\infty}\dfrac{y}{x}=1$, $b=\lim\limits_{x\to +\infty}(y-x)=-6$ nên $y=x-6$ là tiệm cận xiên (khi $x\to +\infty$); $a=\lim\limits_{x\to -\infty}\dfrac{y}{x}=1$, $b=\lim\limits_{x\to -\infty}(y-x)=-6$ nên $y=x-6$ là tiệm cận xiên (khi $x\to -\infty$); $\lim\limits_{x\to {-4}^{+}}y=+\infty$ nên $x=-4$ là tiệm cận đứng.}]{Tổng số đường tiệm cận (đứng, ngang, xiên) của đồ thị hàm số $y=\dfrac{-x^{2}+2x+6}{-x-4}$ là}
  {$2$.}
  {$3$.}{$1$.}{$4$.}

\cauTN[id=C1-B3-D03-TH-007,muc=TH,loigiai={Xét hàm số $y=-2x+2+\dfrac{1}{x}$. Chia đa thức: $y=-2x+2+\dfrac{1}{x}$; $a=\lim\limits_{x\to +\infty}\dfrac{y}{x}=-2$, $b=\lim\limits_{x\to +\infty}(y+2x)=2$ nên $y=-2x+2$ là tiệm cận xiên (khi $x\to +\infty$); $a=\lim\limits_{x\to -\infty}\dfrac{y}{x}=-2$, $b=\lim\limits_{x\to -\infty}(y+2x)=2$ nên $y=-2x+2$ là tiệm cận xiên (khi $x\to -\infty$); $\lim\limits_{x\to {0}^{+}}y=+\infty$ nên $x=0$ là tiệm cận đứng.}]{Tiệm cận xiên của đồ thị hàm số $y=-2x+2+\dfrac{1}{x}$ đi qua điểm nào dưới đây?}
  {$M(-2;6)$.}
  {$M(-2;7)$.}{$M(-2;5)$.}{$M(-2;8)$.}

\cauTN[id=C1-B3-D03-TH-008,muc=TH,loigiai={Xét hàm số $y=\dfrac{-2x^{2}-3x-1}{2x-2}$. Chia đa thức: $y=-x-\dfrac{5}{2}+\dfrac{-6}{2 x - 2}$; $a=\lim\limits_{x\to +\infty}\dfrac{y}{x}=-1$, $b=\lim\limits_{x\to +\infty}(y+x)=-\dfrac{5}{2}$ nên $y=-x-\dfrac{5}{2}$ là tiệm cận xiên (khi $x\to +\infty$); $a=\lim\limits_{x\to -\infty}\dfrac{y}{x}=-1$, $b=\lim\limits_{x\to -\infty}(y+x)=-\dfrac{5}{2}$ nên $y=-x-\dfrac{5}{2}$ là tiệm cận xiên (khi $x\to -\infty$); $\lim\limits_{x\to {1}^{+}}y=-\infty$ nên $x=1$ là tiệm cận đứng.}]{Tiệm cận xiên của đồ thị hàm số $y=\dfrac{-2x^{2}-3x-1}{2x-2}$ là đường thẳng}
  {$y=-x-\dfrac{5}{2}$.}
  {$y=-x+\dfrac{5}{2}$.}{$y=x-\dfrac{5}{2}$.}{$y=-x-\dfrac{3}{2}$.}

\cauTN[id=C1-B3-D03-TH-009,muc=TH,loigiai={Xét hàm số $y=\dfrac{-x^{2}-3x-3}{-x+2}$. Chia đa thức: $y=x+5+\dfrac{-13}{2 - x}$; $a=\lim\limits_{x\to +\infty}\dfrac{y}{x}=1$, $b=\lim\limits_{x\to +\infty}(y-x)=5$ nên $y=x+5$ là tiệm cận xiên (khi $x\to +\infty$); $a=\lim\limits_{x\to -\infty}\dfrac{y}{x}=1$, $b=\lim\limits_{x\to -\infty}(y-x)=5$ nên $y=x+5$ là tiệm cận xiên (khi $x\to -\infty$); $\lim\limits_{x\to {2}^{+}}y=+\infty$ nên $x=2$ là tiệm cận đứng.}]{Tiệm cận xiên của đồ thị hàm số $y=\dfrac{-x^{2}-3x-3}{-x+2}$ là đường thẳng}
  {$y=x+5$.}
  {$y=x-5$.}{$y=-x+5$.}{$y=x+6$.}

\cauTN[id=C1-B3-D03-TH-010,muc=TH,loigiai={Xét hàm số $y=-2x-2+\dfrac{2}{x}$. Chia đa thức: $y=-2x-2+\dfrac{2}{x}$; $a=\lim\limits_{x\to +\infty}\dfrac{y}{x}=-2$, $b=\lim\limits_{x\to +\infty}(y+2x)=-2$ nên $y=-2x-2$ là tiệm cận xiên (khi $x\to +\infty$); $a=\lim\limits_{x\to -\infty}\dfrac{y}{x}=-2$, $b=\lim\limits_{x\to -\infty}(y+2x)=-2$ nên $y=-2x-2$ là tiệm cận xiên (khi $x\to -\infty$); $\lim\limits_{x\to {0}^{+}}y=+\infty$ nên $x=0$ là tiệm cận đứng.}]{Tiệm cận xiên của đồ thị hàm số $y=-2x-2+\dfrac{2}{x}$ là đường thẳng}
  {$y=-2x-2$.}
  {$y=-2x+2$.}{$y=2x-2$.}{$y=-2x-1$.}

\cauTN[id=C1-B3-D03-TH-011,muc=TH,loigiai={Xét hàm số $y=\dfrac{-x^{2}-2x+3}{x-4}$. Chia đa thức: $y=-x-6+\dfrac{-21}{x - 4}$; $a=\lim\limits_{x\to +\infty}\dfrac{y}{x}=-1$, $b=\lim\limits_{x\to +\infty}(y+x)=-6$ nên $y=-x-6$ là tiệm cận xiên (khi $x\to +\infty$); $a=\lim\limits_{x\to -\infty}\dfrac{y}{x}=-1$, $b=\lim\limits_{x\to -\infty}(y+x)=-6$ nên $y=-x-6$ là tiệm cận xiên (khi $x\to -\infty$); $\lim\limits_{x\to {4}^{+}}y=-\infty$ nên $x=4$ là tiệm cận đứng.}]{Tiệm cận xiên của đồ thị hàm số $y=\dfrac{-x^{2}-2x+3}{x-4}$ là đường thẳng}
  {$y=-x-6$.}
  {$y=-x+6$.}{$y=x-6$.}{$y=-x-5$.}

\cauTN[id=C1-B3-D03-TH-012,muc=TH,loigiai={Xét hàm số $y=\dfrac{-x^{3}+3x^{2}-2}{x^{2}+1}$. Chia đa thức: $y=-x+3+\dfrac{x - 5}{x^{2} + 1}$; $a=\lim\limits_{x\to +\infty}\dfrac{y}{x}=-1$, $b=\lim\limits_{x\to +\infty}(y+x)=3$ nên $y=-x+3$ là tiệm cận xiên (khi $x\to +\infty$); $a=\lim\limits_{x\to -\infty}\dfrac{y}{x}=-1$, $b=\lim\limits_{x\to -\infty}(y+x)=3$ nên $y=-x+3$ là tiệm cận xiên (khi $x\to -\infty$).}]{Tiệm cận xiên của đồ thị hàm số $y=\dfrac{-x^{3}+3x^{2}-2}{x^{2}+1}$ là đường thẳng}
  {$y=-x+3$.}
  {$y=-x-3$.}{$y=x+3$.}{$y=-x+4$.}

\cauTN[id=C1-B3-D03-TH-013,muc=TH,loigiai={Xét hàm số $y=\dfrac{-x^{2}+x+3}{x+3}$. Chia đa thức: $y=-x+4+\dfrac{-9}{x + 3}$; $a=\lim\limits_{x\to +\infty}\dfrac{y}{x}=-1$, $b=\lim\limits_{x\to +\infty}(y+x)=4$ nên $y=-x+4$ là tiệm cận xiên (khi $x\to +\infty$); $a=\lim\limits_{x\to -\infty}\dfrac{y}{x}=-1$, $b=\lim\limits_{x\to -\infty}(y+x)=4$ nên $y=-x+4$ là tiệm cận xiên (khi $x\to -\infty$); $\lim\limits_{x\to {-3}^{+}}y=-\infty$ nên $x=-3$ là tiệm cận đứng.}]{Tiệm cận xiên của đồ thị hàm số $y=\dfrac{-x^{2}+x+3}{x+3}$ đi qua điểm nào dưới đây?}
  {$M(-1;5)$.}
  {$M(-1;6)$.}{$M(-1;4)$.}{$M(-1;7)$.}

\cauTN[id=C1-B3-D03-TH-014,muc=TH,loigiai={Xét hàm số $y=\dfrac{1}{2}x+1+\dfrac{4}{x}$. Chia đa thức: $y=\dfrac{1}{2}x+1+\dfrac{8}{2 x}$; $a=\lim\limits_{x\to +\infty}\dfrac{y}{x}=\dfrac{1}{2}$, $b=\lim\limits_{x\to +\infty}(y-\dfrac{1}{2}x)=1$ nên $y=\dfrac{1}{2}x+1$ là tiệm cận xiên (khi $x\to +\infty$); $a=\lim\limits_{x\to -\infty}\dfrac{y}{x}=\dfrac{1}{2}$, $b=\lim\limits_{x\to -\infty}(y-\dfrac{1}{2}x)=1$ nên $y=\dfrac{1}{2}x+1$ là tiệm cận xiên (khi $x\to -\infty$); $\lim\limits_{x\to {0}^{+}}y=+\infty$ nên $x=0$ là tiệm cận đứng.}]{Tiệm cận xiên của đồ thị hàm số $y=\dfrac{1}{2}x+1+\dfrac{4}{x}$ là đường thẳng}
  {$y=\dfrac{1}{2}x+1$.}
  {$y=\dfrac{1}{2}x-1$.}{$y=-\dfrac{1}{2}x+1$.}{$y=\dfrac{1}{2}x+2$.}

\cauTN[id=C1-B3-D03-TH-015,muc=TH,loigiai={Xét hàm số $y=\dfrac{3x^{2}+x+5}{-x}$. Chia đa thức: $y=-3x-1+\dfrac{-5}{x}$; $a=\lim\limits_{x\to +\infty}\dfrac{y}{x}=-3$, $b=\lim\limits_{x\to +\infty}(y+3x)=-1$ nên $y=-3x-1$ là tiệm cận xiên (khi $x\to +\infty$); $a=\lim\limits_{x\to -\infty}\dfrac{y}{x}=-3$, $b=\lim\limits_{x\to -\infty}(y+3x)=-1$ nên $y=-3x-1$ là tiệm cận xiên (khi $x\to -\infty$); $\lim\limits_{x\to {0}^{+}}y=-\infty$ nên $x=0$ là tiệm cận đứng.}]{Tiệm cận xiên của đồ thị hàm số $y=\dfrac{3x^{2}+x+5}{-x}$ là đường thẳng}
  {$y=-3x-1$.}
  {$y=-3x+1$.}{$y=3x-1$.}{$y=-3x$.}

\cauTN[id=C1-B3-D03-TH-016,muc=TH,loigiai={Xét hàm số $y=\dfrac{-x^{2}-3x+3}{x-1}$. Chia đa thức: $y=-x-4+\dfrac{-1}{x - 1}$; $a=\lim\limits_{x\to +\infty}\dfrac{y}{x}=-1$, $b=\lim\limits_{x\to +\infty}(y+x)=-4$ nên $y=-x-4$ là tiệm cận xiên (khi $x\to +\infty$); $a=\lim\limits_{x\to -\infty}\dfrac{y}{x}=-1$, $b=\lim\limits_{x\to -\infty}(y+x)=-4$ nên $y=-x-4$ là tiệm cận xiên (khi $x\to -\infty$); $\lim\limits_{x\to {1}^{+}}y=-\infty$ nên $x=1$ là tiệm cận đứng.}]{Tiệm cận xiên của đồ thị hàm số $y=\dfrac{-x^{2}-3x+3}{x-1}$ là đường thẳng}
  {$y=-x-4$.}
  {$y=-x+4$.}{$y=x-4$.}{$y=-x-3$.}

\cauTN[id=C1-B3-D03-TH-017,muc=TH,loigiai={Xét hàm số $y=\dfrac{-2x^{2}-6}{-x-3}$. Chia đa thức: $y=2x-6+\dfrac{-24}{- x - 3}$; $a=\lim\limits_{x\to +\infty}\dfrac{y}{x}=2$, $b=\lim\limits_{x\to +\infty}(y-2x)=-6$ nên $y=2x-6$ là tiệm cận xiên (khi $x\to +\infty$); $a=\lim\limits_{x\to -\infty}\dfrac{y}{x}=2$, $b=\lim\limits_{x\to -\infty}(y-2x)=-6$ nên $y=2x-6$ là tiệm cận xiên (khi $x\to -\infty$); $\lim\limits_{x\to {-3}^{+}}y=+\infty$ nên $x=-3$ là tiệm cận đứng.}]{Tổng số đường tiệm cận (đứng, ngang, xiên) của đồ thị hàm số $y=\dfrac{-2x^{2}-6}{-x-3}$ là}
  {$2$.}
  {$3$.}{$1$.}{$4$.}

\cauTN[id=C1-B3-D03-TH-018,muc=TH,loigiai={Xét hàm số $y=\dfrac{x^{3}-x^{2}-1}{x^{2}-9}$. Chia đa thức: $y=x-1+\dfrac{9 x - 10}{x^{2} - 9}$; $a=\lim\limits_{x\to +\infty}\dfrac{y}{x}=1$, $b=\lim\limits_{x\to +\infty}(y-x)=-1$ nên $y=x-1$ là tiệm cận xiên (khi $x\to +\infty$); $a=\lim\limits_{x\to -\infty}\dfrac{y}{x}=1$, $b=\lim\limits_{x\to -\infty}(y-x)=-1$ nên $y=x-1$ là tiệm cận xiên (khi $x\to -\infty$); $\lim\limits_{x\to {-3}^{+}}y=+\infty$ nên $x=-3$ là tiệm cận đứng; $\lim\limits_{x\to {3}^{+}}y=+\infty$ nên $x=3$ là tiệm cận đứng.}]{Tiệm cận xiên của đồ thị hàm số $y=\dfrac{x^{3}-x^{2}-1}{x^{2}-9}$ là đường thẳng}
  {$y=x-1$.}
  {$y=x+1$.}{$y=-x-1$.}{$y=x$.}

\cauTN[id=C1-B3-D03-TH-019,muc=TH,loigiai={Xét hàm số $y=\dfrac{2x^{2}-3x+5}{-x+1}$. Chia đa thức: $y=-2x+1+\dfrac{4}{1 - x}$; $a=\lim\limits_{x\to +\infty}\dfrac{y}{x}=-2$, $b=\lim\limits_{x\to +\infty}(y+2x)=1$ nên $y=-2x+1$ là tiệm cận xiên (khi $x\to +\infty$); $a=\lim\limits_{x\to -\infty}\dfrac{y}{x}=-2$, $b=\lim\limits_{x\to -\infty}(y+2x)=1$ nên $y=-2x+1$ là tiệm cận xiên (khi $x\to -\infty$); $\lim\limits_{x\to {1}^{+}}y=-\infty$ nên $x=1$ là tiệm cận đứng.}]{Tiệm cận xiên của đồ thị hàm số $y=\dfrac{2x^{2}-3x+5}{-x+1}$ là đường thẳng}
  {$y=-2x+1$.}
  {$y=-2x-1$.}{$y=2x+1$.}{$y=-2x+2$.}

\cauTN[id=C1-B3-D03-TH-020,muc=TH,loigiai={Xét hàm số $y=\dfrac{2x^{2}-5x-1}{2x-2}$. Chia đa thức: $y=x-\dfrac{3}{2}+\dfrac{-4}{2 x - 2}$; $a=\lim\limits_{x\to +\infty}\dfrac{y}{x}=1$, $b=\lim\limits_{x\to +\infty}(y-x)=-\dfrac{3}{2}$ nên $y=x-\dfrac{3}{2}$ là tiệm cận xiên (khi $x\to +\infty$); $a=\lim\limits_{x\to -\infty}\dfrac{y}{x}=1$, $b=\lim\limits_{x\to -\infty}(y-x)=-\dfrac{3}{2}$ nên $y=x-\dfrac{3}{2}$ là tiệm cận xiên (khi $x\to -\infty$); $\lim\limits_{x\to {1}^{+}}y=-\infty$ nên $x=1$ là tiệm cận đứng.}]{Tiệm cận xiên của đồ thị hàm số $y=\dfrac{2x^{2}-5x-1}{2x-2}$ đi qua điểm nào dưới đây?}
  {$M(-1;-\dfrac{5}{2})$.}
  {$M(-1;-\dfrac{3}{2})$.}{$M(-1;-\dfrac{7}{2})$.}{$M(-1;-\dfrac{1}{2})$.}

\cauTN[id=C1-B3-D03-TH-021,muc=TH,loigiai={Xét hàm số $y=\dfrac{-2x^{2}+5x+1}{x-1}$. Chia đa thức: $y=-2x+3+\dfrac{4}{x - 1}$; $a=\lim\limits_{x\to +\infty}\dfrac{y}{x}=-2$, $b=\lim\limits_{x\to +\infty}(y+2x)=3$ nên $y=-2x+3$ là tiệm cận xiên (khi $x\to +\infty$); $a=\lim\limits_{x\to -\infty}\dfrac{y}{x}=-2$, $b=\lim\limits_{x\to -\infty}(y+2x)=3$ nên $y=-2x+3$ là tiệm cận xiên (khi $x\to -\infty$); $\lim\limits_{x\to {1}^{+}}y=+\infty$ nên $x=1$ là tiệm cận đứng.}]{Tiệm cận xiên của đồ thị hàm số $y=\dfrac{-2x^{2}+5x+1}{x-1}$ là đường thẳng}
  {$y=-2x+3$.}
  {$y=-2x-3$.}{$y=2x+3$.}{$y=-2x+4$.}

\cauTN[id=C1-B3-D03-TH-022,muc=TH,loigiai={Xét hàm số $y=\dfrac{3x^{2}+4x-5}{-x-3}$. Chia đa thức: $y=-3x+5+\dfrac{10}{- x - 3}$; $a=\lim\limits_{x\to +\infty}\dfrac{y}{x}=-3$, $b=\lim\limits_{x\to +\infty}(y+3x)=5$ nên $y=-3x+5$ là tiệm cận xiên (khi $x\to +\infty$); $a=\lim\limits_{x\to -\infty}\dfrac{y}{x}=-3$, $b=\lim\limits_{x\to -\infty}(y+3x)=5$ nên $y=-3x+5$ là tiệm cận xiên (khi $x\to -\infty$); $\lim\limits_{x\to {-3}^{+}}y=-\infty$ nên $x=-3$ là tiệm cận đứng.}]{Tiệm cận xiên của đồ thị hàm số $y=\dfrac{3x^{2}+4x-5}{-x-3}$ là đường thẳng}
  {$y=-3x+5$.}
  {$y=-3x-5$.}{$y=3x+5$.}{$y=-3x+6$.}

\cauTN[id=C1-B3-D03-TH-023,muc=TH,loigiai={Xét hàm số $y=-x+1+\dfrac{3}{x-1}$. Chia đa thức: $y=-x+1+\dfrac{3}{x - 1}$; $a=\lim\limits_{x\to +\infty}\dfrac{y}{x}=-1$, $b=\lim\limits_{x\to +\infty}(y+x)=1$ nên $y=-x+1$ là tiệm cận xiên (khi $x\to +\infty$); $a=\lim\limits_{x\to -\infty}\dfrac{y}{x}=-1$, $b=\lim\limits_{x\to -\infty}(y+x)=1$ nên $y=-x+1$ là tiệm cận xiên (khi $x\to -\infty$); $\lim\limits_{x\to {1}^{+}}y=+\infty$ nên $x=1$ là tiệm cận đứng.}]{Tổng số đường tiệm cận (đứng, ngang, xiên) của đồ thị hàm số $y=-x+1+\dfrac{3}{x-1}$ là}
  {$2$.}
  {$3$.}{$1$.}{$4$.}

\cauTN[id=C1-B3-D03-TH-024,muc=TH,loigiai={Xét hàm số $y=\dfrac{-x^{3}+x^{2}-1}{x^{2}+2}$. Chia đa thức: $y=-x+1+\dfrac{2 x - 3}{x^{2} + 2}$; $a=\lim\limits_{x\to +\infty}\dfrac{y}{x}=-1$, $b=\lim\limits_{x\to +\infty}(y+x)=1$ nên $y=-x+1$ là tiệm cận xiên (khi $x\to +\infty$); $a=\lim\limits_{x\to -\infty}\dfrac{y}{x}=-1$, $b=\lim\limits_{x\to -\infty}(y+x)=1$ nên $y=-x+1$ là tiệm cận xiên (khi $x\to -\infty$).}]{Tổng số đường tiệm cận (đứng, ngang, xiên) của đồ thị hàm số $y=\dfrac{-x^{3}+x^{2}-1}{x^{2}+2}$ là}
  {$1$.}
  {$2$.}{$0$.}{$3$.}

\cauTN[id=C1-B3-D03-TH-025,muc=TH,loigiai={Xét hàm số $y=\dfrac{2x^{3}+x^{2}+3}{x^{2}-4}$. Chia đa thức: $y=2x+1+\dfrac{8 x + 7}{x^{2} - 4}$; $a=\lim\limits_{x\to +\infty}\dfrac{y}{x}=2$, $b=\lim\limits_{x\to +\infty}(y-2x)=1$ nên $y=2x+1$ là tiệm cận xiên (khi $x\to +\infty$); $a=\lim\limits_{x\to -\infty}\dfrac{y}{x}=2$, $b=\lim\limits_{x\to -\infty}(y-2x)=1$ nên $y=2x+1$ là tiệm cận xiên (khi $x\to -\infty$); $\lim\limits_{x\to {-2}^{+}}y=+\infty$ nên $x=-2$ là tiệm cận đứng; $\lim\limits_{x\to {2}^{+}}y=+\infty$ nên $x=2$ là tiệm cận đứng.}]{Tiệm cận xiên của đồ thị hàm số $y=\dfrac{2x^{3}+x^{2}+3}{x^{2}-4}$ đi qua điểm nào dưới đây?}
  {$M(-2;-3)$.}
  {$M(-2;-2)$.}{$M(-2;-4)$.}{$M(-2;-1)$.}

\cauTN[id=C1-B3-D03-TH-026,muc=TH,loigiai={Xét hàm số $y=\dfrac{2x^{3}}{x^{2}+1}$. Chia đa thức: $y=2x+\dfrac{- 2 x}{x^{2} + 1}$; $a=\lim\limits_{x\to +\infty}\dfrac{y}{x}=2$, $b=\lim\limits_{x\to +\infty}(y-2x)=0$ nên $y=2x$ là tiệm cận xiên (khi $x\to +\infty$); $a=\lim\limits_{x\to -\infty}\dfrac{y}{x}=2$, $b=\lim\limits_{x\to -\infty}(y-2x)=0$ nên $y=2x$ là tiệm cận xiên (khi $x\to -\infty$).}]{Tổng số đường tiệm cận (đứng, ngang, xiên) của đồ thị hàm số $y=\dfrac{2x^{3}}{x^{2}+1}$ là}
  {$1$.}
  {$2$.}{$0$.}{$3$.}

\cauTN[id=C1-B3-D03-TH-027,muc=TH,loigiai={Xét hàm số $y=\dfrac{2x^{3}+2x^{2}+1}{x^{2}+4}$. Chia đa thức: $y=2x+2+\dfrac{- 8 x - 7}{x^{2} + 4}$; $a=\lim\limits_{x\to +\infty}\dfrac{y}{x}=2$, $b=\lim\limits_{x\to +\infty}(y-2x)=2$ nên $y=2x+2$ là tiệm cận xiên (khi $x\to +\infty$); $a=\lim\limits_{x\to -\infty}\dfrac{y}{x}=2$, $b=\lim\limits_{x\to -\infty}(y-2x)=2$ nên $y=2x+2$ là tiệm cận xiên (khi $x\to -\infty$).}]{Tiệm cận xiên của đồ thị hàm số $y=\dfrac{2x^{3}+2x^{2}+1}{x^{2}+4}$ là đường thẳng}
  {$y=2x+2$.}
  {$y=2x-2$.}{$y=-2x+2$.}{$y=2x+3$.}

\cauTN[id=C1-B3-D03-TH-028,muc=TH,loigiai={Xét hàm số $y=\dfrac{-x^{2}-5x-4}{-x+2}$. Chia đa thức: $y=x+7+\dfrac{-18}{2 - x}$; $a=\lim\limits_{x\to +\infty}\dfrac{y}{x}=1$, $b=\lim\limits_{x\to +\infty}(y-x)=7$ nên $y=x+7$ là tiệm cận xiên (khi $x\to +\infty$); $a=\lim\limits_{x\to -\infty}\dfrac{y}{x}=1$, $b=\lim\limits_{x\to -\infty}(y-x)=7$ nên $y=x+7$ là tiệm cận xiên (khi $x\to -\infty$); $\lim\limits_{x\to {2}^{+}}y=+\infty$ nên $x=2$ là tiệm cận đứng.}]{Tiệm cận xiên của đồ thị hàm số $y=\dfrac{-x^{2}-5x-4}{-x+2}$ là đường thẳng}
  {$y=x+7$.}
  {$y=x-7$.}{$y=-x+7$.}{$y=x+8$.}

\cauTN[id=C1-B3-D03-TH-029,muc=TH,loigiai={Xét hàm số $y=3x+\dfrac{1}{x-2}$. Chia đa thức: $y=3x+\dfrac{1}{x - 2}$; $a=\lim\limits_{x\to +\infty}\dfrac{y}{x}=3$, $b=\lim\limits_{x\to +\infty}(y-3x)=0$ nên $y=3x$ là tiệm cận xiên (khi $x\to +\infty$); $a=\lim\limits_{x\to -\infty}\dfrac{y}{x}=3$, $b=\lim\limits_{x\to -\infty}(y-3x)=0$ nên $y=3x$ là tiệm cận xiên (khi $x\to -\infty$); $\lim\limits_{x\to {2}^{+}}y=+\infty$ nên $x=2$ là tiệm cận đứng.}]{Tiệm cận xiên của đồ thị hàm số $y=3x+\dfrac{1}{x-2}$ là đường thẳng}
  {$y=3x$.}
  {$y=-3x$.}{$y=3x+1$.}{$y=3x-1$.}

\cauTN[id=C1-B3-D03-TH-030,muc=TH,loigiai={Xét hàm số $y=\dfrac{3x^{2}-5x+4}{-x-4}$. Chia đa thức: $y=-3x+17+\dfrac{72}{- x - 4}$; $a=\lim\limits_{x\to +\infty}\dfrac{y}{x}=-3$, $b=\lim\limits_{x\to +\infty}(y+3x)=17$ nên $y=-3x+17$ là tiệm cận xiên (khi $x\to +\infty$); $a=\lim\limits_{x\to -\infty}\dfrac{y}{x}=-3$, $b=\lim\limits_{x\to -\infty}(y+3x)=17$ nên $y=-3x+17$ là tiệm cận xiên (khi $x\to -\infty$); $\lim\limits_{x\to {-4}^{+}}y=-\infty$ nên $x=-4$ là tiệm cận đứng.}]{Tiệm cận xiên của đồ thị hàm số $y=\dfrac{3x^{2}-5x+4}{-x-4}$ đi qua điểm nào dưới đây?}
  {$M(-1;20)$.}
  {$M(-1;21)$.}{$M(-1;19)$.}{$M(-1;22)$.}

\cauTN[id=C1-B3-D03-TH-031,muc=TH,loigiai={Xét hàm số $y=\dfrac{-x^{3}+3x^{2}+3}{x^{2}-1}$. Chia đa thức: $y=-x+3+\dfrac{6 - x}{x^{2} - 1}$; $a=\lim\limits_{x\to +\infty}\dfrac{y}{x}=-1$, $b=\lim\limits_{x\to +\infty}(y+x)=3$ nên $y=-x+3$ là tiệm cận xiên (khi $x\to +\infty$); $a=\lim\limits_{x\to -\infty}\dfrac{y}{x}=-1$, $b=\lim\limits_{x\to -\infty}(y+x)=3$ nên $y=-x+3$ là tiệm cận xiên (khi $x\to -\infty$); $\lim\limits_{x\to {-1}^{+}}y=-\infty$ nên $x=-1$ là tiệm cận đứng; $\lim\limits_{x\to {1}^{+}}y=+\infty$ nên $x=1$ là tiệm cận đứng.}]{Tiệm cận xiên của đồ thị hàm số $y=\dfrac{-x^{3}+3x^{2}+3}{x^{2}-1}$ là đường thẳng}
  {$y=-x+3$.}
  {$y=-x-3$.}{$y=x+3$.}{$y=-x+4$.}

\cauTN[id=C1-B3-D03-TH-032,muc=TH,loigiai={Xét hàm số $y=\dfrac{2x^{2}-2x+4}{x+4}$. Chia đa thức: $y=2x-10+\dfrac{44}{x + 4}$; $a=\lim\limits_{x\to +\infty}\dfrac{y}{x}=2$, $b=\lim\limits_{x\to +\infty}(y-2x)=-10$ nên $y=2x-10$ là tiệm cận xiên (khi $x\to +\infty$); $a=\lim\limits_{x\to -\infty}\dfrac{y}{x}=2$, $b=\lim\limits_{x\to -\infty}(y-2x)=-10$ nên $y=2x-10$ là tiệm cận xiên (khi $x\to -\infty$); $\lim\limits_{x\to {-4}^{+}}y=+\infty$ nên $x=-4$ là tiệm cận đứng.}]{Tiệm cận xiên của đồ thị hàm số $y=\dfrac{2x^{2}-2x+4}{x+4}$ đi qua điểm nào dưới đây?}
  {$M(1;-8)$.}
  {$M(1;-7)$.}{$M(1;-9)$.}{$M(1;-6)$.}

\cauTN[id=C1-B3-D03-TH-033,muc=TH,loigiai={Xét hàm số $y=2x+1+\dfrac{4}{x-1}$. Chia đa thức: $y=2x+1+\dfrac{4}{x - 1}$; $a=\lim\limits_{x\to +\infty}\dfrac{y}{x}=2$, $b=\lim\limits_{x\to +\infty}(y-2x)=1$ nên $y=2x+1$ là tiệm cận xiên (khi $x\to +\infty$); $a=\lim\limits_{x\to -\infty}\dfrac{y}{x}=2$, $b=\lim\limits_{x\to -\infty}(y-2x)=1$ nên $y=2x+1$ là tiệm cận xiên (khi $x\to -\infty$); $\lim\limits_{x\to {1}^{+}}y=+\infty$ nên $x=1$ là tiệm cận đứng.}]{Tiệm cận xiên của đồ thị hàm số $y=2x+1+\dfrac{4}{x-1}$ là đường thẳng}
  {$y=2x+1$.}
  {$y=2x-1$.}{$y=-2x+1$.}{$y=2x+2$.}

\cauTN[id=C1-B3-D03-TH-034,muc=TH,loigiai={Xét hàm số $y=\dfrac{-x^{3}+x^{2}+1}{x^{2}-9}$. Chia đa thức: $y=-x+1+\dfrac{10 - 9 x}{x^{2} - 9}$; $a=\lim\limits_{x\to +\infty}\dfrac{y}{x}=-1$, $b=\lim\limits_{x\to +\infty}(y+x)=1$ nên $y=-x+1$ là tiệm cận xiên (khi $x\to +\infty$); $a=\lim\limits_{x\to -\infty}\dfrac{y}{x}=-1$, $b=\lim\limits_{x\to -\infty}(y+x)=1$ nên $y=-x+1$ là tiệm cận xiên (khi $x\to -\infty$); $\lim\limits_{x\to {-3}^{+}}y=-\infty$ nên $x=-3$ là tiệm cận đứng; $\lim\limits_{x\to {3}^{+}}y=-\infty$ nên $x=3$ là tiệm cận đứng.}]{Tổng số đường tiệm cận (đứng, ngang, xiên) của đồ thị hàm số $y=\dfrac{-x^{3}+x^{2}+1}{x^{2}-9}$ là}
  {$3$.}
  {$4$.}{$2$.}{$5$.}

\cauTN[id=C1-B3-D03-TH-035,muc=TH,loigiai={Xét hàm số $y=\dfrac{2x^{2}+2x-4}{x+1}$. Chia đa thức: $y=2x+\dfrac{-4}{x + 1}$; $a=\lim\limits_{x\to +\infty}\dfrac{y}{x}=2$, $b=\lim\limits_{x\to +\infty}(y-2x)=0$ nên $y=2x$ là tiệm cận xiên (khi $x\to +\infty$); $a=\lim\limits_{x\to -\infty}\dfrac{y}{x}=2$, $b=\lim\limits_{x\to -\infty}(y-2x)=0$ nên $y=2x$ là tiệm cận xiên (khi $x\to -\infty$); $\lim\limits_{x\to {-1}^{+}}y=-\infty$ nên $x=-1$ là tiệm cận đứng.}]{Tiệm cận xiên của đồ thị hàm số $y=\dfrac{2x^{2}+2x-4}{x+1}$ là đường thẳng}
  {$y=2x$.}
  {$y=-2x$.}{$y=2x+1$.}{$y=2x-1$.}

\cauTN[id=C1-B3-D03-TH-036,muc=TH,loigiai={Xét hàm số $y=\dfrac{x^{3}-2x^{2}-2}{x^{2}+2}$. Chia đa thức: $y=x-2+\dfrac{2 - 2 x}{x^{2} + 2}$; $a=\lim\limits_{x\to +\infty}\dfrac{y}{x}=1$, $b=\lim\limits_{x\to +\infty}(y-x)=-2$ nên $y=x-2$ là tiệm cận xiên (khi $x\to +\infty$); $a=\lim\limits_{x\to -\infty}\dfrac{y}{x}=1$, $b=\lim\limits_{x\to -\infty}(y-x)=-2$ nên $y=x-2$ là tiệm cận xiên (khi $x\to -\infty$).}]{Tiệm cận xiên của đồ thị hàm số $y=\dfrac{x^{3}-2x^{2}-2}{x^{2}+2}$ đi qua điểm nào dưới đây?}
  {$M(3;1)$.}
  {$M(3;2)$.}{$M(3;0)$.}{$M(3;3)$.}

\cauTN[id=C1-B3-D03-TH-037,muc=TH,loigiai={Xét hàm số $y=3x-3+\dfrac{3}{x-2}$. Chia đa thức: $y=3x-3+\dfrac{3}{x - 2}$; $a=\lim\limits_{x\to +\infty}\dfrac{y}{x}=3$, $b=\lim\limits_{x\to +\infty}(y-3x)=-3$ nên $y=3x-3$ là tiệm cận xiên (khi $x\to +\infty$); $a=\lim\limits_{x\to -\infty}\dfrac{y}{x}=3$, $b=\lim\limits_{x\to -\infty}(y-3x)=-3$ nên $y=3x-3$ là tiệm cận xiên (khi $x\to -\infty$); $\lim\limits_{x\to {2}^{+}}y=+\infty$ nên $x=2$ là tiệm cận đứng.}]{Tiệm cận xiên của đồ thị hàm số $y=3x-3+\dfrac{3}{x-2}$ là đường thẳng}
  {$y=3x-3$.}
  {$y=3x+3$.}{$y=-3x-3$.}{$y=3x-2$.}

\cauTN[id=C1-B3-D03-TH-038,muc=TH,loigiai={Xét hàm số $y=\dfrac{-x^{2}-3x-1}{-x-2}$. Chia đa thức: $y=x+1+\dfrac{1}{- x - 2}$; $a=\lim\limits_{x\to +\infty}\dfrac{y}{x}=1$, $b=\lim\limits_{x\to +\infty}(y-x)=1$ nên $y=x+1$ là tiệm cận xiên (khi $x\to +\infty$); $a=\lim\limits_{x\to -\infty}\dfrac{y}{x}=1$, $b=\lim\limits_{x\to -\infty}(y-x)=1$ nên $y=x+1$ là tiệm cận xiên (khi $x\to -\infty$); $\lim\limits_{x\to {-2}^{+}}y=-\infty$ nên $x=-2$ là tiệm cận đứng.}]{Tiệm cận xiên của đồ thị hàm số $y=\dfrac{-x^{2}-3x-1}{-x-2}$ là đường thẳng}
  {$y=x+1$.}
  {$y=x-1$.}{$y=-x+1$.}{$y=x+2$.}

\cauTN[id=C1-B3-D03-TH-039,muc=TH,loigiai={Xét hàm số $y=\dfrac{-x^{2}+5x+6}{x-1}$. Chia đa thức: $y=-x+4+\dfrac{10}{x - 1}$; $a=\lim\limits_{x\to +\infty}\dfrac{y}{x}=-1$, $b=\lim\limits_{x\to +\infty}(y+x)=4$ nên $y=-x+4$ là tiệm cận xiên (khi $x\to +\infty$); $a=\lim\limits_{x\to -\infty}\dfrac{y}{x}=-1$, $b=\lim\limits_{x\to -\infty}(y+x)=4$ nên $y=-x+4$ là tiệm cận xiên (khi $x\to -\infty$); $\lim\limits_{x\to {1}^{+}}y=+\infty$ nên $x=1$ là tiệm cận đứng.}]{Tiệm cận xiên của đồ thị hàm số $y=\dfrac{-x^{2}+5x+6}{x-1}$ đi qua điểm nào dưới đây?}
  {$M(2;2)$.}
  {$M(2;3)$.}{$M(2;1)$.}{$M(2;4)$.}

\cauTN[id=C1-B3-D03-TH-040,muc=TH,loigiai={Xét hàm số $y=-2x-4+\dfrac{3}{x-3}$. Chia đa thức: $y=-2x-4+\dfrac{3}{x - 3}$; $a=\lim\limits_{x\to +\infty}\dfrac{y}{x}=-2$, $b=\lim\limits_{x\to +\infty}(y+2x)=-4$ nên $y=-2x-4$ là tiệm cận xiên (khi $x\to +\infty$); $a=\lim\limits_{x\to -\infty}\dfrac{y}{x}=-2$, $b=\lim\limits_{x\to -\infty}(y+2x)=-4$ nên $y=-2x-4$ là tiệm cận xiên (khi $x\to -\infty$); $\lim\limits_{x\to {3}^{+}}y=+\infty$ nên $x=3$ là tiệm cận đứng.}]{Tiệm cận xiên của đồ thị hàm số $y=-2x-4+\dfrac{3}{x-3}$ đi qua điểm nào dưới đây?}
  {$M(-2;0)$.}
  {$M(-2;1)$.}{$M(-2;-1)$.}{$M(-2;2)$.}

\cauTN[id=C1-B3-D03-TH-041,muc=TH,loigiai={Xét hàm số $y=\dfrac{x^{3}-2x^{2}-2}{x^{2}+4}$. Chia đa thức: $y=x-2+\dfrac{6 - 4 x}{x^{2} + 4}$; $a=\lim\limits_{x\to +\infty}\dfrac{y}{x}=1$, $b=\lim\limits_{x\to +\infty}(y-x)=-2$ nên $y=x-2$ là tiệm cận xiên (khi $x\to +\infty$); $a=\lim\limits_{x\to -\infty}\dfrac{y}{x}=1$, $b=\lim\limits_{x\to -\infty}(y-x)=-2$ nên $y=x-2$ là tiệm cận xiên (khi $x\to -\infty$).}]{Tiệm cận xiên của đồ thị hàm số $y=\dfrac{x^{3}-2x^{2}-2}{x^{2}+4}$ là đường thẳng}
  {$y=x-2$.}
  {$y=x+2$.}{$y=-x-2$.}{$y=x-1$.}

\cauTN[id=C1-B3-D03-TH-042,muc=TH,loigiai={Xét hàm số $y=x-4+\dfrac{3}{x+2}$. Chia đa thức: $y=x-4+\dfrac{3}{x + 2}$; $a=\lim\limits_{x\to +\infty}\dfrac{y}{x}=1$, $b=\lim\limits_{x\to +\infty}(y-x)=-4$ nên $y=x-4$ là tiệm cận xiên (khi $x\to +\infty$); $a=\lim\limits_{x\to -\infty}\dfrac{y}{x}=1$, $b=\lim\limits_{x\to -\infty}(y-x)=-4$ nên $y=x-4$ là tiệm cận xiên (khi $x\to -\infty$); $\lim\limits_{x\to {-2}^{+}}y=+\infty$ nên $x=-2$ là tiệm cận đứng.}]{Tổng số đường tiệm cận (đứng, ngang, xiên) của đồ thị hàm số $y=x-4+\dfrac{3}{x+2}$ là}
  {$2$.}
  {$3$.}{$1$.}{$4$.}

\cauTN[id=C1-B3-D03-TH-043,muc=TH,loigiai={Xét hàm số $y=x+2+\dfrac{1}{x}$. Chia đa thức: $y=x+2+\dfrac{1}{x}$; $a=\lim\limits_{x\to +\infty}\dfrac{y}{x}=1$, $b=\lim\limits_{x\to +\infty}(y-x)=2$ nên $y=x+2$ là tiệm cận xiên (khi $x\to +\infty$); $a=\lim\limits_{x\to -\infty}\dfrac{y}{x}=1$, $b=\lim\limits_{x\to -\infty}(y-x)=2$ nên $y=x+2$ là tiệm cận xiên (khi $x\to -\infty$); $\lim\limits_{x\to {0}^{+}}y=+\infty$ nên $x=0$ là tiệm cận đứng.}]{Tổng số đường tiệm cận (đứng, ngang, xiên) của đồ thị hàm số $y=x+2+\dfrac{1}{x}$ là}
  {$2$.}
  {$3$.}{$1$.}{$4$.}

\cauTN[id=C1-B3-D03-TH-044,muc=TH,loigiai={Xét hàm số $y=\dfrac{3x^{2}-4x-6}{-x+2}$. Chia đa thức: $y=-3x-2+\dfrac{-2}{2 - x}$; $a=\lim\limits_{x\to +\infty}\dfrac{y}{x}=-3$, $b=\lim\limits_{x\to +\infty}(y+3x)=-2$ nên $y=-3x-2$ là tiệm cận xiên (khi $x\to +\infty$); $a=\lim\limits_{x\to -\infty}\dfrac{y}{x}=-3$, $b=\lim\limits_{x\to -\infty}(y+3x)=-2$ nên $y=-3x-2$ là tiệm cận xiên (khi $x\to -\infty$); $\lim\limits_{x\to {2}^{+}}y=+\infty$ nên $x=2$ là tiệm cận đứng.}]{Tiệm cận xiên của đồ thị hàm số $y=\dfrac{3x^{2}-4x-6}{-x+2}$ là đường thẳng}
  {$y=-3x-2$.}
  {$y=-3x+2$.}{$y=3x-2$.}{$y=-3x-1$.}

\cauTN[id=C1-B3-D03-TH-045,muc=TH,loigiai={Xét hàm số $y=-x-4+\dfrac{3}{x-3}$. Chia đa thức: $y=-x-4+\dfrac{3}{x - 3}$; $a=\lim\limits_{x\to +\infty}\dfrac{y}{x}=-1$, $b=\lim\limits_{x\to +\infty}(y+x)=-4$ nên $y=-x-4$ là tiệm cận xiên (khi $x\to +\infty$); $a=\lim\limits_{x\to -\infty}\dfrac{y}{x}=-1$, $b=\lim\limits_{x\to -\infty}(y+x)=-4$ nên $y=-x-4$ là tiệm cận xiên (khi $x\to -\infty$); $\lim\limits_{x\to {3}^{+}}y=+\infty$ nên $x=3$ là tiệm cận đứng.}]{Tiệm cận xiên của đồ thị hàm số $y=-x-4+\dfrac{3}{x-3}$ đi qua điểm nào dưới đây?}
  {$M(3;-7)$.}
  {$M(3;-6)$.}{$M(3;-8)$.}{$M(3;-5)$.}

\cauTN[id=C1-B3-D03-TH-046,muc=TH,loigiai={Xét hàm số $y=\dfrac{2x^{3}+2x^{2}}{x^{2}+2}$. Chia đa thức: $y=2x+2+\dfrac{- 4 x - 4}{x^{2} + 2}$; $a=\lim\limits_{x\to +\infty}\dfrac{y}{x}=2$, $b=\lim\limits_{x\to +\infty}(y-2x)=2$ nên $y=2x+2$ là tiệm cận xiên (khi $x\to +\infty$); $a=\lim\limits_{x\to -\infty}\dfrac{y}{x}=2$, $b=\lim\limits_{x\to -\infty}(y-2x)=2$ nên $y=2x+2$ là tiệm cận xiên (khi $x\to -\infty$).}]{Tiệm cận xiên của đồ thị hàm số $y=\dfrac{2x^{3}+2x^{2}}{x^{2}+2}$ là đường thẳng}
  {$y=2x+2$.}
  {$y=2x-2$.}{$y=-2x+2$.}{$y=2x+3$.}

\cauTN[id=C1-B3-D03-TH-047,muc=TH,loigiai={Xét hàm số $y=\dfrac{2x^{2}-x+2}{2x+3}$. Chia đa thức: $y=x-2+\dfrac{8}{2 x + 3}$; $a=\lim\limits_{x\to +\infty}\dfrac{y}{x}=1$, $b=\lim\limits_{x\to +\infty}(y-x)=-2$ nên $y=x-2$ là tiệm cận xiên (khi $x\to +\infty$); $a=\lim\limits_{x\to -\infty}\dfrac{y}{x}=1$, $b=\lim\limits_{x\to -\infty}(y-x)=-2$ nên $y=x-2$ là tiệm cận xiên (khi $x\to -\infty$); $\lim\limits_{x\to {-\dfrac{3}{2}}^{+}}y=+\infty$ nên $x=-\dfrac{3}{2}$ là tiệm cận đứng.}]{Tiệm cận xiên của đồ thị hàm số $y=\dfrac{2x^{2}-x+2}{2x+3}$ đi qua điểm nào dưới đây?}
  {$M(-2;-4)$.}
  {$M(-2;-3)$.}{$M(-2;-5)$.}{$M(-2;-2)$.}

\cauTN[id=C1-B3-D03-TH-048,muc=TH,loigiai={Xét hàm số $y=\dfrac{x^{2}-3x-5}{2x-1}$. Chia đa thức: $y=\dfrac{1}{2}x-\dfrac{5}{4}+\dfrac{- \dfrac{25}{4}}{2 x - 1}$; $a=\lim\limits_{x\to +\infty}\dfrac{y}{x}=\dfrac{1}{2}$, $b=\lim\limits_{x\to +\infty}(y-\dfrac{1}{2}x)=-\dfrac{5}{4}$ nên $y=\dfrac{1}{2}x-\dfrac{5}{4}$ là tiệm cận xiên (khi $x\to +\infty$); $a=\lim\limits_{x\to -\infty}\dfrac{y}{x}=\dfrac{1}{2}$, $b=\lim\limits_{x\to -\infty}(y-\dfrac{1}{2}x)=-\dfrac{5}{4}$ nên $y=\dfrac{1}{2}x-\dfrac{5}{4}$ là tiệm cận xiên (khi $x\to -\infty$); $\lim\limits_{x\to {\dfrac{1}{2}}^{+}}y=-\infty$ nên $x=\dfrac{1}{2}$ là tiệm cận đứng.}]{Tiệm cận xiên của đồ thị hàm số $y=\dfrac{x^{2}-3x-5}{2x-1}$ là đường thẳng}
  {$y=\dfrac{1}{2}x-\dfrac{5}{4}$.}
  {$y=\dfrac{1}{2}x+\dfrac{5}{4}$.}{$y=-\dfrac{1}{2}x-\dfrac{5}{4}$.}{$y=\dfrac{1}{2}x-\dfrac{1}{4}$.}

\cauDS[id=C1-B3-D03-TH-049,muc=TH,loigiai={Xét hàm số $y=2x-4+\dfrac{4}{x+2}$. Chia đa thức: $y=2x-4+\dfrac{4}{x + 2}$; $a=\lim\limits_{x\to +\infty}\dfrac{y}{x}=2$, $b=\lim\limits_{x\to +\infty}(y-2x)=-4$ nên $y=2x-4$ là tiệm cận xiên (khi $x\to +\infty$); $a=\lim\limits_{x\to -\infty}\dfrac{y}{x}=2$, $b=\lim\limits_{x\to -\infty}(y-2x)=-4$ nên $y=2x-4$ là tiệm cận xiên (khi $x\to -\infty$); $\lim\limits_{x\to {-2}^{+}}y=+\infty$ nên $x=-2$ là tiệm cận đứng. a) Đúng. b) Sai. c) Đúng. d) Đúng.}]{Cho hàm số $y=2x-4+\dfrac{4}{x+2}$. Xét tính đúng sai của các mệnh đề sau:}{%
  \y{Đường thẳng $y=2x-4$ là tiệm cận xiên của đồ thị hàm số.}{d}
  \y{Đường thẳng $y=2x-2$ là tiệm cận xiên của đồ thị hàm số.}{s}
  \y{Đồ thị hàm số có tiệm cận đứng $x=-2$.}{d}
  \y{Đồ thị hàm số có tất cả $2$ đường tiệm cận.}{d}}

\cauDS[id=C1-B3-D03-TH-050,muc=TH,loigiai={Xét hàm số $y=3x-1-\dfrac{1}{x-1}$. Chia đa thức: $y=3x-1+\dfrac{-1}{x - 1}$; $a=\lim\limits_{x\to +\infty}\dfrac{y}{x}=3$, $b=\lim\limits_{x\to +\infty}(y-3x)=-1$ nên $y=3x-1$ là tiệm cận xiên (khi $x\to +\infty$); $a=\lim\limits_{x\to -\infty}\dfrac{y}{x}=3$, $b=\lim\limits_{x\to -\infty}(y-3x)=-1$ nên $y=3x-1$ là tiệm cận xiên (khi $x\to -\infty$); $\lim\limits_{x\to {1}^{+}}y=-\infty$ nên $x=1$ là tiệm cận đứng. a) Sai. b) Sai: giới hạn của $y$ tại vô cực không hữu hạn. c) Đúng. d) Đúng.}]{Cho hàm số $y=3x-1-\dfrac{1}{x-1}$. Xét tính đúng sai của các mệnh đề sau:}{%
  \y{Đường thẳng $y=3x-2$ là tiệm cận xiên của đồ thị hàm số.}{s}
  \y{Đồ thị hàm số có tiệm cận ngang.}{s}
  \y{Đồ thị hàm số có tiệm cận đứng $x=1$.}{d}
  \y{Đồ thị hàm số có tất cả $2$ đường tiệm cận.}{d}}

\cauDS[id=C1-B3-D03-TH-051,muc=TH,loigiai={Xét hàm số $y=-2x+2+\dfrac{3}{x-1}$. Chia đa thức: $y=-2x+2+\dfrac{3}{x - 1}$; $a=\lim\limits_{x\to +\infty}\dfrac{y}{x}=-2$, $b=\lim\limits_{x\to +\infty}(y+2x)=2$ nên $y=-2x+2$ là tiệm cận xiên (khi $x\to +\infty$); $a=\lim\limits_{x\to -\infty}\dfrac{y}{x}=-2$, $b=\lim\limits_{x\to -\infty}(y+2x)=2$ nên $y=-2x+2$ là tiệm cận xiên (khi $x\to -\infty$); $\lim\limits_{x\to {1}^{+}}y=+\infty$ nên $x=1$ là tiệm cận đứng. a) Đúng. b) Sai: giới hạn của $y$ tại vô cực không hữu hạn. c) Đúng. d) Sai: có $2$ đường tiệm cận.}]{Cho hàm số $y=-2x+2+\dfrac{3}{x-1}$. Xét tính đúng sai của các mệnh đề sau:}{%
  \y{Giao điểm của tiệm cận đứng và tiệm cận xiên là $I(1;0)$.}{d}
  \y{Đồ thị hàm số có tiệm cận ngang.}{s}
  \y{Đường thẳng $y=-2x+2$ là tiệm cận xiên của đồ thị hàm số.}{d}
  \y{Đồ thị hàm số có tất cả $3$ đường tiệm cận.}{s}}

\cauDS[id=C1-B3-D03-TH-052,muc=TH,loigiai={Xét hàm số $y=\dfrac{2x^{3}+2x^{2}-1}{x^{2}+1}$. Chia đa thức: $y=2x+2+\dfrac{- 2 x - 3}{x^{2} + 1}$; $a=\lim\limits_{x\to +\infty}\dfrac{y}{x}=2$, $b=\lim\limits_{x\to +\infty}(y-2x)=2$ nên $y=2x+2$ là tiệm cận xiên (khi $x\to +\infty$); $a=\lim\limits_{x\to -\infty}\dfrac{y}{x}=2$, $b=\lim\limits_{x\to -\infty}(y-2x)=2$ nên $y=2x+2$ là tiệm cận xiên (khi $x\to -\infty$). a) Đúng. b) Đúng. c) Sai. d) Sai: có $1$ đường tiệm cận.}]{Cho hàm số $y=\dfrac{2x^{3}+2x^{2}-1}{x^{2}+1}$. Xét tính đúng sai của các mệnh đề sau:}{%
  \y{Đường thẳng $y=2x+2$ là tiệm cận xiên của đồ thị hàm số.}{d}
  \y{Đồ thị hàm số không có tiệm cận đứng.}{d}
  \y{Đường thẳng $y=2x+4$ là tiệm cận xiên của đồ thị hàm số.}{s}
  \y{Đồ thị hàm số có tất cả $2$ đường tiệm cận.}{s}}

\cauDS[id=C1-B3-D03-TH-053,muc=TH,loigiai={Xét hàm số $y=\dfrac{x^{3}+x^{2}+3}{x^{2}+1}$. Chia đa thức: $y=x+1+\dfrac{2 - x}{x^{2} + 1}$; $a=\lim\limits_{x\to +\infty}\dfrac{y}{x}=1$, $b=\lim\limits_{x\to +\infty}(y-x)=1$ nên $y=x+1$ là tiệm cận xiên (khi $x\to +\infty$); $a=\lim\limits_{x\to -\infty}\dfrac{y}{x}=1$, $b=\lim\limits_{x\to -\infty}(y-x)=1$ nên $y=x+1$ là tiệm cận xiên (khi $x\to -\infty$). a) Đúng. b) Đúng. c) Đúng. d) Sai: giới hạn của $y$ tại vô cực không hữu hạn.}]{Cho hàm số $y=\dfrac{x^{3}+x^{2}+3}{x^{2}+1}$. Xét tính đúng sai của các mệnh đề sau:}{%
  \y{Đường thẳng $y=x+1$ là tiệm cận xiên của đồ thị hàm số.}{d}
  \y{Đồ thị hàm số không có tiệm cận đứng.}{d}
  \y{Đồ thị hàm số có tất cả $1$ đường tiệm cận.}{d}
  \y{Đồ thị hàm số có tiệm cận ngang.}{s}}

\cauDS[id=C1-B3-D03-TH-054,muc=TH,loigiai={Xét hàm số $y=\dfrac{-x^{3}+3x^{2}+1}{x^{2}+1}$. Chia đa thức: $y=-x+3+\dfrac{x - 2}{x^{2} + 1}$; $a=\lim\limits_{x\to +\infty}\dfrac{y}{x}=-1$, $b=\lim\limits_{x\to +\infty}(y+x)=3$ nên $y=-x+3$ là tiệm cận xiên (khi $x\to +\infty$); $a=\lim\limits_{x\to -\infty}\dfrac{y}{x}=-1$, $b=\lim\limits_{x\to -\infty}(y+x)=3$ nên $y=-x+3$ là tiệm cận xiên (khi $x\to -\infty$). a) Sai. b) Sai: có $1$ đường tiệm cận. c) Đúng. d) Đúng.}]{Cho hàm số $y=\dfrac{-x^{3}+3x^{2}+1}{x^{2}+1}$. Xét tính đúng sai của các mệnh đề sau:}{%
  \y{Đường thẳng $y=-x+2$ là tiệm cận xiên của đồ thị hàm số.}{s}
  \y{Đồ thị hàm số có tất cả $2$ đường tiệm cận.}{s}
  \y{Đường thẳng $y=-x+3$ là tiệm cận xiên của đồ thị hàm số.}{d}
  \y{Đồ thị hàm số có tất cả $1$ đường tiệm cận.}{d}}

\cauDS[id=C1-B3-D03-TH-055,muc=TH,loigiai={Xét hàm số $y=\dfrac{2x^{2}-2x-2}{-x-2}$. Chia đa thức: $y=-2x+6+\dfrac{10}{- x - 2}$; $a=\lim\limits_{x\to +\infty}\dfrac{y}{x}=-2$, $b=\lim\limits_{x\to +\infty}(y+2x)=6$ nên $y=-2x+6$ là tiệm cận xiên (khi $x\to +\infty$); $a=\lim\limits_{x\to -\infty}\dfrac{y}{x}=-2$, $b=\lim\limits_{x\to -\infty}(y+2x)=6$ nên $y=-2x+6$ là tiệm cận xiên (khi $x\to -\infty$); $\lim\limits_{x\to {-2}^{+}}y=-\infty$ nên $x=-2$ là tiệm cận đứng. a) Đúng. b) Đúng. c) Sai: có $2$ đường tiệm cận. d) Sai: tiệm cận xiên là $y=-2x+6$.}]{Cho hàm số $y=\dfrac{2x^{2}-2x-2}{-x-2}$. Xét tính đúng sai của các mệnh đề sau:}{%
  \y{Đường thẳng $y=-2x+6$ là tiệm cận xiên của đồ thị hàm số.}{d}
  \y{Đồ thị hàm số có tất cả $2$ đường tiệm cận.}{d}
  \y{Đồ thị hàm số có tất cả $3$ đường tiệm cận.}{s}
  \y{Đường thẳng $y=-2x$ là tiệm cận xiên của đồ thị hàm số.}{s}}

\cauDS[id=C1-B3-D03-TH-056,muc=TH,loigiai={Xét hàm số $y=3x+\dfrac{4}{x-2}$. Chia đa thức: $y=3x+\dfrac{4}{x - 2}$; $a=\lim\limits_{x\to +\infty}\dfrac{y}{x}=3$, $b=\lim\limits_{x\to +\infty}(y-3x)=0$ nên $y=3x$ là tiệm cận xiên (khi $x\to +\infty$); $a=\lim\limits_{x\to -\infty}\dfrac{y}{x}=3$, $b=\lim\limits_{x\to -\infty}(y-3x)=0$ nên $y=3x$ là tiệm cận xiên (khi $x\to -\infty$); $\lim\limits_{x\to {2}^{+}}y=+\infty$ nên $x=2$ là tiệm cận đứng. a) Đúng. b) Sai. c) Đúng. d) Đúng.}]{Cho hàm số $y=3x+\dfrac{4}{x-2}$. Xét tính đúng sai của các mệnh đề sau:}{%
  \y{Đường thẳng $y=3x$ là tiệm cận xiên của đồ thị hàm số.}{d}
  \y{Đường thẳng $y=3x+2$ là tiệm cận xiên của đồ thị hàm số.}{s}
  \y{Đồ thị hàm số có tất cả $2$ đường tiệm cận.}{d}
  \y{Giao điểm của tiệm cận đứng và tiệm cận xiên là $I(2;6)$.}{d}}

\cauDS[id=C1-B3-D03-TH-057,muc=TH,loigiai={Xét hàm số $y=\dfrac{2x^{2}+5x+1}{x-4}$. Chia đa thức: $y=2x+13+\dfrac{53}{x - 4}$; $a=\lim\limits_{x\to +\infty}\dfrac{y}{x}=2$, $b=\lim\limits_{x\to +\infty}(y-2x)=13$ nên $y=2x+13$ là tiệm cận xiên (khi $x\to +\infty$); $a=\lim\limits_{x\to -\infty}\dfrac{y}{x}=2$, $b=\lim\limits_{x\to -\infty}(y-2x)=13$ nên $y=2x+13$ là tiệm cận xiên (khi $x\to -\infty$); $\lim\limits_{x\to {4}^{+}}y=+\infty$ nên $x=4$ là tiệm cận đứng. a) Sai: giới hạn của $y$ tại vô cực không hữu hạn. b) Đúng. c) Sai. d) Đúng.}]{Cho hàm số $y=\dfrac{2x^{2}+5x+1}{x-4}$. Xét tính đúng sai của các mệnh đề sau:}{%
  \y{Đồ thị hàm số có tiệm cận ngang.}{s}
  \y{Đường thẳng $y=2x+13$ là tiệm cận xiên của đồ thị hàm số.}{d}
  \y{Đường thẳng $y=2x+14$ là tiệm cận xiên của đồ thị hàm số.}{s}
  \y{Đồ thị hàm số có tiệm cận đứng $x=4$.}{d}}

\cauDS[id=C1-B3-D03-TH-058,muc=TH,loigiai={Xét hàm số $y=\dfrac{3x^{2}-2x-6}{-x+3}$. Chia đa thức: $y=-3x-7+\dfrac{15}{3 - x}$; $a=\lim\limits_{x\to +\infty}\dfrac{y}{x}=-3$, $b=\lim\limits_{x\to +\infty}(y+3x)=-7$ nên $y=-3x-7$ là tiệm cận xiên (khi $x\to +\infty$); $a=\lim\limits_{x\to -\infty}\dfrac{y}{x}=-3$, $b=\lim\limits_{x\to -\infty}(y+3x)=-7$ nên $y=-3x-7$ là tiệm cận xiên (khi $x\to -\infty$); $\lim\limits_{x\to {3}^{+}}y=-\infty$ nên $x=3$ là tiệm cận đứng. a) Sai. b) Đúng. c) Đúng. d) Sai: tiệm cận xiên là $y=-3x-7$.}]{Cho hàm số $y=\dfrac{3x^{2}-2x-6}{-x+3}$. Xét tính đúng sai của các mệnh đề sau:}{%
  \y{Đường thẳng $y=-3x-8$ là tiệm cận xiên của đồ thị hàm số.}{s}
  \y{Đường thẳng $y=-3x-7$ là tiệm cận xiên của đồ thị hàm số.}{d}
  \y{Đồ thị hàm số có tiệm cận đứng $x=3$.}{d}
  \y{Đường thẳng $y=-3x$ là tiệm cận xiên của đồ thị hàm số.}{s}}

\cauDS[id=C1-B3-D03-TH-059,muc=TH,loigiai={Xét hàm số $y=\dfrac{2x^{2}+5x+3}{x+2}$. Chia đa thức: $y=2x+1+\dfrac{1}{x + 2}$; $a=\lim\limits_{x\to +\infty}\dfrac{y}{x}=2$, $b=\lim\limits_{x\to +\infty}(y-2x)=1$ nên $y=2x+1$ là tiệm cận xiên (khi $x\to +\infty$); $a=\lim\limits_{x\to -\infty}\dfrac{y}{x}=2$, $b=\lim\limits_{x\to -\infty}(y-2x)=1$ nên $y=2x+1$ là tiệm cận xiên (khi $x\to -\infty$); $\lim\limits_{x\to {-2}^{+}}y=+\infty$ nên $x=-2$ là tiệm cận đứng. a) Sai: giới hạn của $y$ tại vô cực không hữu hạn. b) Đúng. c) Đúng. d) Sai.}]{Cho hàm số $y=\dfrac{2x^{2}+5x+3}{x+2}$. Xét tính đúng sai của các mệnh đề sau:}{%
  \y{Đồ thị hàm số có tiệm cận ngang.}{s}
  \y{Đồ thị hàm số có tất cả $2$ đường tiệm cận.}{d}
  \y{Đồ thị hàm số có tiệm cận đứng $x=-2$.}{d}
  \y{Đường thẳng $y=2x$ là tiệm cận xiên của đồ thị hàm số.}{s}}

\cauDS[id=C1-B3-D03-TH-060,muc=TH,loigiai={Xét hàm số $y=-2x-2-\dfrac{2}{x-1}$. Chia đa thức: $y=-2x-2+\dfrac{-2}{x - 1}$; $a=\lim\limits_{x\to +\infty}\dfrac{y}{x}=-2$, $b=\lim\limits_{x\to +\infty}(y+2x)=-2$ nên $y=-2x-2$ là tiệm cận xiên (khi $x\to +\infty$); $a=\lim\limits_{x\to -\infty}\dfrac{y}{x}=-2$, $b=\lim\limits_{x\to -\infty}(y+2x)=-2$ nên $y=-2x-2$ là tiệm cận xiên (khi $x\to -\infty$); $\lim\limits_{x\to {1}^{+}}y=-\infty$ nên $x=1$ là tiệm cận đứng. a) Sai. b) Sai: giới hạn của $y$ tại vô cực không hữu hạn. c) Đúng. d) Đúng.}]{Cho hàm số $y=-2x-2-\dfrac{2}{x-1}$. Xét tính đúng sai của các mệnh đề sau:}{%
  \y{Đường thẳng $y=-2x-1$ là tiệm cận xiên của đồ thị hàm số.}{s}
  \y{Đồ thị hàm số có tiệm cận ngang.}{s}
  \y{Đồ thị hàm số có tất cả $2$ đường tiệm cận.}{d}
  \y{Đường thẳng $y=-2x-2$ là tiệm cận xiên của đồ thị hàm số.}{d}}

\cauDS[id=C1-B3-D03-TH-061,muc=TH,loigiai={Xét hàm số $y=\dfrac{x^{3}+x^{2}-3}{x^{2}+1}$. Chia đa thức: $y=x+1+\dfrac{- x - 4}{x^{2} + 1}$; $a=\lim\limits_{x\to +\infty}\dfrac{y}{x}=1$, $b=\lim\limits_{x\to +\infty}(y-x)=1$ nên $y=x+1$ là tiệm cận xiên (khi $x\to +\infty$); $a=\lim\limits_{x\to -\infty}\dfrac{y}{x}=1$, $b=\lim\limits_{x\to -\infty}(y-x)=1$ nên $y=x+1$ là tiệm cận xiên (khi $x\to -\infty$). a) Sai: giới hạn của $y$ tại vô cực không hữu hạn. b) Đúng. c) Sai: tiệm cận xiên là $y=x+1$. d) Sai: có $1$ đường tiệm cận.}]{Cho hàm số $y=\dfrac{x^{3}+x^{2}-3}{x^{2}+1}$. Xét tính đúng sai của các mệnh đề sau:}{%
  \y{Đồ thị hàm số có tiệm cận ngang.}{s}
  \y{Đồ thị hàm số có tất cả $1$ đường tiệm cận.}{d}
  \y{Đường thẳng $y=x$ là tiệm cận xiên của đồ thị hàm số.}{s}
  \y{Đồ thị hàm số có tất cả $2$ đường tiệm cận.}{s}}

\cauDS[id=C1-B3-D03-TH-062,muc=TH,loigiai={Xét hàm số $y=\dfrac{2x^{2}+2x-3}{x+2}$. Chia đa thức: $y=2x-2+\dfrac{1}{x + 2}$; $a=\lim\limits_{x\to +\infty}\dfrac{y}{x}=2$, $b=\lim\limits_{x\to +\infty}(y-2x)=-2$ nên $y=2x-2$ là tiệm cận xiên (khi $x\to +\infty$); $a=\lim\limits_{x\to -\infty}\dfrac{y}{x}=2$, $b=\lim\limits_{x\to -\infty}(y-2x)=-2$ nên $y=2x-2$ là tiệm cận xiên (khi $x\to -\infty$); $\lim\limits_{x\to {-2}^{+}}y=+\infty$ nên $x=-2$ là tiệm cận đứng. a) Đúng. b) Sai. c) Sai: giới hạn của $y$ tại vô cực không hữu hạn. d) Đúng.}]{Cho hàm số $y=\dfrac{2x^{2}+2x-3}{x+2}$. Xét tính đúng sai của các mệnh đề sau:}{%
  \y{Đường thẳng $y=2x-2$ là tiệm cận xiên của đồ thị hàm số.}{d}
  \y{Đường thẳng $y=2x-3$ là tiệm cận xiên của đồ thị hàm số.}{s}
  \y{Đồ thị hàm số có tiệm cận ngang.}{s}
  \y{Giao điểm của tiệm cận đứng và tiệm cận xiên là $I(-2;-6)$.}{d}}

\cauDS[id=C1-B3-D03-TH-063,muc=TH,loigiai={Xét hàm số $y=\dfrac{x^{3}-2x^{2}-3}{x^{2}-4}$. Chia đa thức: $y=x-2+\dfrac{4 x - 11}{x^{2} - 4}$; $a=\lim\limits_{x\to +\infty}\dfrac{y}{x}=1$, $b=\lim\limits_{x\to +\infty}(y-x)=-2$ nên $y=x-2$ là tiệm cận xiên (khi $x\to +\infty$); $a=\lim\limits_{x\to -\infty}\dfrac{y}{x}=1$, $b=\lim\limits_{x\to -\infty}(y-x)=-2$ nên $y=x-2$ là tiệm cận xiên (khi $x\to -\infty$); $\lim\limits_{x\to {-2}^{+}}y=+\infty$ nên $x=-2$ là tiệm cận đứng; $\lim\limits_{x\to {2}^{+}}y=-\infty$ nên $x=2$ là tiệm cận đứng. a) Đúng. b) Đúng. c) Đúng. d) Sai: giới hạn của $y$ tại vô cực không hữu hạn.}]{Cho hàm số $y=\dfrac{x^{3}-2x^{2}-3}{x^{2}-4}$. Xét tính đúng sai của các mệnh đề sau:}{%
  \y{Đồ thị hàm số có tất cả $3$ đường tiệm cận.}{d}
  \y{Đồ thị hàm số có tiệm cận đứng $x=-2$.}{d}
  \y{Đường thẳng $y=x-2$ là tiệm cận xiên của đồ thị hàm số.}{d}
  \y{Đồ thị hàm số có tiệm cận ngang.}{s}}

\cauDS[id=C1-B3-D03-TH-064,muc=TH,loigiai={Xét hàm số $y=-x+1+\dfrac{2}{x-1}$. Chia đa thức: $y=-x+1+\dfrac{2}{x - 1}$; $a=\lim\limits_{x\to +\infty}\dfrac{y}{x}=-1$, $b=\lim\limits_{x\to +\infty}(y+x)=1$ nên $y=-x+1$ là tiệm cận xiên (khi $x\to +\infty$); $a=\lim\limits_{x\to -\infty}\dfrac{y}{x}=-1$, $b=\lim\limits_{x\to -\infty}(y+x)=1$ nên $y=-x+1$ là tiệm cận xiên (khi $x\to -\infty$); $\lim\limits_{x\to {1}^{+}}y=+\infty$ nên $x=1$ là tiệm cận đứng. a) Sai: có $2$ đường tiệm cận. b) Sai. c) Đúng. d) Sai: tiệm cận xiên là $y=-x+1$.}]{Cho hàm số $y=-x+1+\dfrac{2}{x-1}$. Xét tính đúng sai của các mệnh đề sau:}{%
  \y{Đồ thị hàm số có tất cả $3$ đường tiệm cận.}{s}
  \y{Đường thẳng $y=-x+3$ là tiệm cận xiên của đồ thị hàm số.}{s}
  \y{Đồ thị hàm số có tất cả $2$ đường tiệm cận.}{d}
  \y{Đường thẳng $y=-x$ là tiệm cận xiên của đồ thị hàm số.}{s}}

\cauTLN[id=C1-B3-D03-TH-065,muc=TH,loigiai={Xét hàm số $y=\dfrac{3x^{2}+4x-3}{x-3}$. Chia đa thức: $y=3x+13+\dfrac{36}{x - 3}$; $a=\lim\limits_{x\to +\infty}\dfrac{y}{x}=3$, $b=\lim\limits_{x\to +\infty}(y-3x)=13$ nên $y=3x+13$ là tiệm cận xiên (khi $x\to +\infty$); $a=\lim\limits_{x\to -\infty}\dfrac{y}{x}=3$, $b=\lim\limits_{x\to -\infty}(y-3x)=13$ nên $y=3x+13$ là tiệm cận xiên (khi $x\to -\infty$); $\lim\limits_{x\to {3}^{+}}y=+\infty$ nên $x=3$ là tiệm cận đứng. Kết quả là $16$.}]{Tiệm cận xiên của đồ thị hàm số $y=\dfrac{3x^{2}+4x-3}{x-3}$ có phương trình $y=ax+b$. Tính $a+b$.}{16}

\cauTLN[id=C1-B3-D03-TH-066,muc=TH,loigiai={Xét hàm số $y=\dfrac{3x^{2}+6}{x+3}$. Chia đa thức: $y=3x-9+\dfrac{33}{x + 3}$; $a=\lim\limits_{x\to +\infty}\dfrac{y}{x}=3$, $b=\lim\limits_{x\to +\infty}(y-3x)=-9$ nên $y=3x-9$ là tiệm cận xiên (khi $x\to +\infty$); $a=\lim\limits_{x\to -\infty}\dfrac{y}{x}=3$, $b=\lim\limits_{x\to -\infty}(y-3x)=-9$ nên $y=3x-9$ là tiệm cận xiên (khi $x\to -\infty$); $\lim\limits_{x\to {-3}^{+}}y=+\infty$ nên $x=-3$ là tiệm cận đứng. Kết quả là $-9$.}]{Tiệm cận xiên của đồ thị hàm số $y=\dfrac{3x^{2}+6}{x+3}$ cắt trục tung tại điểm có tung độ bằng bao nhiêu?}{-9}

\cauTLN[id=C1-B3-D03-TH-067,muc=TH,loigiai={Xét hàm số $y=3x+4+\dfrac{4}{x-3}$. Chia đa thức: $y=3x+4+\dfrac{4}{x - 3}$; $a=\lim\limits_{x\to +\infty}\dfrac{y}{x}=3$, $b=\lim\limits_{x\to +\infty}(y-3x)=4$ nên $y=3x+4$ là tiệm cận xiên (khi $x\to +\infty$); $a=\lim\limits_{x\to -\infty}\dfrac{y}{x}=3$, $b=\lim\limits_{x\to -\infty}(y-3x)=4$ nên $y=3x+4$ là tiệm cận xiên (khi $x\to -\infty$); $\lim\limits_{x\to {3}^{+}}y=+\infty$ nên $x=3$ là tiệm cận đứng. Kết quả là $4$.}]{Tiệm cận xiên của đồ thị hàm số $y=3x+4+\dfrac{4}{x-3}$ cắt trục tung tại điểm có tung độ bằng bao nhiêu?}{4}

\cauTLN[id=C1-B3-D03-TH-068,muc=TH,loigiai={Xét hàm số $y=\dfrac{-x^{2}+4x+2}{2x-4}$. Chia đa thức: $y=-\dfrac{1}{2}x+1+\dfrac{6}{2 x - 4}$; $a=\lim\limits_{x\to +\infty}\dfrac{y}{x}=-\dfrac{1}{2}$, $b=\lim\limits_{x\to +\infty}(y+\dfrac{1}{2}x)=1$ nên $y=-\dfrac{1}{2}x+1$ là tiệm cận xiên (khi $x\to +\infty$); $a=\lim\limits_{x\to -\infty}\dfrac{y}{x}=-\dfrac{1}{2}$, $b=\lim\limits_{x\to -\infty}(y+\dfrac{1}{2}x)=1$ nên $y=-\dfrac{1}{2}x+1$ là tiệm cận xiên (khi $x\to -\infty$); $\lim\limits_{x\to {2}^{+}}y=+\infty$ nên $x=2$ là tiệm cận đứng. Kết quả là $-2$.}]{Tiệm cận xiên của đồ thị hàm số $y=\dfrac{-x^{2}+4x+2}{2x-4}$ có phương trình $y=ax+b$. Tính $2a-b$.}{-2}

\cauTLN[id=C1-B3-D03-TH-069,muc=TH,loigiai={Xét hàm số $y=\dfrac{2x^{2}+3x+5}{2x-4}$. Chia đa thức: $y=x+\dfrac{7}{2}+\dfrac{19}{2 x - 4}$; $a=\lim\limits_{x\to +\infty}\dfrac{y}{x}=1$, $b=\lim\limits_{x\to +\infty}(y-x)=\dfrac{7}{2}$ nên $y=x+\dfrac{7}{2}$ là tiệm cận xiên (khi $x\to +\infty$); $a=\lim\limits_{x\to -\infty}\dfrac{y}{x}=1$, $b=\lim\limits_{x\to -\infty}(y-x)=\dfrac{7}{2}$ nên $y=x+\dfrac{7}{2}$ là tiệm cận xiên (khi $x\to -\infty$); $\lim\limits_{x\to {2}^{+}}y=+\infty$ nên $x=2$ là tiệm cận đứng. Kết quả là $\dfrac{15}{2}$.}]{Giao điểm của tiệm cận đứng và tiệm cận xiên của đồ thị hàm số $y=\dfrac{2x^{2}+3x+5}{2x-4}$ là $I(x_0;y_0)$. Tính $x_0+y_0$.}{7,5}

\cauTLN[id=C1-B3-D03-TH-070,muc=TH,loigiai={Xét hàm số $y=\dfrac{-x^{2}-3x-6}{x-4}$. Chia đa thức: $y=-x-7+\dfrac{-34}{x - 4}$; $a=\lim\limits_{x\to +\infty}\dfrac{y}{x}=-1$, $b=\lim\limits_{x\to +\infty}(y+x)=-7$ nên $y=-x-7$ là tiệm cận xiên (khi $x\to +\infty$); $a=\lim\limits_{x\to -\infty}\dfrac{y}{x}=-1$, $b=\lim\limits_{x\to -\infty}(y+x)=-7$ nên $y=-x-7$ là tiệm cận xiên (khi $x\to -\infty$); $\lim\limits_{x\to {4}^{+}}y=-\infty$ nên $x=4$ là tiệm cận đứng. Kết quả là $-7$.}]{Giao điểm của tiệm cận đứng và tiệm cận xiên của đồ thị hàm số $y=\dfrac{-x^{2}-3x-6}{x-4}$ là $I(x_0;y_0)$. Tính $x_0+y_0$.}{-7}

\cauTLN[id=C1-B3-D03-TH-071,muc=TH,loigiai={Xét hàm số $y=\dfrac{2x^{3}+3}{x^{2}+2}$. Chia đa thức: $y=2x+\dfrac{3 - 4 x}{x^{2} + 2}$; $a=\lim\limits_{x\to +\infty}\dfrac{y}{x}=2$, $b=\lim\limits_{x\to +\infty}(y-2x)=0$ nên $y=2x$ là tiệm cận xiên (khi $x\to +\infty$); $a=\lim\limits_{x\to -\infty}\dfrac{y}{x}=2$, $b=\lim\limits_{x\to -\infty}(y-2x)=0$ nên $y=2x$ là tiệm cận xiên (khi $x\to -\infty$). Kết quả là $0$.}]{Tiệm cận xiên của đồ thị hàm số $y=\dfrac{2x^{3}+3}{x^{2}+2}$ cắt trục tung tại điểm có tung độ bằng bao nhiêu?}{0}

\cauTLN[id=C1-B3-D03-TH-072,muc=TH,loigiai={Xét hàm số $y=\dfrac{2x^{3}+x^{2}+1}{x^{2}-4}$. Chia đa thức: $y=2x+1+\dfrac{8 x + 5}{x^{2} - 4}$; $a=\lim\limits_{x\to +\infty}\dfrac{y}{x}=2$, $b=\lim\limits_{x\to +\infty}(y-2x)=1$ nên $y=2x+1$ là tiệm cận xiên (khi $x\to +\infty$); $a=\lim\limits_{x\to -\infty}\dfrac{y}{x}=2$, $b=\lim\limits_{x\to -\infty}(y-2x)=1$ nên $y=2x+1$ là tiệm cận xiên (khi $x\to -\infty$); $\lim\limits_{x\to {-2}^{+}}y=+\infty$ nên $x=-2$ là tiệm cận đứng; $\lim\limits_{x\to {2}^{+}}y=+\infty$ nên $x=2$ là tiệm cận đứng. Kết quả là $3$.}]{Tiệm cận xiên của đồ thị hàm số $y=\dfrac{2x^{3}+x^{2}+1}{x^{2}-4}$ có phương trình $y=ax+b$. Tính $2a-b$.}{3}

\cauTLN[id=C1-B3-D03-TH-073,muc=TH,loigiai={Xét hàm số $y=\dfrac{2x^{2}-2x-6}{x+4}$. Chia đa thức: $y=2x-10+\dfrac{34}{x + 4}$; $a=\lim\limits_{x\to +\infty}\dfrac{y}{x}=2$, $b=\lim\limits_{x\to +\infty}(y-2x)=-10$ nên $y=2x-10$ là tiệm cận xiên (khi $x\to +\infty$); $a=\lim\limits_{x\to -\infty}\dfrac{y}{x}=2$, $b=\lim\limits_{x\to -\infty}(y-2x)=-10$ nên $y=2x-10$ là tiệm cận xiên (khi $x\to -\infty$); $\lim\limits_{x\to {-4}^{+}}y=+\infty$ nên $x=-4$ là tiệm cận đứng. Kết quả là $-8$.}]{Tiệm cận xiên của đồ thị hàm số $y=\dfrac{2x^{2}-2x-6}{x+4}$ có phương trình $y=ax+b$. Tính $a+b$.}{-8}

\cauTLN[id=C1-B3-D03-TH-074,muc=TH,loigiai={Xét hàm số $y=\dfrac{-x^{3}}{x^{2}-4}$. Chia đa thức: $y=-x+\dfrac{- 4 x}{x^{2} - 4}$; $a=\lim\limits_{x\to +\infty}\dfrac{y}{x}=-1$, $b=\lim\limits_{x\to +\infty}(y+x)=0$ nên $y=-x$ là tiệm cận xiên (khi $x\to +\infty$); $a=\lim\limits_{x\to -\infty}\dfrac{y}{x}=-1$, $b=\lim\limits_{x\to -\infty}(y+x)=0$ nên $y=-x$ là tiệm cận xiên (khi $x\to -\infty$); $\lim\limits_{x\to {-2}^{+}}y=-\infty$ nên $x=-2$ là tiệm cận đứng; $\lim\limits_{x\to {2}^{+}}y=-\infty$ nên $x=2$ là tiệm cận đứng. Kết quả là $0$.}]{Giao điểm của tiệm cận đứng và tiệm cận xiên của đồ thị hàm số $y=\dfrac{-x^{3}}{x^{2}-4}$ là $I(x_0;y_0)$. Tính $x_0+y_0$.}{0}

\cauTLN[id=C1-B3-D03-TH-075,muc=TH,loigiai={Xét hàm số $y=\dfrac{-x^{2}-3x-1}{x-1}$. Chia đa thức: $y=-x-4+\dfrac{-5}{x - 1}$; $a=\lim\limits_{x\to +\infty}\dfrac{y}{x}=-1$, $b=\lim\limits_{x\to +\infty}(y+x)=-4$ nên $y=-x-4$ là tiệm cận xiên (khi $x\to +\infty$); $a=\lim\limits_{x\to -\infty}\dfrac{y}{x}=-1$, $b=\lim\limits_{x\to -\infty}(y+x)=-4$ nên $y=-x-4$ là tiệm cận xiên (khi $x\to -\infty$); $\lim\limits_{x\to {1}^{+}}y=-\infty$ nên $x=1$ là tiệm cận đứng. Kết quả là $2$.}]{Tiệm cận xiên của đồ thị hàm số $y=\dfrac{-x^{2}-3x-1}{x-1}$ có phương trình $y=ax+b$. Tính $2a-b$.}{2}

\cauTLN[id=C1-B3-D03-TH-076,muc=TH,loigiai={Xét hàm số $y=\dfrac{3x^{2}+4x-1}{x-4}$. Chia đa thức: $y=3x+16+\dfrac{63}{x - 4}$; $a=\lim\limits_{x\to +\infty}\dfrac{y}{x}=3$, $b=\lim\limits_{x\to +\infty}(y-3x)=16$ nên $y=3x+16$ là tiệm cận xiên (khi $x\to +\infty$); $a=\lim\limits_{x\to -\infty}\dfrac{y}{x}=3$, $b=\lim\limits_{x\to -\infty}(y-3x)=16$ nên $y=3x+16$ là tiệm cận xiên (khi $x\to -\infty$); $\lim\limits_{x\to {4}^{+}}y=+\infty$ nên $x=4$ là tiệm cận đứng. Kết quả là $19$.}]{Tiệm cận xiên của đồ thị hàm số $y=\dfrac{3x^{2}+4x-1}{x-4}$ có phương trình $y=ax+b$. Tính $a+b$.}{19}

\cauTLN[id=C1-B3-D03-TH-077,muc=TH,loigiai={Xét hàm số $y=\dfrac{2x^{2}-5x-6}{x-3}$. Chia đa thức: $y=2x+1+\dfrac{-3}{x - 3}$; $a=\lim\limits_{x\to +\infty}\dfrac{y}{x}=2$, $b=\lim\limits_{x\to +\infty}(y-2x)=1$ nên $y=2x+1$ là tiệm cận xiên (khi $x\to +\infty$); $a=\lim\limits_{x\to -\infty}\dfrac{y}{x}=2$, $b=\lim\limits_{x\to -\infty}(y-2x)=1$ nên $y=2x+1$ là tiệm cận xiên (khi $x\to -\infty$); $\lim\limits_{x\to {3}^{+}}y=-\infty$ nên $x=3$ là tiệm cận đứng. Kết quả là $10$.}]{Giao điểm của tiệm cận đứng và tiệm cận xiên của đồ thị hàm số $y=\dfrac{2x^{2}-5x-6}{x-3}$ là $I(x_0;y_0)$. Tính $x_0+y_0$.}{10}

\cauTLN[id=C1-B3-D03-TH-078,muc=TH,loigiai={Xét hàm số $y=\dfrac{x^{3}+2x^{2}-3}{x^{2}-4}$. Chia đa thức: $y=x+2+\dfrac{4 x + 5}{x^{2} - 4}$; $a=\lim\limits_{x\to +\infty}\dfrac{y}{x}=1$, $b=\lim\limits_{x\to +\infty}(y-x)=2$ nên $y=x+2$ là tiệm cận xiên (khi $x\to +\infty$); $a=\lim\limits_{x\to -\infty}\dfrac{y}{x}=1$, $b=\lim\limits_{x\to -\infty}(y-x)=2$ nên $y=x+2$ là tiệm cận xiên (khi $x\to -\infty$); $\lim\limits_{x\to {-2}^{+}}y=+\infty$ nên $x=-2$ là tiệm cận đứng; $\lim\limits_{x\to {2}^{+}}y=+\infty$ nên $x=2$ là tiệm cận đứng. Kết quả là $-2$.}]{Giao điểm của tiệm cận đứng và tiệm cận xiên của đồ thị hàm số $y=\dfrac{x^{3}+2x^{2}-3}{x^{2}-4}$ là $I(x_0;y_0)$. Tính $x_0+y_0$.}{-2}

\cauTLN[id=C1-B3-D03-TH-079,muc=TH,loigiai={Xét hàm số $y=2x-2+\dfrac{1}{x+3}$. Chia đa thức: $y=2x-2+\dfrac{1}{x + 3}$; $a=\lim\limits_{x\to +\infty}\dfrac{y}{x}=2$, $b=\lim\limits_{x\to +\infty}(y-2x)=-2$ nên $y=2x-2$ là tiệm cận xiên (khi $x\to +\infty$); $a=\lim\limits_{x\to -\infty}\dfrac{y}{x}=2$, $b=\lim\limits_{x\to -\infty}(y-2x)=-2$ nên $y=2x-2$ là tiệm cận xiên (khi $x\to -\infty$); $\lim\limits_{x\to {-3}^{+}}y=+\infty$ nên $x=-3$ là tiệm cận đứng. Kết quả là $-11$.}]{Giao điểm của tiệm cận đứng và tiệm cận xiên của đồ thị hàm số $y=2x-2+\dfrac{1}{x+3}$ là $I(x_0;y_0)$. Tính $x_0+y_0$.}{-11}

\cauTLN[id=C1-B3-D03-TH-080,muc=TH,loigiai={Xét hàm số $y=-x+\dfrac{4}{x-2}$. Chia đa thức: $y=-x+\dfrac{4}{x - 2}$; $a=\lim\limits_{x\to +\infty}\dfrac{y}{x}=-1$, $b=\lim\limits_{x\to +\infty}(y+x)=0$ nên $y=-x$ là tiệm cận xiên (khi $x\to +\infty$); $a=\lim\limits_{x\to -\infty}\dfrac{y}{x}=-1$, $b=\lim\limits_{x\to -\infty}(y+x)=0$ nên $y=-x$ là tiệm cận xiên (khi $x\to -\infty$); $\lim\limits_{x\to {2}^{+}}y=+\infty$ nên $x=2$ là tiệm cận đứng. Kết quả là $0$.}]{Tiệm cận xiên của đồ thị hàm số $y=-x+\dfrac{4}{x-2}$ cắt trục tung tại điểm có tung độ bằng bao nhiêu?}{0}

\cauTN[id=C1-B3-D03-VD-001,muc=VD,loigiai={Xét hàm số $y=x-\sqrt{x^{2}-6x}$. $\lim\limits_{x\to +\infty}y=3$ nên $y=3$ là tiệm cận ngang; $a=\lim\limits_{x\to -\infty}\dfrac{y}{x}=2$, $b=\lim\limits_{x\to -\infty}(y-2x)=-3$ nên $y=2x-3$ là tiệm cận xiên (khi $x\to -\infty$).}]{Tổng số đường tiệm cận (đứng, ngang, xiên) của đồ thị hàm số $y=x-\sqrt{x^{2}-6x}$ là}
  {$2$.}
  {$3$.}{$1$.}{$4$.}

\cauTN[id=C1-B3-D03-VD-002,muc=VD,loigiai={Xét hàm số $y=\dfrac{3x^{2}+x+2}{-x+3}$. Chia đa thức: $y=-3x-10+\dfrac{32}{3 - x}$; $a=\lim\limits_{x\to +\infty}\dfrac{y}{x}=-3$, $b=\lim\limits_{x\to +\infty}(y+3x)=-10$ nên $y=-3x-10$ là tiệm cận xiên (khi $x\to +\infty$); $a=\lim\limits_{x\to -\infty}\dfrac{y}{x}=-3$, $b=\lim\limits_{x\to -\infty}(y+3x)=-10$ nên $y=-3x-10$ là tiệm cận xiên (khi $x\to -\infty$); $\lim\limits_{x\to {3}^{+}}y=-\infty$ nên $x=3$ là tiệm cận đứng.}]{Tổng số đường tiệm cận (đứng, ngang, xiên) của đồ thị hàm số $y=\dfrac{3x^{2}+x+2}{-x+3}$ là}
  {$2$.}
  {$3$.}{$1$.}{$4$.}

\cauTN[id=C1-B3-D03-VD-003,muc=VD,loigiai={Xét hàm số $y=x-\sqrt{x^{2}-2x}$. $\lim\limits_{x\to +\infty}y=1$ nên $y=1$ là tiệm cận ngang; $a=\lim\limits_{x\to -\infty}\dfrac{y}{x}=2$, $b=\lim\limits_{x\to -\infty}(y-2x)=-1$ nên $y=2x-1$ là tiệm cận xiên (khi $x\to -\infty$).}]{Tổng số đường tiệm cận (đứng, ngang, xiên) của đồ thị hàm số $y=x-\sqrt{x^{2}-2x}$ là}
  {$2$.}
  {$3$.}{$1$.}{$4$.}

\cauTN[id=C1-B3-D03-VD-004,muc=VD,loigiai={Xét hàm số $y=\dfrac{x^{3}+x^{2}}{x^{2}+1}$. Chia đa thức: $y=x+1+\dfrac{- x - 1}{x^{2} + 1}$; $a=\lim\limits_{x\to +\infty}\dfrac{y}{x}=1$, $b=\lim\limits_{x\to +\infty}(y-x)=1$ nên $y=x+1$ là tiệm cận xiên (khi $x\to +\infty$); $a=\lim\limits_{x\to -\infty}\dfrac{y}{x}=1$, $b=\lim\limits_{x\to -\infty}(y-x)=1$ nên $y=x+1$ là tiệm cận xiên (khi $x\to -\infty$).}]{Tổng số đường tiệm cận (đứng, ngang, xiên) của đồ thị hàm số $y=\dfrac{x^{3}+x^{2}}{x^{2}+1}$ là}
  {$1$.}
  {$2$.}{$0$.}{$3$.}

\cauTN[id=C1-B3-D03-VD-005,muc=VD,loigiai={Xét hàm số $y=\dfrac{x^{3}-1}{x^{2}+1}$. Chia đa thức: $y=x+\dfrac{- x - 1}{x^{2} + 1}$; $a=\lim\limits_{x\to +\infty}\dfrac{y}{x}=1$, $b=\lim\limits_{x\to +\infty}(y-x)=0$ nên $y=x$ là tiệm cận xiên (khi $x\to +\infty$); $a=\lim\limits_{x\to -\infty}\dfrac{y}{x}=1$, $b=\lim\limits_{x\to -\infty}(y-x)=0$ nên $y=x$ là tiệm cận xiên (khi $x\to -\infty$).}]{Tiệm cận xiên của đồ thị hàm số $y=\dfrac{x^{3}-1}{x^{2}+1}$ đi qua điểm nào dưới đây?}
  {$M(-2;-2)$.}
  {$M(-2;-1)$.}{$M(-2;-3)$.}{$M(-2;0)$.}

\cauTN[id=C1-B3-D03-VD-006,muc=VD,loigiai={Xét hàm số $y=x-\sqrt{x^{2}-4x}$. $\lim\limits_{x\to +\infty}y=2$ nên $y=2$ là tiệm cận ngang; $a=\lim\limits_{x\to -\infty}\dfrac{y}{x}=2$, $b=\lim\limits_{x\to -\infty}(y-2x)=-2$ nên $y=2x-2$ là tiệm cận xiên (khi $x\to -\infty$).}]{Tổng số đường tiệm cận (đứng, ngang, xiên) của đồ thị hàm số $y=x-\sqrt{x^{2}-4x}$ là}
  {$2$.}
  {$3$.}{$1$.}{$4$.}

\cauTN[id=C1-B3-D03-VD-007,muc=VD,loigiai={Xét hàm số $y=\dfrac{2x^{2}+5x-6}{x}$. Chia đa thức: $y=2x+5+\dfrac{-6}{x}$; $a=\lim\limits_{x\to +\infty}\dfrac{y}{x}=2$, $b=\lim\limits_{x\to +\infty}(y-2x)=5$ nên $y=2x+5$ là tiệm cận xiên (khi $x\to +\infty$); $a=\lim\limits_{x\to -\infty}\dfrac{y}{x}=2$, $b=\lim\limits_{x\to -\infty}(y-2x)=5$ nên $y=2x+5$ là tiệm cận xiên (khi $x\to -\infty$); $\lim\limits_{x\to {0}^{+}}y=-\infty$ nên $x=0$ là tiệm cận đứng.}]{Tổng số đường tiệm cận (đứng, ngang, xiên) của đồ thị hàm số $y=\dfrac{2x^{2}+5x-6}{x}$ là}
  {$2$.}
  {$3$.}{$1$.}{$4$.}

\cauTN[id=C1-B3-D03-VD-008,muc=VD,loigiai={Xét hàm số $y=\dfrac{-x^{3}+x^{2}-2}{x^{2}+1}$. Chia đa thức: $y=-x+1+\dfrac{x - 3}{x^{2} + 1}$; $a=\lim\limits_{x\to +\infty}\dfrac{y}{x}=-1$, $b=\lim\limits_{x\to +\infty}(y+x)=1$ nên $y=-x+1$ là tiệm cận xiên (khi $x\to +\infty$); $a=\lim\limits_{x\to -\infty}\dfrac{y}{x}=-1$, $b=\lim\limits_{x\to -\infty}(y+x)=1$ nên $y=-x+1$ là tiệm cận xiên (khi $x\to -\infty$).}]{Tổng số đường tiệm cận (đứng, ngang, xiên) của đồ thị hàm số $y=\dfrac{-x^{3}+x^{2}-2}{x^{2}+1}$ là}
  {$1$.}
  {$2$.}{$0$.}{$3$.}

\cauTN[id=C1-B3-D03-VD-009,muc=VD,loigiai={Xét hàm số $y=\dfrac{-x^{2}+5}{x+3}$. Chia đa thức: $y=-x+3+\dfrac{-4}{x + 3}$; $a=\lim\limits_{x\to +\infty}\dfrac{y}{x}=-1$, $b=\lim\limits_{x\to +\infty}(y+x)=3$ nên $y=-x+3$ là tiệm cận xiên (khi $x\to +\infty$); $a=\lim\limits_{x\to -\infty}\dfrac{y}{x}=-1$, $b=\lim\limits_{x\to -\infty}(y+x)=3$ nên $y=-x+3$ là tiệm cận xiên (khi $x\to -\infty$); $\lim\limits_{x\to {-3}^{+}}y=-\infty$ nên $x=-3$ là tiệm cận đứng.}]{Tiệm cận xiên của đồ thị hàm số $y=\dfrac{-x^{2}+5}{x+3}$ đi qua điểm nào dưới đây?}
  {$M(2;1)$.}
  {$M(2;2)$.}{$M(2;0)$.}{$M(2;3)$.}

\cauTN[id=C1-B3-D03-VD-010,muc=VD,loigiai={Xét hàm số $y=\dfrac{-x^{2}-4x-1}{2x}$. Chia đa thức: $y=-\dfrac{1}{2}x-2+\dfrac{-1}{2 x}$; $a=\lim\limits_{x\to +\infty}\dfrac{y}{x}=-\dfrac{1}{2}$, $b=\lim\limits_{x\to +\infty}(y+\dfrac{1}{2}x)=-2$ nên $y=-\dfrac{1}{2}x-2$ là tiệm cận xiên (khi $x\to +\infty$); $a=\lim\limits_{x\to -\infty}\dfrac{y}{x}=-\dfrac{1}{2}$, $b=\lim\limits_{x\to -\infty}(y+\dfrac{1}{2}x)=-2$ nên $y=-\dfrac{1}{2}x-2$ là tiệm cận xiên (khi $x\to -\infty$); $\lim\limits_{x\to {0}^{+}}y=-\infty$ nên $x=0$ là tiệm cận đứng.}]{Tiệm cận xiên của đồ thị hàm số $y=\dfrac{-x^{2}-4x-1}{2x}$ đi qua điểm nào dưới đây?}
  {$M(3;-\dfrac{7}{2})$.}
  {$M(3;-\dfrac{5}{2})$.}{$M(3;-\dfrac{9}{2})$.}{$M(3;-\dfrac{3}{2})$.}

\cauTN[id=C1-B3-D03-VD-011,muc=VD,loigiai={Xét hàm số $y=\sqrt{x^{2}+6x+10}$. $a=\lim\limits_{x\to +\infty}\dfrac{y}{x}=1$, $b=\lim\limits_{x\to +\infty}(y-x)=3$ nên $y=x+3$ là tiệm cận xiên (khi $x\to +\infty$); $a=\lim\limits_{x\to -\infty}\dfrac{y}{x}=-1$, $b=\lim\limits_{x\to -\infty}(y+x)=-3$ nên $y=-x-3$ là tiệm cận xiên (khi $x\to -\infty$).}]{Tổng số đường tiệm cận (đứng, ngang, xiên) của đồ thị hàm số $y=\sqrt{x^{2}+6x+10}$ là}
  {$2$.}
  {$3$.}{$1$.}{$4$.}

\cauTN[id=C1-B3-D03-VD-012,muc=VD,loigiai={Xét hàm số $y=\sqrt{x^{2}+4x+7}$. $a=\lim\limits_{x\to +\infty}\dfrac{y}{x}=1$, $b=\lim\limits_{x\to +\infty}(y-x)=2$ nên $y=x+2$ là tiệm cận xiên (khi $x\to +\infty$); $a=\lim\limits_{x\to -\infty}\dfrac{y}{x}=-1$, $b=\lim\limits_{x\to -\infty}(y+x)=-2$ nên $y=-x-2$ là tiệm cận xiên (khi $x\to -\infty$).}]{Tổng số đường tiệm cận (đứng, ngang, xiên) của đồ thị hàm số $y=\sqrt{x^{2}+4x+7}$ là}
  {$2$.}
  {$3$.}{$1$.}{$4$.}

\cauTN[id=C1-B3-D03-VD-013,muc=VD,loigiai={Xét hàm số $y=2x+\sqrt{x^2+9}$. $a=\lim\limits_{x\to +\infty}\dfrac{y}{x}=3$, $b=\lim\limits_{x\to +\infty}(y-3x)=0$ nên $y=3x$ là tiệm cận xiên (khi $x\to +\infty$); $a=\lim\limits_{x\to -\infty}\dfrac{y}{x}=1$, $b=\lim\limits_{x\to -\infty}(y-x)=0$ nên $y=x$ là tiệm cận xiên (khi $x\to -\infty$).}]{Tổng số đường tiệm cận (đứng, ngang, xiên) của đồ thị hàm số $y=2x+\sqrt{x^2+9}$ là}
  {$2$.}
  {$3$.}{$1$.}{$4$.}

\cauTN[id=C1-B3-D03-VD-014,muc=VD,loigiai={Xét hàm số $y=\dfrac{-x^{3}-x^{2}+1}{x^{2}-9}$. Chia đa thức: $y=-x-1+\dfrac{- 9 x - 8}{x^{2} - 9}$; $a=\lim\limits_{x\to +\infty}\dfrac{y}{x}=-1$, $b=\lim\limits_{x\to +\infty}(y+x)=-1$ nên $y=-x-1$ là tiệm cận xiên (khi $x\to +\infty$); $a=\lim\limits_{x\to -\infty}\dfrac{y}{x}=-1$, $b=\lim\limits_{x\to -\infty}(y+x)=-1$ nên $y=-x-1$ là tiệm cận xiên (khi $x\to -\infty$); $\lim\limits_{x\to {-3}^{+}}y=-\infty$ nên $x=-3$ là tiệm cận đứng; $\lim\limits_{x\to {3}^{+}}y=-\infty$ nên $x=3$ là tiệm cận đứng.}]{Tiệm cận xiên của đồ thị hàm số $y=\dfrac{-x^{3}-x^{2}+1}{x^{2}-9}$ đi qua điểm nào dưới đây?}
  {$M(2;-3)$.}
  {$M(2;-2)$.}{$M(2;-4)$.}{$M(2;-1)$.}

\cauTN[id=C1-B3-D03-VD-015,muc=VD,loigiai={Xét hàm số $y=\dfrac{2x^{2}-2x-6}{x-2}$. Chia đa thức: $y=2x+2+\dfrac{-2}{x - 2}$; $a=\lim\limits_{x\to +\infty}\dfrac{y}{x}=2$, $b=\lim\limits_{x\to +\infty}(y-2x)=2$ nên $y=2x+2$ là tiệm cận xiên (khi $x\to +\infty$); $a=\lim\limits_{x\to -\infty}\dfrac{y}{x}=2$, $b=\lim\limits_{x\to -\infty}(y-2x)=2$ nên $y=2x+2$ là tiệm cận xiên (khi $x\to -\infty$); $\lim\limits_{x\to {2}^{+}}y=-\infty$ nên $x=2$ là tiệm cận đứng.}]{Tổng số đường tiệm cận (đứng, ngang, xiên) của đồ thị hàm số $y=\dfrac{2x^{2}-2x-6}{x-2}$ là}
  {$2$.}
  {$3$.}{$1$.}{$4$.}

\cauTN[id=C1-B3-D03-VD-016,muc=VD,loigiai={Xét hàm số $y=\dfrac{-2x^{2}+4x-2}{-x}$. Chia đa thức: $y=2x-4+\dfrac{2}{x}$; $a=\lim\limits_{x\to +\infty}\dfrac{y}{x}=2$, $b=\lim\limits_{x\to +\infty}(y-2x)=-4$ nên $y=2x-4$ là tiệm cận xiên (khi $x\to +\infty$); $a=\lim\limits_{x\to -\infty}\dfrac{y}{x}=2$, $b=\lim\limits_{x\to -\infty}(y-2x)=-4$ nên $y=2x-4$ là tiệm cận xiên (khi $x\to -\infty$); $\lim\limits_{x\to {0}^{+}}y=+\infty$ nên $x=0$ là tiệm cận đứng.}]{Tiệm cận xiên của đồ thị hàm số $y=\dfrac{-2x^{2}+4x-2}{-x}$ là đường thẳng}
  {$y=2x-4$.}
  {$y=2x+4$.}{$y=-2x-4$.}{$y=2x-3$.}

\cauTN[id=C1-B3-D03-VD-017,muc=VD,loigiai={Xét hàm số $y=x-\sqrt{x^{2}-2x}$. $\lim\limits_{x\to +\infty}y=1$ nên $y=1$ là tiệm cận ngang; $a=\lim\limits_{x\to -\infty}\dfrac{y}{x}=2$, $b=\lim\limits_{x\to -\infty}(y-2x)=-1$ nên $y=2x-1$ là tiệm cận xiên (khi $x\to -\infty$).}]{Tiệm cận xiên của đồ thị hàm số $y=x-\sqrt{x^{2}-2x}$ đi qua điểm nào dưới đây?}
  {$M(2;3)$.}
  {$M(2;4)$.}{$M(2;2)$.}{$M(2;5)$.}

\cauTN[id=C1-B3-D03-VD-018,muc=VD,loigiai={Xét hàm số $y=x-\sqrt{x^{2}-6x}$. $\lim\limits_{x\to +\infty}y=3$ nên $y=3$ là tiệm cận ngang; $a=\lim\limits_{x\to -\infty}\dfrac{y}{x}=2$, $b=\lim\limits_{x\to -\infty}(y-2x)=-3$ nên $y=2x-3$ là tiệm cận xiên (khi $x\to -\infty$).}]{Tiệm cận xiên của đồ thị hàm số $y=x-\sqrt{x^{2}-6x}$ là đường thẳng}
  {$y=2x-3$.}
  {$y=2x+3$.}{$y=-2x-3$.}{$y=2x-2$.}

\cauTN[id=C1-B3-D03-VD-019,muc=VD,loigiai={Xét hàm số $y=\dfrac{2x^{2}-5x+1}{x-1}$. Chia đa thức: $y=2x-3+\dfrac{-2}{x - 1}$; $a=\lim\limits_{x\to +\infty}\dfrac{y}{x}=2$, $b=\lim\limits_{x\to +\infty}(y-2x)=-3$ nên $y=2x-3$ là tiệm cận xiên (khi $x\to +\infty$); $a=\lim\limits_{x\to -\infty}\dfrac{y}{x}=2$, $b=\lim\limits_{x\to -\infty}(y-2x)=-3$ nên $y=2x-3$ là tiệm cận xiên (khi $x\to -\infty$); $\lim\limits_{x\to {1}^{+}}y=-\infty$ nên $x=1$ là tiệm cận đứng.}]{Tiệm cận xiên của đồ thị hàm số $y=\dfrac{2x^{2}-5x+1}{x-1}$ là đường thẳng}
  {$y=2x-3$.}
  {$y=2x+3$.}{$y=-2x-3$.}{$y=2x-2$.}

\cauTN[id=C1-B3-D03-VD-020,muc=VD,loigiai={Xét hàm số $y=\sqrt{x^{2}-2x+1}$. $a=\lim\limits_{x\to +\infty}\dfrac{y}{x}=1$, $b=\lim\limits_{x\to +\infty}(y-x)=-1$ nên $y=x-1$ là tiệm cận xiên (khi $x\to +\infty$); $a=\lim\limits_{x\to -\infty}\dfrac{y}{x}=-1$, $b=\lim\limits_{x\to -\infty}(y+x)=1$ nên $y=-x+1$ là tiệm cận xiên (khi $x\to -\infty$).}]{Một tiệm cận xiên của đồ thị hàm số $y=\sqrt{x^{2}-2x+1}$ đi qua điểm nào dưới đây?}
  {$M(3;2)$.}
  {$M(3;3)$.}{$M(3;1)$.}{$M(3;4)$.}

\cauTN[id=C1-B3-D03-VD-021,muc=VD,loigiai={Xét hàm số $y=\sqrt{x^{2}-2x+4}$. $a=\lim\limits_{x\to +\infty}\dfrac{y}{x}=1$, $b=\lim\limits_{x\to +\infty}(y-x)=-1$ nên $y=x-1$ là tiệm cận xiên (khi $x\to +\infty$); $a=\lim\limits_{x\to -\infty}\dfrac{y}{x}=-1$, $b=\lim\limits_{x\to -\infty}(y+x)=1$ nên $y=-x+1$ là tiệm cận xiên (khi $x\to -\infty$).}]{Tổng số đường tiệm cận (đứng, ngang, xiên) của đồ thị hàm số $y=\sqrt{x^{2}-2x+4}$ là}
  {$2$.}
  {$3$.}{$1$.}{$4$.}

\cauTN[id=C1-B3-D03-VD-022,muc=VD,loigiai={Xét hàm số $y=\sqrt{x^{2}-4x+4}$. $a=\lim\limits_{x\to +\infty}\dfrac{y}{x}=1$, $b=\lim\limits_{x\to +\infty}(y-x)=-2$ nên $y=x-2$ là tiệm cận xiên (khi $x\to +\infty$); $a=\lim\limits_{x\to -\infty}\dfrac{y}{x}=-1$, $b=\lim\limits_{x\to -\infty}(y+x)=2$ nên $y=-x+2$ là tiệm cận xiên (khi $x\to -\infty$).}]{Tổng số đường tiệm cận (đứng, ngang, xiên) của đồ thị hàm số $y=\sqrt{x^{2}-4x+4}$ là}
  {$2$.}
  {$3$.}{$1$.}{$4$.}

\cauTN[id=C1-B3-D03-VD-023,muc=VD,loigiai={Xét hàm số $y=\dfrac{-x^{2}+5x+5}{-x-3}$. Chia đa thức: $y=x-8+\dfrac{-19}{- x - 3}$; $a=\lim\limits_{x\to +\infty}\dfrac{y}{x}=1$, $b=\lim\limits_{x\to +\infty}(y-x)=-8$ nên $y=x-8$ là tiệm cận xiên (khi $x\to +\infty$); $a=\lim\limits_{x\to -\infty}\dfrac{y}{x}=1$, $b=\lim\limits_{x\to -\infty}(y-x)=-8$ nên $y=x-8$ là tiệm cận xiên (khi $x\to -\infty$); $\lim\limits_{x\to {-3}^{+}}y=+\infty$ nên $x=-3$ là tiệm cận đứng.}]{Tiệm cận xiên của đồ thị hàm số $y=\dfrac{-x^{2}+5x+5}{-x-3}$ là đường thẳng}
  {$y=x-8$.}
  {$y=x+8$.}{$y=-x-8$.}{$y=x-7$.}

\cauTN[id=C1-B3-D03-VD-024,muc=VD,loigiai={Xét hàm số $y=\sqrt{x^{2}+6x+12}$. $a=\lim\limits_{x\to +\infty}\dfrac{y}{x}=1$, $b=\lim\limits_{x\to +\infty}(y-x)=3$ nên $y=x+3$ là tiệm cận xiên (khi $x\to +\infty$); $a=\lim\limits_{x\to -\infty}\dfrac{y}{x}=-1$, $b=\lim\limits_{x\to -\infty}(y+x)=-3$ nên $y=-x-3$ là tiệm cận xiên (khi $x\to -\infty$).}]{Tổng số đường tiệm cận (đứng, ngang, xiên) của đồ thị hàm số $y=\sqrt{x^{2}+6x+12}$ là}
  {$2$.}
  {$3$.}{$1$.}{$4$.}

\cauTN[id=C1-B3-D03-VD-025,muc=VD,loigiai={Xét hàm số $y=x-\sqrt{x^{2}-6x}$. $\lim\limits_{x\to +\infty}y=3$ nên $y=3$ là tiệm cận ngang; $a=\lim\limits_{x\to -\infty}\dfrac{y}{x}=2$, $b=\lim\limits_{x\to -\infty}(y-2x)=-3$ nên $y=2x-3$ là tiệm cận xiên (khi $x\to -\infty$).}]{Tiệm cận xiên của đồ thị hàm số $y=x-\sqrt{x^{2}-6x}$ đi qua điểm nào dưới đây?}
  {$M(-2;-7)$.}
  {$M(-2;-6)$.}{$M(-2;-8)$.}{$M(-2;-5)$.}

\cauTN[id=C1-B3-D03-VD-026,muc=VD,loigiai={Xét hàm số $y=\dfrac{x^{3}}{x^{2}-4}$. Chia đa thức: $y=x+\dfrac{4 x}{x^{2} - 4}$; $a=\lim\limits_{x\to +\infty}\dfrac{y}{x}=1$, $b=\lim\limits_{x\to +\infty}(y-x)=0$ nên $y=x$ là tiệm cận xiên (khi $x\to +\infty$); $a=\lim\limits_{x\to -\infty}\dfrac{y}{x}=1$, $b=\lim\limits_{x\to -\infty}(y-x)=0$ nên $y=x$ là tiệm cận xiên (khi $x\to -\infty$); $\lim\limits_{x\to {-2}^{+}}y=+\infty$ nên $x=-2$ là tiệm cận đứng; $\lim\limits_{x\to {2}^{+}}y=+\infty$ nên $x=2$ là tiệm cận đứng.}]{Tổng số đường tiệm cận (đứng, ngang, xiên) của đồ thị hàm số $y=\dfrac{x^{3}}{x^{2}-4}$ là}
  {$3$.}
  {$4$.}{$2$.}{$5$.}

\cauTN[id=C1-B3-D03-VD-027,muc=VD,loigiai={Xét hàm số $y=\dfrac{x^{2}+2x-6}{-x+3}$. Chia đa thức: $y=-x-5+\dfrac{9}{3 - x}$; $a=\lim\limits_{x\to +\infty}\dfrac{y}{x}=-1$, $b=\lim\limits_{x\to +\infty}(y+x)=-5$ nên $y=-x-5$ là tiệm cận xiên (khi $x\to +\infty$); $a=\lim\limits_{x\to -\infty}\dfrac{y}{x}=-1$, $b=\lim\limits_{x\to -\infty}(y+x)=-5$ nên $y=-x-5$ là tiệm cận xiên (khi $x\to -\infty$); $\lim\limits_{x\to {3}^{+}}y=-\infty$ nên $x=3$ là tiệm cận đứng.}]{Tiệm cận xiên của đồ thị hàm số $y=\dfrac{x^{2}+2x-6}{-x+3}$ đi qua điểm nào dưới đây?}
  {$M(2;-7)$.}
  {$M(2;-6)$.}{$M(2;-8)$.}{$M(2;-5)$.}

\cauTN[id=C1-B3-D03-VD-028,muc=VD,loigiai={Xét hàm số $y=\sqrt{x^{2}-4x+7}$. $a=\lim\limits_{x\to +\infty}\dfrac{y}{x}=1$, $b=\lim\limits_{x\to +\infty}(y-x)=-2$ nên $y=x-2$ là tiệm cận xiên (khi $x\to +\infty$); $a=\lim\limits_{x\to -\infty}\dfrac{y}{x}=-1$, $b=\lim\limits_{x\to -\infty}(y+x)=2$ nên $y=-x+2$ là tiệm cận xiên (khi $x\to -\infty$).}]{Tổng số đường tiệm cận (đứng, ngang, xiên) của đồ thị hàm số $y=\sqrt{x^{2}-4x+7}$ là}
  {$2$.}
  {$3$.}{$1$.}{$4$.}

\cauTN[id=C1-B3-D03-VD-029,muc=VD,loigiai={Xét hàm số $y=2x+\sqrt{x^2+4}$. $a=\lim\limits_{x\to +\infty}\dfrac{y}{x}=3$, $b=\lim\limits_{x\to +\infty}(y-3x)=0$ nên $y=3x$ là tiệm cận xiên (khi $x\to +\infty$); $a=\lim\limits_{x\to -\infty}\dfrac{y}{x}=1$, $b=\lim\limits_{x\to -\infty}(y-x)=0$ nên $y=x$ là tiệm cận xiên (khi $x\to -\infty$).}]{Tổng số đường tiệm cận (đứng, ngang, xiên) của đồ thị hàm số $y=2x+\sqrt{x^2+4}$ là}
  {$2$.}
  {$3$.}{$1$.}{$4$.}

\cauTN[id=C1-B3-D03-VD-030,muc=VD,loigiai={Xét hàm số $y=\sqrt{x^{2}+4x+4}$. $a=\lim\limits_{x\to +\infty}\dfrac{y}{x}=1$, $b=\lim\limits_{x\to +\infty}(y-x)=2$ nên $y=x+2$ là tiệm cận xiên (khi $x\to +\infty$); $a=\lim\limits_{x\to -\infty}\dfrac{y}{x}=-1$, $b=\lim\limits_{x\to -\infty}(y+x)=-2$ nên $y=-x-2$ là tiệm cận xiên (khi $x\to -\infty$).}]{Tổng số đường tiệm cận (đứng, ngang, xiên) của đồ thị hàm số $y=\sqrt{x^{2}+4x+4}$ là}
  {$2$.}
  {$3$.}{$1$.}{$4$.}

\cauTN[id=C1-B3-D03-VD-031,muc=VD,loigiai={Xét hàm số $y=\dfrac{-x^{3}-x^{2}+1}{x^{2}-4}$. Chia đa thức: $y=-x-1+\dfrac{- 4 x - 3}{x^{2} - 4}$; $a=\lim\limits_{x\to +\infty}\dfrac{y}{x}=-1$, $b=\lim\limits_{x\to +\infty}(y+x)=-1$ nên $y=-x-1$ là tiệm cận xiên (khi $x\to +\infty$); $a=\lim\limits_{x\to -\infty}\dfrac{y}{x}=-1$, $b=\lim\limits_{x\to -\infty}(y+x)=-1$ nên $y=-x-1$ là tiệm cận xiên (khi $x\to -\infty$); $\lim\limits_{x\to {-2}^{+}}y=-\infty$ nên $x=-2$ là tiệm cận đứng; $\lim\limits_{x\to {2}^{+}}y=-\infty$ nên $x=2$ là tiệm cận đứng.}]{Tiệm cận xiên của đồ thị hàm số $y=\dfrac{-x^{3}-x^{2}+1}{x^{2}-4}$ đi qua điểm nào dưới đây?}
  {$M(2;-3)$.}
  {$M(2;-2)$.}{$M(2;-4)$.}{$M(2;-1)$.}

\cauTN[id=C1-B3-D03-VD-032,muc=VD,loigiai={Xét hàm số $y=\dfrac{-x^{3}+x^{2}-2}{x^{2}-4}$. Chia đa thức: $y=-x+1+\dfrac{2 - 4 x}{x^{2} - 4}$; $a=\lim\limits_{x\to +\infty}\dfrac{y}{x}=-1$, $b=\lim\limits_{x\to +\infty}(y+x)=1$ nên $y=-x+1$ là tiệm cận xiên (khi $x\to +\infty$); $a=\lim\limits_{x\to -\infty}\dfrac{y}{x}=-1$, $b=\lim\limits_{x\to -\infty}(y+x)=1$ nên $y=-x+1$ là tiệm cận xiên (khi $x\to -\infty$); $\lim\limits_{x\to {-2}^{+}}y=-\infty$ nên $x=-2$ là tiệm cận đứng; $\lim\limits_{x\to {2}^{+}}y=-\infty$ nên $x=2$ là tiệm cận đứng.}]{Tiệm cận xiên của đồ thị hàm số $y=\dfrac{-x^{3}+x^{2}-2}{x^{2}-4}$ là đường thẳng}
  {$y=-x+1$.}
  {$y=-x-1$.}{$y=x+1$.}{$y=-x+2$.}

\cauTN[id=C1-B3-D03-VD-033,muc=VD,loigiai={Xét hàm số $y=\dfrac{2x^{2}-3x+5}{2x-4}$. Chia đa thức: $y=x+\dfrac{1}{2}+\dfrac{7}{2 x - 4}$; $a=\lim\limits_{x\to +\infty}\dfrac{y}{x}=1$, $b=\lim\limits_{x\to +\infty}(y-x)=\dfrac{1}{2}$ nên $y=x+\dfrac{1}{2}$ là tiệm cận xiên (khi $x\to +\infty$); $a=\lim\limits_{x\to -\infty}\dfrac{y}{x}=1$, $b=\lim\limits_{x\to -\infty}(y-x)=\dfrac{1}{2}$ nên $y=x+\dfrac{1}{2}$ là tiệm cận xiên (khi $x\to -\infty$); $\lim\limits_{x\to {2}^{+}}y=+\infty$ nên $x=2$ là tiệm cận đứng.}]{Tiệm cận xiên của đồ thị hàm số $y=\dfrac{2x^{2}-3x+5}{2x-4}$ đi qua điểm nào dưới đây?}
  {$M(-2;-\dfrac{3}{2})$.}
  {$M(-2;-\dfrac{1}{2})$.}{$M(-2;-\dfrac{5}{2})$.}{$M(-2;\dfrac{1}{2})$.}

\cauTN[id=C1-B3-D03-VD-034,muc=VD,loigiai={Xét hàm số $y=\dfrac{x^{3}-2x^{2}-1}{x^{2}-1}$. Chia đa thức: $y=x-2+\dfrac{x - 3}{x^{2} - 1}$; $a=\lim\limits_{x\to +\infty}\dfrac{y}{x}=1$, $b=\lim\limits_{x\to +\infty}(y-x)=-2$ nên $y=x-2$ là tiệm cận xiên (khi $x\to +\infty$); $a=\lim\limits_{x\to -\infty}\dfrac{y}{x}=1$, $b=\lim\limits_{x\to -\infty}(y-x)=-2$ nên $y=x-2$ là tiệm cận xiên (khi $x\to -\infty$); $\lim\limits_{x\to {-1}^{+}}y=+\infty$ nên $x=-1$ là tiệm cận đứng; $\lim\limits_{x\to {1}^{+}}y=-\infty$ nên $x=1$ là tiệm cận đứng.}]{Tiệm cận xiên của đồ thị hàm số $y=\dfrac{x^{3}-2x^{2}-1}{x^{2}-1}$ đi qua điểm nào dưới đây?}
  {$M(3;1)$.}
  {$M(3;2)$.}{$M(3;0)$.}{$M(3;3)$.}

\cauTN[id=C1-B3-D03-VD-035,muc=VD,loigiai={Xét hàm số $y=3x+\sqrt{x^2+1}$. $a=\lim\limits_{x\to +\infty}\dfrac{y}{x}=4$, $b=\lim\limits_{x\to +\infty}(y-4x)=0$ nên $y=4x$ là tiệm cận xiên (khi $x\to +\infty$); $a=\lim\limits_{x\to -\infty}\dfrac{y}{x}=2$, $b=\lim\limits_{x\to -\infty}(y-2x)=0$ nên $y=2x$ là tiệm cận xiên (khi $x\to -\infty$).}]{Tổng số đường tiệm cận (đứng, ngang, xiên) của đồ thị hàm số $y=3x+\sqrt{x^2+1}$ là}
  {$2$.}
  {$3$.}{$1$.}{$4$.}

\cauTN[id=C1-B3-D03-VD-036,muc=VD,loigiai={Xét hàm số $y=\sqrt{x^{2}+6x+9}$. $a=\lim\limits_{x\to +\infty}\dfrac{y}{x}=1$, $b=\lim\limits_{x\to +\infty}(y-x)=3$ nên $y=x+3$ là tiệm cận xiên (khi $x\to +\infty$); $a=\lim\limits_{x\to -\infty}\dfrac{y}{x}=-1$, $b=\lim\limits_{x\to -\infty}(y+x)=-3$ nên $y=-x-3$ là tiệm cận xiên (khi $x\to -\infty$).}]{Tổng số đường tiệm cận (đứng, ngang, xiên) của đồ thị hàm số $y=\sqrt{x^{2}+6x+9}$ là}
  {$2$.}
  {$3$.}{$1$.}{$4$.}

\cauDS[id=C1-B3-D03-VD-037,muc=VD,loigiai={Xét hàm số $y=\dfrac{x^{2}-4x-1}{x}$. Chia đa thức: $y=x-4+\dfrac{-1}{x}$; $a=\lim\limits_{x\to +\infty}\dfrac{y}{x}=1$, $b=\lim\limits_{x\to +\infty}(y-x)=-4$ nên $y=x-4$ là tiệm cận xiên (khi $x\to +\infty$); $a=\lim\limits_{x\to -\infty}\dfrac{y}{x}=1$, $b=\lim\limits_{x\to -\infty}(y-x)=-4$ nên $y=x-4$ là tiệm cận xiên (khi $x\to -\infty$); $\lim\limits_{x\to {0}^{+}}y=-\infty$ nên $x=0$ là tiệm cận đứng. a) Đúng. b) Sai. c) Đúng. d) Đúng.}]{Cho hàm số $y=\dfrac{x^{2}-4x-1}{x}$. Xét tính đúng sai của các mệnh đề sau:}{%
  \y{Giao điểm của tiệm cận đứng và tiệm cận xiên là $I(0;-4)$.}{d}
  \y{Đường thẳng $y=x-2$ là tiệm cận xiên của đồ thị hàm số.}{s}
  \y{Đồ thị hàm số có tiệm cận đứng $x=0$.}{d}
  \y{Đồ thị hàm số có tất cả $2$ đường tiệm cận.}{d}}

\cauDS[id=C1-B3-D03-VD-038,muc=VD,loigiai={Xét hàm số $y=\sqrt{x^{2}+6x+14}$. $a=\lim\limits_{x\to +\infty}\dfrac{y}{x}=1$, $b=\lim\limits_{x\to +\infty}(y-x)=3$ nên $y=x+3$ là tiệm cận xiên (khi $x\to +\infty$); $a=\lim\limits_{x\to -\infty}\dfrac{y}{x}=-1$, $b=\lim\limits_{x\to -\infty}(y+x)=-3$ nên $y=-x-3$ là tiệm cận xiên (khi $x\to -\infty$). a) Sai. b) Đúng. c) Đúng. d) Sai: có $2$ đường tiệm cận.}]{Cho hàm số $y=\sqrt{x^{2}+6x+14}$. Xét tính đúng sai của các mệnh đề sau:}{%
  \y{Đường thẳng $y=-x-1$ là tiệm cận xiên của đồ thị hàm số.}{s}
  \y{Đồ thị hàm số không có tiệm cận đứng.}{d}
  \y{Đường thẳng $y=x+3$ là tiệm cận xiên của đồ thị hàm số.}{d}
  \y{Đồ thị hàm số có tất cả $3$ đường tiệm cận.}{s}}

\cauDS[id=C1-B3-D03-VD-039,muc=VD,loigiai={Xét hàm số $y=\dfrac{x^{2}-x+3}{-x-2}$. Chia đa thức: $y=-x+3+\dfrac{9}{- x - 2}$; $a=\lim\limits_{x\to +\infty}\dfrac{y}{x}=-1$, $b=\lim\limits_{x\to +\infty}(y+x)=3$ nên $y=-x+3$ là tiệm cận xiên (khi $x\to +\infty$); $a=\lim\limits_{x\to -\infty}\dfrac{y}{x}=-1$, $b=\lim\limits_{x\to -\infty}(y+x)=3$ nên $y=-x+3$ là tiệm cận xiên (khi $x\to -\infty$); $\lim\limits_{x\to {-2}^{+}}y=-\infty$ nên $x=-2$ là tiệm cận đứng. a) Đúng. b) Sai. c) Đúng. d) Đúng.}]{Cho hàm số $y=\dfrac{x^{2}-x+3}{-x-2}$. Xét tính đúng sai của các mệnh đề sau:}{%
  \y{Đồ thị hàm số có tất cả $2$ đường tiệm cận.}{d}
  \y{Đường thẳng $y=-x+5$ là tiệm cận xiên của đồ thị hàm số.}{s}
  \y{Đường thẳng $y=-x+3$ là tiệm cận xiên của đồ thị hàm số.}{d}
  \y{Giao điểm của tiệm cận đứng và tiệm cận xiên là $I(-2;5)$.}{d}}

\cauDS[id=C1-B3-D03-VD-040,muc=VD,loigiai={Xét hàm số $y=x-\sqrt{x^{2}-4x}$. $\lim\limits_{x\to +\infty}y=2$ nên $y=2$ là tiệm cận ngang; $a=\lim\limits_{x\to -\infty}\dfrac{y}{x}=2$, $b=\lim\limits_{x\to -\infty}(y-2x)=-2$ nên $y=2x-2$ là tiệm cận xiên (khi $x\to -\infty$). a) Sai. b) Sai: có $2$ đường tiệm cận. c) Đúng. d) Đúng.}]{Cho hàm số $y=x-\sqrt{x^{2}-4x}$. Xét tính đúng sai của các mệnh đề sau:}{%
  \y{Đường thẳng $y=2x$ là tiệm cận xiên của đồ thị hàm số.}{s}
  \y{Đồ thị hàm số có tất cả $3$ đường tiệm cận.}{s}
  \y{Đường thẳng $y=2x-2$ là tiệm cận xiên của đồ thị hàm số.}{d}
  \y{Đồ thị hàm số có tiệm cận ngang $y=2$.}{d}}

\cauDS[id=C1-B3-D03-VD-041,muc=VD,loigiai={Xét hàm số $y=\sqrt{x^{2}-4x+5}$. $a=\lim\limits_{x\to +\infty}\dfrac{y}{x}=1$, $b=\lim\limits_{x\to +\infty}(y-x)=-2$ nên $y=x-2$ là tiệm cận xiên (khi $x\to +\infty$); $a=\lim\limits_{x\to -\infty}\dfrac{y}{x}=-1$, $b=\lim\limits_{x\to -\infty}(y+x)=2$ nên $y=-x+2$ là tiệm cận xiên (khi $x\to -\infty$). a) Sai. b) Đúng. c) Sai: tiệm cận xiên là $y=x-2$. d) Đúng.}]{Cho hàm số $y=\sqrt{x^{2}-4x+5}$. Xét tính đúng sai của các mệnh đề sau:}{%
  \y{Đường thẳng $y=-x+3$ là tiệm cận xiên của đồ thị hàm số.}{s}
  \y{Đồ thị hàm số có tất cả $2$ đường tiệm cận.}{d}
  \y{Đường thẳng $y=x$ là tiệm cận xiên của đồ thị hàm số.}{s}
  \y{Đồ thị hàm số không có tiệm cận đứng.}{d}}

\cauDS[id=C1-B3-D03-VD-042,muc=VD,loigiai={Xét hàm số $y=-2x+\sqrt{x^2+9}$. $a=\lim\limits_{x\to +\infty}\dfrac{y}{x}=-1$, $b=\lim\limits_{x\to +\infty}(y+x)=0$ nên $y=-x$ là tiệm cận xiên (khi $x\to +\infty$); $a=\lim\limits_{x\to -\infty}\dfrac{y}{x}=-3$, $b=\lim\limits_{x\to -\infty}(y+3x)=0$ nên $y=-3x$ là tiệm cận xiên (khi $x\to -\infty$). a) Sai. b) Sai: giới hạn của $y$ tại vô cực không hữu hạn. c) Đúng. d) Đúng.}]{Cho hàm số $y=-2x+\sqrt{x^2+9}$. Xét tính đúng sai của các mệnh đề sau:}{%
  \y{Đường thẳng $y=-x+1$ là tiệm cận xiên của đồ thị hàm số.}{s}
  \y{Đồ thị hàm số có tiệm cận ngang.}{s}
  \y{Đồ thị hàm số không có tiệm cận đứng.}{d}
  \y{Đồ thị hàm số có tất cả $2$ đường tiệm cận.}{d}}

\cauDS[id=C1-B3-D03-VD-043,muc=VD,loigiai={Xét hàm số $y=\dfrac{-x^{2}+3x+1}{2x-4}$. Chia đa thức: $y=-\dfrac{1}{2}x+\dfrac{1}{2}+\dfrac{3}{2 x - 4}$; $a=\lim\limits_{x\to +\infty}\dfrac{y}{x}=-\dfrac{1}{2}$, $b=\lim\limits_{x\to +\infty}(y+\dfrac{1}{2}x)=\dfrac{1}{2}$ nên $y=-\dfrac{1}{2}x+\dfrac{1}{2}$ là tiệm cận xiên (khi $x\to +\infty$); $a=\lim\limits_{x\to -\infty}\dfrac{y}{x}=-\dfrac{1}{2}$, $b=\lim\limits_{x\to -\infty}(y+\dfrac{1}{2}x)=\dfrac{1}{2}$ nên $y=-\dfrac{1}{2}x+\dfrac{1}{2}$ là tiệm cận xiên (khi $x\to -\infty$); $\lim\limits_{x\to {2}^{+}}y=+\infty$ nên $x=2$ là tiệm cận đứng. a) Sai: giới hạn của $y$ tại vô cực không hữu hạn. b) Đúng. c) Đúng. d) Sai.}]{Cho hàm số $y=\dfrac{-x^{2}+3x+1}{2x-4}$. Xét tính đúng sai của các mệnh đề sau:}{%
  \y{Đồ thị hàm số có tiệm cận ngang.}{s}
  \y{Giao điểm của tiệm cận đứng và tiệm cận xiên là $I(2;-\dfrac{1}{2})$.}{d}
  \y{Đồ thị hàm số có tất cả $2$ đường tiệm cận.}{d}
  \y{Đường thẳng $y=-\dfrac{1}{2}x+\dfrac{5}{2}$ là tiệm cận xiên của đồ thị hàm số.}{s}}

\cauDS[id=C1-B3-D03-VD-044,muc=VD,loigiai={Xét hàm số $y=\dfrac{-x^{3}+2x^{2}-1}{x^{2}+4}$. Chia đa thức: $y=-x+2+\dfrac{4 x - 9}{x^{2} + 4}$; $a=\lim\limits_{x\to +\infty}\dfrac{y}{x}=-1$, $b=\lim\limits_{x\to +\infty}(y+x)=2$ nên $y=-x+2$ là tiệm cận xiên (khi $x\to +\infty$); $a=\lim\limits_{x\to -\infty}\dfrac{y}{x}=-1$, $b=\lim\limits_{x\to -\infty}(y+x)=2$ nên $y=-x+2$ là tiệm cận xiên (khi $x\to -\infty$). a) Đúng. b) Sai: có $1$ đường tiệm cận. c) Sai. d) Đúng.}]{Cho hàm số $y=\dfrac{-x^{3}+2x^{2}-1}{x^{2}+4}$. Xét tính đúng sai của các mệnh đề sau:}{%
  \y{Đồ thị hàm số không có tiệm cận đứng.}{d}
  \y{Đồ thị hàm số có tất cả $2$ đường tiệm cận.}{s}
  \y{Đường thẳng $y=-x+1$ là tiệm cận xiên của đồ thị hàm số.}{s}
  \y{Đường thẳng $y=-x+2$ là tiệm cận xiên của đồ thị hàm số.}{d}}

\cauDS[id=C1-B3-D03-VD-045,muc=VD,loigiai={Xét hàm số $y=\dfrac{2x^{3}-x^{2}+3}{x^{2}-9}$. Chia đa thức: $y=2x-1+\dfrac{18 x - 6}{x^{2} - 9}$; $a=\lim\limits_{x\to +\infty}\dfrac{y}{x}=2$, $b=\lim\limits_{x\to +\infty}(y-2x)=-1$ nên $y=2x-1$ là tiệm cận xiên (khi $x\to +\infty$); $a=\lim\limits_{x\to -\infty}\dfrac{y}{x}=2$, $b=\lim\limits_{x\to -\infty}(y-2x)=-1$ nên $y=2x-1$ là tiệm cận xiên (khi $x\to -\infty$); $\lim\limits_{x\to {-3}^{+}}y=+\infty$ nên $x=-3$ là tiệm cận đứng; $\lim\limits_{x\to {3}^{+}}y=+\infty$ nên $x=3$ là tiệm cận đứng. a) Đúng. b) Sai. c) Sai: giới hạn của $y$ tại vô cực không hữu hạn. d) Đúng.}]{Cho hàm số $y=\dfrac{2x^{3}-x^{2}+3}{x^{2}-9}$. Xét tính đúng sai của các mệnh đề sau:}{%
  \y{Đồ thị hàm số có tiệm cận đứng $x=-3$.}{d}
  \y{Đường thẳng $y=2x-2$ là tiệm cận xiên của đồ thị hàm số.}{s}
  \y{Đồ thị hàm số có tiệm cận ngang.}{s}
  \y{Giao điểm của tiệm cận đứng và tiệm cận xiên là $I(-3;-7)$.}{d}}

\cauDS[id=C1-B3-D03-VD-046,muc=VD,loigiai={Xét hàm số $y=\sqrt{x^{2}-4x+4}$. $a=\lim\limits_{x\to +\infty}\dfrac{y}{x}=1$, $b=\lim\limits_{x\to +\infty}(y-x)=-2$ nên $y=x-2$ là tiệm cận xiên (khi $x\to +\infty$); $a=\lim\limits_{x\to -\infty}\dfrac{y}{x}=-1$, $b=\lim\limits_{x\to -\infty}(y+x)=2$ nên $y=-x+2$ là tiệm cận xiên (khi $x\to -\infty$). a) Đúng. b) Sai: có $2$ đường tiệm cận. c) Đúng. d) Sai: giới hạn của $y$ tại vô cực không hữu hạn.}]{Cho hàm số $y=\sqrt{x^{2}-4x+4}$. Xét tính đúng sai của các mệnh đề sau:}{%
  \y{Đồ thị hàm số có tất cả $2$ đường tiệm cận.}{d}
  \y{Đồ thị hàm số có tất cả $3$ đường tiệm cận.}{s}
  \y{Đồ thị hàm số không có tiệm cận đứng.}{d}
  \y{Đồ thị hàm số có tiệm cận ngang.}{s}}

\cauDS[id=C1-B3-D03-VD-047,muc=VD,loigiai={Xét hàm số $y=\dfrac{x^{3}-3x^{2}-3}{x^{2}+1}$. Chia đa thức: $y=x-3+\dfrac{- x}{x^{2} + 1}$; $a=\lim\limits_{x\to +\infty}\dfrac{y}{x}=1$, $b=\lim\limits_{x\to +\infty}(y-x)=-3$ nên $y=x-3$ là tiệm cận xiên (khi $x\to +\infty$); $a=\lim\limits_{x\to -\infty}\dfrac{y}{x}=1$, $b=\lim\limits_{x\to -\infty}(y-x)=-3$ nên $y=x-3$ là tiệm cận xiên (khi $x\to -\infty$). a) Sai: giới hạn của $y$ tại vô cực không hữu hạn. b) Sai. c) Đúng. d) Đúng.}]{Cho hàm số $y=\dfrac{x^{3}-3x^{2}-3}{x^{2}+1}$. Xét tính đúng sai của các mệnh đề sau:}{%
  \y{Đồ thị hàm số có tiệm cận ngang.}{s}
  \y{Đường thẳng $y=x-1$ là tiệm cận xiên của đồ thị hàm số.}{s}
  \y{Đồ thị hàm số không có tiệm cận đứng.}{d}
  \y{Đường thẳng $y=x-3$ là tiệm cận xiên của đồ thị hàm số.}{d}}

\cauDS[id=C1-B3-D03-VD-048,muc=VD,loigiai={Xét hàm số $y=\dfrac{-x^{2}+5x+5}{x}$. Chia đa thức: $y=-x+5+\dfrac{5}{x}$; $a=\lim\limits_{x\to +\infty}\dfrac{y}{x}=-1$, $b=\lim\limits_{x\to +\infty}(y+x)=5$ nên $y=-x+5$ là tiệm cận xiên (khi $x\to +\infty$); $a=\lim\limits_{x\to -\infty}\dfrac{y}{x}=-1$, $b=\lim\limits_{x\to -\infty}(y+x)=5$ nên $y=-x+5$ là tiệm cận xiên (khi $x\to -\infty$); $\lim\limits_{x\to {0}^{+}}y=+\infty$ nên $x=0$ là tiệm cận đứng. a) Sai: giới hạn của $y$ tại vô cực không hữu hạn. b) Đúng. c) Đúng. d) Sai.}]{Cho hàm số $y=\dfrac{-x^{2}+5x+5}{x}$. Xét tính đúng sai của các mệnh đề sau:}{%
  \y{Đồ thị hàm số có tiệm cận ngang.}{s}
  \y{Đồ thị hàm số có tất cả $2$ đường tiệm cận.}{d}
  \y{Đường thẳng $y=-x+5$ là tiệm cận xiên của đồ thị hàm số.}{d}
  \y{Đường thẳng $y=-x+4$ là tiệm cận xiên của đồ thị hàm số.}{s}}

\cauTLN[id=C1-B3-D03-VD-049,muc=VD,loigiai={Xét hàm số $y=\dfrac{x^{3}}{x^{2}-9}$. Chia đa thức: $y=x+\dfrac{9 x}{x^{2} - 9}$; $a=\lim\limits_{x\to +\infty}\dfrac{y}{x}=1$, $b=\lim\limits_{x\to +\infty}(y-x)=0$ nên $y=x$ là tiệm cận xiên (khi $x\to +\infty$); $a=\lim\limits_{x\to -\infty}\dfrac{y}{x}=1$, $b=\lim\limits_{x\to -\infty}(y-x)=0$ nên $y=x$ là tiệm cận xiên (khi $x\to -\infty$); $\lim\limits_{x\to {-3}^{+}}y=+\infty$ nên $x=-3$ là tiệm cận đứng; $\lim\limits_{x\to {3}^{+}}y=+\infty$ nên $x=3$ là tiệm cận đứng. Kết quả là $2$.}]{Tiệm cận xiên của đồ thị hàm số $y=\dfrac{x^{3}}{x^{2}-9}$ có phương trình $y=ax+b$. Tính $2a-b$.}{2}

\cauTLN[id=C1-B3-D03-VD-050,muc=VD,loigiai={Xét hàm số $y=3x+\sqrt{x^2+4}$. $a=\lim\limits_{x\to +\infty}\dfrac{y}{x}=4$, $b=\lim\limits_{x\to +\infty}(y-4x)=0$ nên $y=4x$ là tiệm cận xiên (khi $x\to +\infty$); $a=\lim\limits_{x\to -\infty}\dfrac{y}{x}=2$, $b=\lim\limits_{x\to -\infty}(y-2x)=0$ nên $y=2x$ là tiệm cận xiên (khi $x\to -\infty$). Kết quả là $6$.}]{Tính tổng các hệ số góc của các đường tiệm cận xiên của đồ thị hàm số $y=3x+\sqrt{x^2+4}$.}{6}

\cauTLN[id=C1-B3-D03-VD-051,muc=VD,loigiai={Xét hàm số $y=x-\sqrt{x^{2}-4x}$. $\lim\limits_{x\to +\infty}y=2$ nên $y=2$ là tiệm cận ngang; $a=\lim\limits_{x\to -\infty}\dfrac{y}{x}=2$, $b=\lim\limits_{x\to -\infty}(y-2x)=-2$ nên $y=2x-2$ là tiệm cận xiên (khi $x\to -\infty$). Kết quả là $-2$.}]{Tiệm cận xiên của đồ thị hàm số $y=x-\sqrt{x^{2}-4x}$ cắt trục tung tại điểm có tung độ bằng bao nhiêu?}{-2}

\cauTLN[id=C1-B3-D03-VD-052,muc=VD,loigiai={Xét hàm số $y=\dfrac{-2x^{2}+2x-3}{x+1}$. Chia đa thức: $y=-2x+4+\dfrac{-7}{x + 1}$; $a=\lim\limits_{x\to +\infty}\dfrac{y}{x}=-2$, $b=\lim\limits_{x\to +\infty}(y+2x)=4$ nên $y=-2x+4$ là tiệm cận xiên (khi $x\to +\infty$); $a=\lim\limits_{x\to -\infty}\dfrac{y}{x}=-2$, $b=\lim\limits_{x\to -\infty}(y+2x)=4$ nên $y=-2x+4$ là tiệm cận xiên (khi $x\to -\infty$); $\lim\limits_{x\to {-1}^{+}}y=-\infty$ nên $x=-1$ là tiệm cận đứng. Kết quả là $2$.}]{Tiệm cận xiên của đồ thị hàm số $y=\dfrac{-2x^{2}+2x-3}{x+1}$ có phương trình $y=ax+b$. Tính $a+b$.}{2}

\cauTLN[id=C1-B3-D03-VD-053,muc=VD,loigiai={Xét hàm số $y=2x+\sqrt{x^2+4}$. $a=\lim\limits_{x\to +\infty}\dfrac{y}{x}=3$, $b=\lim\limits_{x\to +\infty}(y-3x)=0$ nên $y=3x$ là tiệm cận xiên (khi $x\to +\infty$); $a=\lim\limits_{x\to -\infty}\dfrac{y}{x}=1$, $b=\lim\limits_{x\to -\infty}(y-x)=0$ nên $y=x$ là tiệm cận xiên (khi $x\to -\infty$). Kết quả là $2$.}]{Đồ thị hàm số $y=2x+\sqrt{x^2+4}$ có bao nhiêu đường tiệm cận xiên?}{2}

\cauTLN[id=C1-B3-D03-VD-054,muc=VD,loigiai={Xét hàm số $y=\sqrt{x^{2}+4x+7}$. $a=\lim\limits_{x\to +\infty}\dfrac{y}{x}=1$, $b=\lim\limits_{x\to +\infty}(y-x)=2$ nên $y=x+2$ là tiệm cận xiên (khi $x\to +\infty$); $a=\lim\limits_{x\to -\infty}\dfrac{y}{x}=-1$, $b=\lim\limits_{x\to -\infty}(y+x)=-2$ nên $y=-x-2$ là tiệm cận xiên (khi $x\to -\infty$). Kết quả là $0$.}]{Tính tổng các hệ số góc của các đường tiệm cận xiên của đồ thị hàm số $y=\sqrt{x^{2}+4x+7}$.}{0}

\cauTLN[id=C1-B3-D03-VD-055,muc=VD,loigiai={Xét hàm số $y=x-\sqrt{x^{2}-2x}$. $\lim\limits_{x\to +\infty}y=1$ nên $y=1$ là tiệm cận ngang; $a=\lim\limits_{x\to -\infty}\dfrac{y}{x}=2$, $b=\lim\limits_{x\to -\infty}(y-2x)=-1$ nên $y=2x-1$ là tiệm cận xiên (khi $x\to -\infty$). Kết quả là $5$.}]{Tiệm cận xiên của đồ thị hàm số $y=x-\sqrt{x^{2}-2x}$ có phương trình $y=ax+b$. Tính $2a-b$.}{5}

\cauTLN[id=C1-B3-D03-VD-056,muc=VD,loigiai={Xét hàm số $y=\dfrac{-x^{3}-2}{x^{2}-1}$. Chia đa thức: $y=-x+\dfrac{- x - 2}{x^{2} - 1}$; $a=\lim\limits_{x\to +\infty}\dfrac{y}{x}=-1$, $b=\lim\limits_{x\to +\infty}(y+x)=0$ nên $y=-x$ là tiệm cận xiên (khi $x\to +\infty$); $a=\lim\limits_{x\to -\infty}\dfrac{y}{x}=-1$, $b=\lim\limits_{x\to -\infty}(y+x)=0$ nên $y=-x$ là tiệm cận xiên (khi $x\to -\infty$); $\lim\limits_{x\to {-1}^{+}}y=+\infty$ nên $x=-1$ là tiệm cận đứng; $\lim\limits_{x\to {1}^{+}}y=-\infty$ nên $x=1$ là tiệm cận đứng. Kết quả là $-2$.}]{Tiệm cận xiên của đồ thị hàm số $y=\dfrac{-x^{3}-2}{x^{2}-1}$ có phương trình $y=ax+b$. Tính $2a-b$.}{-2}

\cauTLN[id=C1-B3-D03-VD-057,muc=VD,loigiai={Xét hàm số $y=2x+\sqrt{x^2+9}$. $a=\lim\limits_{x\to +\infty}\dfrac{y}{x}=3$, $b=\lim\limits_{x\to +\infty}(y-3x)=0$ nên $y=3x$ là tiệm cận xiên (khi $x\to +\infty$); $a=\lim\limits_{x\to -\infty}\dfrac{y}{x}=1$, $b=\lim\limits_{x\to -\infty}(y-x)=0$ nên $y=x$ là tiệm cận xiên (khi $x\to -\infty$). Kết quả là $4$.}]{Tính tổng các hệ số góc của các đường tiệm cận xiên của đồ thị hàm số $y=2x+\sqrt{x^2+9}$.}{4}

\cauTLN[id=C1-B3-D03-VD-058,muc=VD,loigiai={Xét hàm số $y=\sqrt{x^{2}-2x+6}$. $a=\lim\limits_{x\to +\infty}\dfrac{y}{x}=1$, $b=\lim\limits_{x\to +\infty}(y-x)=-1$ nên $y=x-1$ là tiệm cận xiên (khi $x\to +\infty$); $a=\lim\limits_{x\to -\infty}\dfrac{y}{x}=-1$, $b=\lim\limits_{x\to -\infty}(y+x)=1$ nên $y=-x+1$ là tiệm cận xiên (khi $x\to -\infty$). Kết quả là $0$.}]{Tính tổng các hệ số góc của các đường tiệm cận xiên của đồ thị hàm số $y=\sqrt{x^{2}-2x+6}$.}{0}

\cauTLN[id=C1-B3-D03-VD-059,muc=VD,loigiai={Xét hàm số $y=x-\sqrt{x^{2}-2x}$. $\lim\limits_{x\to +\infty}y=1$ nên $y=1$ là tiệm cận ngang; $a=\lim\limits_{x\to -\infty}\dfrac{y}{x}=2$, $b=\lim\limits_{x\to -\infty}(y-2x)=-1$ nên $y=2x-1$ là tiệm cận xiên (khi $x\to -\infty$). Kết quả là $-1$.}]{Tiệm cận xiên của đồ thị hàm số $y=x-\sqrt{x^{2}-2x}$ cắt trục tung tại điểm có tung độ bằng bao nhiêu?}{-1}

\cauTLN[id=C1-B3-D03-VD-060,muc=VD,loigiai={Xét hàm số $y=-2x+\sqrt{x^2+1}$. $a=\lim\limits_{x\to +\infty}\dfrac{y}{x}=-1$, $b=\lim\limits_{x\to +\infty}(y+x)=0$ nên $y=-x$ là tiệm cận xiên (khi $x\to +\infty$); $a=\lim\limits_{x\to -\infty}\dfrac{y}{x}=-3$, $b=\lim\limits_{x\to -\infty}(y+3x)=0$ nên $y=-3x$ là tiệm cận xiên (khi $x\to -\infty$). Kết quả là $2$.}]{Đồ thị hàm số $y=-2x+\sqrt{x^2+1}$ có bao nhiêu đường tiệm cận xiên?}{2}
